%
% Tex input file for CGNS Mid-Level Library Manual
%
% To generate a DVI file, then a PostScript file named midlevel.ps:
%
%     latex midlevel.tex
%     dvips midlevel.dvi -o
%
% To generate a PDF file named midlevel.pdf:
%
%     pdflatex midlevel.tex


\documentclass[twoside,fleqn]{article}

%
% Packages
%
%\usepackage{times}			% use PostScript Times-Roman for text
\usepackage{indentfirst}		% indent first paragraph in sections
%\usepackage[first,light]{draftcopy}	% add DRAFT watermark
\usepackage[tbtags]{amsmath}		% use ams math package
\usepackage{graphicx}			% for including eps files
%\usepackage{subfigure}			% for sub-figures
\usepackage[bf]{caption}		% for more flexible captions
   \setlength{\captionmargin}{18pt}	% caption margins
\usepackage{array}			% use new array and table package
   \setlength{\extrarowheight}{2pt}	% increase row heights in tables
\usepackage{tabularx}			% use tabularx package
\usepackage{longtable}			% use longtable package
\usepackage{alltt}			% verbatim input with font changes
%\usepackage{lscape}			% for landscape pages
%\usepackage{mdwtab}			% use mdwtab array and table package
%\usepackage{tabls}			% access \hline[dimen] feature
%\usepackage{dcolumn}			% allow table alignment at dec points
%   \newcolumntype{d}[1]{D{.}{.}{#1}}	% for tables, arg is # decimal places
\usepackage{natbib}			% for author-year citations
   \setlength{\bibhang}{0pt}		% don't use hanging indent format
\usepackage[normalem]{ulem}		% allow underlining
\usepackage{calc}			% calculation package
\usepackage{ifthen}			% control structures
\usepackage{fancyhdr}			% fancy headers and footers
   \fancyhf{}
   \renewcommand{\headrulewidth}{0pt}
   \renewcommand{\footrulewidth}{0pt}
\usepackage{titling}			% more title page control
\usepackage{xspace}			% common-sense spacing after text macro
\usepackage{mdwlist}			% more flexible description lists
\usepackage{colortbl}			% colored table cells
\usepackage{color}			% colors and grays
   \definecolor{subcolor}{rgb}{0.808,0.851,1.0}
   \definecolor{input}{rgb}{0.200,0.200,1.0}
   \definecolor{output}{rgb}{0.600,0.0,0.0}
\usepackage{url}			% typesetting urls
   \urlstyle{rm}
%\usepackage[T1]{fontenc}		% for better-looking "_" and tt "{"
\usepackage{ae}				% used instead of above for PDF output
\usepackage{hyperref}			% for hypertext links in PDF
   \hypersetup{letterpaper,plainpages=false,
               pdftitle={CGNS Mid-Level Library},
               pdfauthor={CGNS Project Group},
               colorlinks,
               linkcolor=blue,citecolor=blue,filecolor=red,pagecolor=blue,
               urlcolor=red}
   \renewcommand{\sectionautorefname}{Section}
   \renewcommand{\subsectionautorefname}{Section}
   \renewcommand{\subsubsectionautorefname}{Section}

%
% Page layout
%
\normalsize
\setlength{\oddsidemargin}{0.5in}	% left margins-1.0in for odd/even pages
\setlength{\evensidemargin}{0.0in}
\setlength{\textwidth}{6.0in}		% text width
%\setlength{\marginparwidth}{0in}	% margin note space parameters
%\setlength{\marginparsep}{0in}
%\setlength{\oddsidemargin}{0.25in}	% left margins-1.0in for odd/even pages
%\setlength{\evensidemargin}{0.75in}
%\setlength{\textwidth}{5.5in}		% text width
\setlength{\marginparwidth}{1.1in}	% margin note space parameters
\setlength{\marginparsep}{0.2in}
\reversemarginpar
\setlength{\topmargin}{-0.5in}		% top margin-1.0in
\setlength{\headheight}{0.25in}		% header space parameters
\setlength{\headsep}{0.5in}
\setlength{\textheight}{8.5in}		% text height
\setlength{\topskip}{\baselineskip}	% dist from top of body to 1st baseline
\setlength{\footskip}{0.75in}		% dist from bottom of body to footer
\raggedbottom				% to avoid stretching vertical space
\pagestyle{fancy}			% using fancyhdr package
\fancypagestyle{plain}{%		% page # in corner for plain style
   \fancyhf{}
   \fancyfoot[LE,RO]{\bfseries \thepage}}

%
% Misc spacing parameters
%
%\setlength{\parskip}{\baselineskip}	% blank space between paragraphs
\setlength{\parskip}{0.5\baselineskip}	% blank space between paragraphs
\setlength{\doublerulesep}{0pt}		% for wider \hlines
\setlength{\fboxsep}{5pt}		% for more margin in boxes
\newlength{\saveparindent}		% save for use where LaTeX changes them
\newlength{\saveparskip}		% save for use where LaTeX changes them
\newlength{\savebaselineskip}
\setlength{\saveparindent}{\parindent}
\setlength{\saveparskip}{\parskip}
\setlength{\savebaselineskip}{\baselineskip}
\newlength{\tmplength}			% for temp use wherever needed
\newlength{\Pwidth}			% for use in tables
\newlength{\Pwidtha}

%
% User-defined environments
%
% "Definition" list with variable-length entries, indented. Requires
% mdwlist package. (See LaTeX Companion, p 64, and mdwlist documentation.)
\newenvironment{Ventryi}[1]%
   {\settowidth{\tmplength}{\hspace{\saveparindent}#1\hspace{1em}}%
    \begin{basedescript}{%
           \desclabelwidth{\tmplength}%
           \desclabelstyle{\nextlinelabel}%
           \renewcommand{\makelabel}[1]{\hspace{\saveparindent}##1\hfil}}}
   {\end{basedescript}}
% Indented definition list with variable-length entries, compact. Requires
% mdwlist package. (See LaTeX Companion, p 64, and mdwlist documentation.)
\newenvironment{Ventryic}[1]%
   {\settowidth{\tmplength}{\hspace{\saveparindent}#1\hspace{1em}}%
    \begin{basedescript}{%
           \desclabelwidth{\tmplength}%
           \desclabelstyle{\nextlinelabel}%
           \renewcommand{\makelabel}[1]{\hspace{\saveparindent}##1\hfil}
           \setlength{\topsep}{0in}%
           \setlength{\parsep}{0in}%
           \setlength{\itemsep}{0in}}}
   {\end{basedescript}}

% "changes" environment, to identify changed code by color.  The color is
% given by the argument, with a default of red.  When the environment is
% exited, the color is changed to black.
\newenvironment{changes}[1][red]%
   {\color{#1}}
   {\color{black}}

% Adapted from comp.text.tex post by Keith Reckdahl
% (reckdahl@am-sparc7.Stanford.EDU)
% Indents text from the left margin by current value of the list
% length \leftmargin
%\newenvironment{indleft}%
%   {\begin{list}{}
%      {\setlength{\topsep}{0pt}%
%       \setlength{\listparindent}{0pt}%
%       \setlength{\itemindent}{0pt}%
%       \setlength{\parsep}{\parskip}%
%      }%
%      \item[]}%
%   {\end{list}}
% As above, but surrounding an alltt environment
\newenvironment{indlefttt}%
   {\begin{list}{}
      {\setlength{\topsep}{0pt}%
       \setlength{\listparindent}{0pt}%
       \setlength{\itemindent}{0pt}%
       \setlength{\parsep}{\parskip}%
      }%
      \item[]%
      \begin{alltt}}%
   {\end{alltt}\end{list}}
%
% Framed box for function syntax
% Works
%\newlength{\Headwidtha}
%\newlength{\savearrayrulewidth}
%\newenvironment{fctbox}
%   {\noindent
%    \setlength{\savearrayrulewidth}{\arrayrulewidth}%
%    \setlength{\arrayrulewidth}{0.6pt}%
%    \setlength{\Headwidtha}{\textwidth-3\arrayrulewidth-4\tabcolsep-0.45in}%
%    \setlength{\tmplength}{\tabcolsep+0.3in}%
%    \begin{flushleft}
%    \begin{tabular}{|>{\ttfamily\columncolor{subcolor}}p{\Headwidtha} |%
%	  >{\centering\ttfamily\columncolor{subcolor}[\tabcolsep][\tmplength]}p{0.15in}%
%	  @{\centering}p{0.15in} @{\centering}p{0.15in} |}%
%    \hline
%    \textnormal{\textbf{Functions}} & \textnormal{\textbf{Modes}} & & \\
%    \hline
%   }
%   {\hline
%    \end{tabular}
%    \setlength{\arrayrulewidth}{\savearrayrulewidth}%
%    \end{flushleft}
%   }

% 
\newlength{\Headwidtha}
\newlength{\savearrayrulewidth}
\newenvironment{fctbox}
   {\noindent
    \setlength{\savearrayrulewidth}{\arrayrulewidth}%
    \setlength{\arrayrulewidth}{0.8pt}%
    \settowidth{\tmplength}{\textbf{Modes}}
    \setlength{\Headwidtha}{\textwidth-3\arrayrulewidth-4\tabcolsep-\tmplength}%
    \begin{flushleft}
    \begin{tabular}{|>{\ttfamily\columncolor{subcolor}}p{\Headwidtha}%
                    |>{\ttfamily\columncolor{subcolor}}c |}
    \hline
    \textnormal{\textbf{Functions}} & \textnormal{\textbf{Modes}} \\
    \hline
   }
   {\hline
    \end{tabular}
    \setlength{\arrayrulewidth}{\savearrayrulewidth}%
    \end{flushleft}
   }

% Doesn't work ...
%\newlength{\Headwidtha}
%\newlength{\savearrayrulewidth}
%\newenvironment{fctbox}
%   {\noindent
%    \setlength{\savearrayrulewidth}{\arrayrulewidth}%
%    \setlength{\arrayrulewidth}{0.6pt}%
%    \setlength{\Headwidtha}{\textwidth-3\arrayrulewidth-4\tabcolsep-0.45in}%
%    \setlength{\tmplength}{\tabcolsep+0.3in}%
%    \begin{flushleft}
%    \begin{tabular}{|>{\ttfamily\columncolor{subcolor}}p{\Headwidtha} |%
%          >{\centering\ttfamily\columncolor{subcolor}[\tabcolsep][\tmplength]}p{0.45in} |}%
%    \hline
%    \textnormal{\textbf{Functions}} & \textnormal{\textbf{Modes}} \\
%    \hline
%   }
%   {\hline
%    \end{tabular}
%    \setlength{\arrayrulewidth}{\savearrayrulewidth}%
%    \end{flushleft}
%   }

%
% User-defined commands
%

% Function name initialization
\newcommand{\fctname}{}

% Function header
\newcommand{\fctheader}[3]
   {\newpage%
    \hypertarget{fct:#3}{}%
    \noindent
    \ital{#1} (\texttt{#2})
    \smallskip
    \renewcommand{\fctname}{#1}
%    \fancyfoot[LE]{\bfseries \thepage}
%    \fancyfoot[RO]{\bfseries \thepage}
    \addtocontents{toc}{\protect\contentsline {subsubsection}{\ital{\fctname}\ (\texttt{#2})}{\thepage}{fct:#3}}
   }

% Function header, no new page
\newcommand{\fctheadernnp}[3]
   {\noindent
    \hypertarget{fct:#3}{}%
    \ital{#1} (\texttt{#2})
    \smallskip
    \renewcommand{\fctname}{\ital{#1}}
%    \fancyfoot[LE]{\bfseries \thepage}
%    \fancyfoot[RO]{\bfseries \thepage}
    \addtocontents{toc}{\protect\contentsline {subsubsection}{\fctname\ (\texttt{#2})}{\thepage}{fct:#3}}
   }

% Re-def longtable's caption command to use \captionlabelfont from caption pkg
% (basis lifted from longtable.sty)
\makeatletter
\renewcommand{\LT@makecaption}[3]{%
  \LT@mcol\LT@cols c{\hbox to\z@{\hss\parbox[t]\LTcapwidth{%
    \sbox\@tempboxa{{\captionlabelfont #1{#2: }}#3}%
    \ifdim\wd\@tempboxa>\hsize
      {\captionlabelfont #1{#2: }}#3%
    \else
      \hbox to\hsize{\hfil\box\@tempboxa\hfil}%
    \fi
    \endgraf\vskip\baselineskip}%
  \hss}}}
\makeatother

% Make next page odd, with preceding blank page empty (LaTeX Companion, p 93)
\newcommand{\clearemptydoublepage}{\newpage{\pagestyle{empty}\cleardoublepage}}

% Text subscripts, analogous to \textsuperscript command (from comp.text.tex
% post by rf@cl.cam.ac.uk (Robin Fairbairns))
\makeatletter
\DeclareRobustCommand*\tsub[1]{%
  \@tsub{\selectfont#1}}
\def\@tsub#1{%
  {\m@th\ensuremath{_{\mbox{\fontsize\sf@size\z@#1}}}}}
\makeatother

% Text superscripts, but shorter
\newcommand{\tsup}[1]{\textsuperscript{#1}}

% Shortcuts for specific fonts
\newcommand{\bold}[1]{{\normalfont\textbf{#1}}}  % Bold
\newcommand{\ital}[1]{{\normalfont\textit{#1}}}  % Italic
\newcommand{\key}[1]{{\normalfont\texttt{#1}}}  % Fixed font for keywords

% Examples
\newcommand{\Example}[1]{\noindent\uline{\textit{#1}}}

% Use section titles as marks (i.e., in headers) instead of subsection titles
\renewcommand{\sectionmark}[1]{\markright{\thesection\ \ #1}}
\renewcommand{\subsectionmark}[1]{}

% "Better" treatment of headings for Appendices  (LaTeX Companion, pp 29,30)
\renewcommand{\appendix}{%
   \newcommand{\app}{%
      \secdef\Appendix\sAppendix}%
   \newcommand{\subapp}{%
      \secdef\subAppendix\subsAppendix}%
   \newcounter{app}%
   \newcounter{subapp}[app]%
   \renewcommand{\theapp}{\Alph{app}}%
   \renewcommand{\thesubapp}{\theapp.\arabic{subapp}}%
 }
\newcommand{\Appendix}[2][?]{%   Complex form
   \refstepcounter{app}%
   \addcontentsline{toc}{section}%
      {\protect\bfseries{\appendixname~\theapp.\ \ } #1}%
   {\noindent\Large\bfseries\appendixname\ \theapp.\ \ #2\par}%
   \sectionmark{#1}\vspace{\baselineskip}}
\newcommand{\sAppendix}[1]{%   Starred form
   \refstepcounter{app}%
   {\noindent\Large\bfseries\appendixname\ \theapp.\ \ #1\par}%
   \vspace{\baselineskip}}

\newcommand{\subAppendix}[2][?]{%   Complex form
   \refstepcounter{subapp}%
   \addcontentsline{toc}{subsection}%
       {\thesubapp\ \ #1}%
   \vspace{\baselineskip}
   {\noindent\large\bfseries\thesubapp\ \ #2\par}%
   \vspace{\baselineskip}}
\newcommand{\subsAppendix}[1]{%   Starred form
   \refstepcounter{subapp}%
   \vspace{\baselineskip}
   {\noindent\large\bfseries\thesubapp\ \ #1\par}%
   \vspace{\baselineskip}}

\begin{document}

\pagenumbering{roman}

%\fancyhead[LE]{\bfseries Mid-Level Library}
%\fancyhead[RO]{\bfseries \rightmark}
\fancyfoot[LE,RO]{\bfseries \thepage}

%\setlength{\droptitle}{1.0in}
\pretitle{\begin{flushleft}\LARGE%
          \includegraphics[width=2.0in]{../logo/cgns_bw}\\*[0.25in]}
\posttitle{\par\end{flushleft}\vskip 1.0em}
\title{{\bfseries CFD General Notation System\\
Mid-Level Library}\\*[0.25in]
{\Large Document Version 2.5.3\\
CGNS Version 2.5}}
\author{}
\date{}
\maketitle
\thispagestyle{empty}

%\title{{\LARGE \textbf{CGNS}}\\*[0.5in]
%The CFD General Notation System\\
%Mid-Level Library\footnote
%{This is an unnumbered version of this document, created \today.
%It was derived from the HTML version of the document, originally
%written by Diane Poirier of ICEM CFD Engineering.}}
%\author{CGNS Project Group}
%\date{}
%\maketitle
%\thispagestyle{empty}

\clearemptydoublepage
\setlength{\parskip}{0ex}		% remove blank space between paragraphs
\thispagestyle{plain}
\tableofcontents
\setlength{\parskip}{\saveparskip}	% put it back

\clearemptydoublepage
\pagenumbering{arabic}

\fancyhead[LE]{\bfseries Mid-Level Library}
\fancyhead[RO]{\bfseries \rightmark}

\section{Introduction}
\label{s:intro}
\thispagestyle{plain}

Advanced Data Format (ADF) is a library of basic database management
and I/O subroutines that implements a relatively simple hierarchical
database.
ADF is written in ANSI C to enhance portability of the software, and the
design of the database allows for portability of files from platform to
platform.
There is also a Fortran interface.
The files are self-describing (i.e., it is possible to browse files and
determine their contents) and extensible (i.e., many different pieces of
software on different platforms may add or modify information).

The routines allow the user to construct a tree structure with their data.
(See \autoref{f:example} on p.~\pageref*{f:example}.)
This structure is very similar to the directory structures of the UNIX
or DOS operating systems.
ADF also allows links between nodes within the same file or to different
files.
This feature works somewhat like ``soft links'' in the UNIX operating
system.
The major difference between the aforementioned directory structures and
ADF is that the nodes not only contain information about their children
(next lower-level nodes) but may also contain data.

The installation package includes the source code to ADF, the Fortran
interface, sample files, and a simple file browser.

\subsection{History of Project}

ADF was developed as part of the CFD
General Notation System(CGNS) project.
The purpose of the CGNS project is to define the data models for
Navier-Stokes based Computational Fluid Dynamics (CFD) technology,
develop a set of standard interface data structures for those data
models, and develop the software that will allow implementation of those
data structures in existing and future CFD analysis tools.
The CGNS system consists of a collection of conventions, and software
conforming to these conventions, for the storage and retrieval of CFD
data.
Adherence to these conventions is intended to facilitate the exchange of
CFD data between sites, between application codes such as solvers and
grid generators, and across computing platforms.

Once the data models (called the \href{../sids/sids.pdf}{Standard
Interface Data Structures}, or SIDS) were defined, a project was started
to write software that could faithfully reproduce that information on
disk.
ADF is the result of that project.
While ADF was developed specifically for the CFD process, it is quite
general and has no built-in knowledge of CFD; therefore, it should
be applicable to storing any type of data that lends itself to a
hierarchical definition.

\subsection{Other Software Interfaces and ADF}

Two other software packages were investigated as part of this project.
The first database interface investigated was the
Hierarchical Data Format (HDF)
developed at the National Center for
Supercomputing Applications at the University of Illinois.
This database system has a large user base and support and has been
in existence more than 6 years with utilities and graphical routines
written with both a C and a Fortran interface.
The limitation of HDF was that it was not truly hierarchical despite its
name.
Any hierarchy has to be built using naming conventions.
Since the CGNS data models indicated a natural hierarchical structure,
it seemed appropriate to develop database software that worked in that
mode by design.
The ADF design considerations are summarized in
\hyperref[s:design]{Appendix~\ref*{s:design}}.

The second was the
Common File Format
(CFF) developed by McDonnell Douglas Aerospace.
CFF is a second-generation database management system that provides a
unifying file structure for CFD data.
The purpose of the Common File is to insulate the user from the myriad
of different computer types that make up computer systems so that the
user or application programmer may process a file from another machine
without performing explicit conversions.
CFF was written in Fortran; however, it was felt that portability and
extensibility could be enhanced using C.
Much of the experience gained from the McDonnell Douglas group is
incorporated into ADF due to the cooperative efforts of personnel from
the CFD group at McDonnell Douglas Aerospace.

\subsection{Organization of Manual}

The main section of this manual explains the basics of the
hierarchical structure of an ADF database.
In addition, the concept of a node as
the basic building block of the hierarchy is developed in detail.
The remaining sections of this manual are extensive.
They provide a glossary of terms and
conventions, as well as information related to the
ADF version releases,
version control numbering, and
architectures that are supported by ADF.
There are two examples, one in Fortran and one in C, in
\hyperref[s:sampleFortran]{Appendix~\ref*{s:sampleFortran}} and
\hyperref[s:sampleC]{Appendix~\ref*{s:sampleC}} respectively, that
implement the structure illustrated in \autoref{f:example} on
p.~\pageref*{f:example}.
These examples should help familiarize the new user with ADF.
Design considerations, file optimization, and portability issues are
also discussed.
The individual ADF core routines are described in detail in
\hyperref[s:subs]{Appendix~\ref*{s:subs}}.
They are categorized into
database-level routines,
data structure and management subroutines,
data query subroutines,
data I/O subroutines, and some
miscellaneous utility subroutines.
Numerous examples are included to clarify the use of each subroutine.
\hyperref[s:errors]{Appendix~\ref*{s:errors}} provides a summary of ADF
error messages, and \hyperref[s:defaults]{Appendix~\ref*{s:defaults}}
lists default values for various parameters and limits on dimensions of
arrays.


\clearemptydoublepage
\section{General Remarks}
\label{s:general}
\thispagestyle{plain}

\subsection{Acquiring the Software and Documentation}

The CGNS Mid-Level Library may be downloaded from SourceForge, at
{\itshape\url{http://sourceforge.net/projects/cgns}}.
This manual, as well as the other CGNS documentation, is available in
both HTML and PDF format from the CGNS documentation web site, at
{\itshape\url{http://www.grc.nasa.gov/www/cgns/CGNS_docs_current/}}.

\subsection{Organization of This Manual}

The sections that follow describe the Mid-Level Library functions
in detail.
The first three sections cover some basic file operations (i.e.,
opening and closing a CGNS file, and some configuration options)
(\autoref{s:fileops}), accessing a specific node in a CGNS database
(\autoref{s:navigating}), and error handling (\autoref{s:error}).
The remaining sections describe the functions used to read, write, and
modify nodes and data in a CGNS database.
These sections basically follow the organization used in
the ``Detailed CGNS Node Descriptions'' section of the
\href{../filemap/filemap.pdf}{SIDS File Mapping} manual.

At the start of each sub-section is a \textit{Node} line, listing the
the applicable CGNS node label.

Next is a table illustrating the syntax for the Mid-Level Library
functions.
The C functions are shown in the top half of the table, followed by the
corresponding Fortran routines in the bottom half of the table.
Input variables are shown in an \textcolor{input}{\texttt{upright blue}}
font, and output variables are shown in a
\textcolor{output}{\texttt{\textit{slanted red}}} font.
As of Version 3.1, some of the arguments to the Mid-Level Library have
changed from \texttt{int} to \texttt{cgsize\_t} in order to support 64-bit data.
For each function, the right-hand column lists the modes (read, write,
and/or modify) applicable to that function.

The input and output variables are then listed and defined.

\subsection{Language}

The CGNS Mid-Level Library is written in C, but each function has a
Fortran counterpart.
All function names start with ``\texttt{cg\_}''.
The Fortran functions have the same name as their C counterpart with the
addition of the suffix ``\texttt{\_f}''.

\subsection{Character Strings}

All data structure names and labels in CGNS are limited to 32
characters.
When reading a file, it is advised to pre-allocate the
character string variables to 32 characters in Fortran, and 33
in C (to include the string terminator).
Other character strings, such as the CGNS file name or
descriptor text, are unlimited in length.
The space for unlimited length character strings will be created by the
Mid-Level Library; it is then the responsibility of the application to
release this space by a call to \texttt{cg\_free}, described in
\autoref{s:free}.

\subsection{Error Status}

All C functions return an integer value representing the error
status.
All Fortran functions have an additional parameter, \texttt{ier},
which contains the value of the error status.
An error status different from zero implies that an error occured.
The error message can be printed using the error
handling functions of the CGNS library, described in \autoref{s:error}.
The error codes are coded in the C and Fortran include files
\textit{cgnslib.h} and \textit{cgnslib\_f.h}.

\subsection{Typedefs}
\label{s:typedefs}

Beginning with version 3.1, two new typedef variables have been
introduced to support 64-bit mode. The \texttt{cglong\_t} typedef
is always a 64-bit integer, and \texttt{cgsize\_t} will be either
a 32-bit or 64-bit integer depending on how the library was built.
Many of the C functions in the MLL have been changed
to to use \texttt{cgsize\_t} instead of \texttt{int} in the arguments.
These functions include any that may exceed the \texttt{2Gb} limit
of an \texttt{int}, e.g. zone dimensions, element data, boundary
conditions, and connectivity. In Fortran, all integer data is taken
to be \texttt{integer*4} for 32-bit and \texttt{integer*8} for 64-bit
builds.

Several types of variables are defined using typedefs in the
\textit{cgnslib.h} file.
These are intended to facilitate the implementation of
CGNS in C.
These variable types are defined as an enumeration of key
words admissible for any variable of these types.
The file \textit{cgnslib.h} must be included in any C application
programs which use these data types.

In Fortran, the same key words are defined as integer parameters in the
include file \textit{cgnslib\_f.h}.
Such variables should be declared as \texttt{integer} in Fortran
applications.
The file \textit{cgnslib\_f.h} must be included in any
Fortran application using these key words.

Note that the first two enumerated values in these lists,
\textit{xxxNull} and \textit{xxxUserDefined}, are only available
in the C interface, and are provided in the advent that your
C compiler does strict type checking. In Fortran, these values
are replaced by the numerically equivalent \texttt{CG\_Null}
and \texttt{CG\_UserDefined}. These values are also defined in the
C interface, thus either form may be used. The function prototypes
for the MLL use \texttt{CG\_Null} and \texttt{CG\_UserDefined},
rather than the more specific values.

The list of enumerated values (key words) for each of these variable
types (typedefs) are:
        
{\raggedright
\begin{Ventryi}{\texttt{AverageInterfaceType\_t}}
   \item [\texttt{ZoneType\_t}]
         \texttt{ZoneTypeNull, ZoneTypeUserDefined, Structured, Unstructured}
   \item [\texttt{ElementType\_t}]
         \texttt{ElementTypeNull, ElementTypeUserDefined, NODE, BAR\_2, BAR\_3, TRI\_3, TRI\_6, QUAD\_4,
         QUAD\_8, QUAD\_9, TETRA\_4, TETRA\_10, PYRA\_5,
         PYRA\_14, PENTA\_6, PENTA\_15, PENTA\_18, HEXA\_8,
         HEXA\_20, HEXA\_27, MIXED, PYRA\_13, NGON\_n, NFACE\_n}
   \item [\texttt{DataType\_t}]
         \texttt{DataTypeNull, DataTypeUserDefined, Integer, RealSingle, RealDouble, Character}
   \item [\texttt{DataClass\_t}]
         \texttt{DataClassNull, DataClassUserDefined, Dimensional, NormalizedByDimensional,
         NormalizedByUnknownDimensional, NondimensionalParameter,
         DimensionlessConstant}
   \item [\texttt{MassUnits\_t}]
         \texttt{MassUnitsNull, MassUnitsNullUserDefined, Kilogram, Gram, Slug, PoundMass}
   \item [\texttt{LengthUnits\_t}]
         \texttt{LengthUnitsNull, LengthUnitsUserDefined, Meter, Centimeter, Millimeter, Foot,
         Inch}
   \item [\texttt{TimeUnits\_t}]
         \texttt{TimeUnitsNull, TimeUnitsUserDefined, Second}
   \item [\texttt{TemperatureUnits\_t}]
         \texttt{TemperatureUnitsNull, TemperatureUnitsUserDefined, Kelvin, Celsius, Rankine,
         Fahrenheit}
   \item [\texttt{AngleUnits\_t}]
         \texttt{AngleUnitsNull, AngleUnitsUserDefined, Degree, Radian}
   \item [\texttt{ElectricCurrentUnits\_t}]
         \texttt{ElectricCurrentUnitsNull, ElectricCurrentUnitsUserDefined, Ampere, Abampere, Statampere, Edison,
         auCurrent}
   \item [\texttt{SubstanceAmountUnits\_t}]
         \texttt{SubstanceAmountUnitsNull, SubstanceAmountUnitsUserDefined, Mole, Entities, StandardCubicFoot,
         StandardCubicMeter}
   \item [\texttt{LuminousIntensityUnits\_t}]
         \texttt{LuminousIntensityUnitsNull, LuminousIntensityUnitsUserDefined, Candela, Candle, Carcel, Hefner, Violle}
   \item [\texttt{GoverningEquationsType\_t}]
         \texttt{GoverningEquationsTypeNull, GoverningEquationsTypeUserDefined, FullPotential, Euler, NSLaminar,
         NSTurbulent, NSLaminarIncompressible,
         NSTurbulentIncompressible}
   \item [\texttt{ModelType\_t}]
         \texttt{ModelTypeNull, ModelTypeUserDefined, Ideal, VanderWaals, Constant,
         PowerLaw, SutherlandLaw, ConstantPrandtl, EddyViscosity,
         ReynoldsStress, ReynoldsStressAlgebraic,
         Algebraic\_BaldwinLomax, Algebraic\_CebeciSmith,
         HalfEquation\_JohnsonKing, OneEquation\_BaldwinBarth,
         OneEquation\_SpalartAllmaras, TwoEquation\_JonesLaunder,
         TwoEquation\_MenterSST, TwoEquation\_Wilcox,
	 CaloricallyPerfect, ThermallyPerfect,
	 ConstantDensity, RedlichKwong,
	 Frozen, ThermalEquilib, ThermalNonequilib,
	 ChemicalEquilibCurveFit, ChemicalEquilibMinimization,
	 ChemicalNonequilib, EMElectricField, EMMagneticField,
         EMConductivity, Voltage, Interpolated,
         Equilibrium\_LinRessler, Chemistry\_LinRessler}
   \item [\texttt{GridLocation\_t}]
         \texttt{GridLocationNull, GridLocationUserDefined, Vertex, IFaceCenter, CellCenter, JFaceCenter,
         FaceCenter, KFaceCenter, EdgeCenter}
   \item [\texttt{GridConnectivityType\_t}]
         \texttt{GridConnectivityTypeNull, GridConnectivityTypeUserDefined, Overset, Abutting, Abutting1to1}
   \item [\texttt{PointSetType\_t}]
         \texttt{PointSetTypeNull, PointSetTypeUserDefined, PointList, PointRange, PointListDonor, PointRangeDonor,
         ElementList, ElementRange, CellListDonor}
   \item [\texttt{BCType\_t}]
         \texttt{BCTypeNull, BCTypeUserDefined, BCAxisymmetricWedge,
         BCDegenerateLine, BCExtrapolate,
         BCDegeneratePoint, BCDirichlet, BCFarfield, BCNeumann,
         BCGeneral, BCInflow, BCOutflow, BCInflowSubsonic,
         BCOutflowSubsonic, BCInflowSupersonic,
         BCOutflowSupersonic, BCSymmetryPlane, BCTunnelInflow,
         BCSymmetryPolar, BCTunnelOutflow, BCWallViscous,
         BCWall, BCWallViscousHeatFlux, BCWallInviscid,
         BCWallViscousIsothermal, FamilySpecified}
   \item [\texttt{BCDataType\_t}]
         \texttt{BCDataTypeNull, BCDataTypeUserDefined, Dirichlet, Neumann}
   \item [\texttt{RigidGridMotionType\_t}]
         \texttt{RigidGridMotionTypeNull, RigidGridMotionTypeUserDefined, ConstantRate, VariableRate}
   \item [\texttt{ArbitraryGridMotionType\_t}]
         \texttt{ArbitraryGridMotionTypeNull, ArbitraryGridMotionTypeUserDefined, NonDeformingGrid, DeformingGrid}
   \item [\texttt{SimulationType\_t}]
         \texttt{SimulationTypeNull, SimulationTypeUserDefined, TimeAccurate, NonTimeAccurate}
   \item [\texttt{WallFunctionType\_t}]
         \texttt{WallFunctionTypeNull, WallFunctionTypeUserDefined, Generic}
   \item [\texttt{AreaType\_t}]
         \texttt{AreaTypeNull, AreaTypeUserDefined, BleedArea, CaptureArea}
   \item [\texttt{AverageInterfaceType\_t}]
         \texttt{AverageInterfaceTypeNull, AverageInterfaceTypeUserDefined, AverageAll, AverageCircumferential, AverageRadial, AverageI,
	 AverageJ, AverageK}
\end{Ventryi}
}

\subsection{Character Names For Typedefs}

The CGNS library defines character arrays which map the typedefs above
to character strings.
These are global arrays dimensioned to the size of each list of
typedefs.
To retrieve a character string representation of a typedef, use the
typedef value as an index to the appropiate character array.
For example, to retrieve the string ``\texttt{Meter}''
for the \texttt{LengthUnits\_t Meter} typedef, use
\texttt{LengthUnitsName[Meter]}.
Functions are available to retrieve these names without the need for
direct global data access.
These functions also do bounds checking on the input, and if out of
range, will return the string ``\texttt{<invalid>}''.
An additional benefit is that these will work from within a Windows DLL,
and are thus the recommended access technique.
The routines have the same name as the global data arrays, but with a
``\texttt{cg\_}'' prepended.
For the example above, use ``\texttt{cg\_LengthUnitsName(Meter)}''.

\begin{table}
\begin{center}
\begin{tabular}{l}
\hline
Typedef Name Access Functions \\
\hline
\textcolor{output}{const char *name} = cg\_MassUnitsName(\textcolor{input}{MassUnits\_t type}); \\
\textcolor{output}{const char *name} = cg\_LengthUnitsName(\textcolor{input}{LengthUnits\_t type}); \\
\textcolor{output}{const char *name} = cg\_TimeUnitsName(\textcolor{input}{TimeUnits\_t type}); \\
\textcolor{output}{const char *name} = cg\_TemperatureUnitsName(\textcolor{input}{TemperatureUnits\_t type}); \\
\textcolor{output}{const char *name} = cg\_ElectricCurrentUnitsName(\textcolor{input}{ElectricCurrentUnits\_t type}); \\
\textcolor{output}{const char *name} = cg\_SubstanceAmountUnitsName(\textcolor{input}{SubstanceAmountUnits\_t type}); \\
\textcolor{output}{const char *name} = cg\_LuminousIntensityUnitsName(\textcolor{input}{LuminousIntensityUnits\_t type}); \\
\textcolor{output}{const char *name} = cg\_DataClassName(\textcolor{input}{DataClass\_t type}); \\
\textcolor{output}{const char *name} = cg\_GridLocationName(\textcolor{input}{GridLocation\_t type}); \\
\textcolor{output}{const char *name} = cg\_BCDataTypeName(\textcolor{input}{BCDataType\_t type}); \\
\textcolor{output}{const char *name} = cg\_GridConnectivityTypeName(\textcolor{input}{GridConnectivityType\_t type}); \\
\textcolor{output}{const char *name} = cg\_PointSetTypeName(\textcolor{input}{PointSetType\_t type}); \\
\textcolor{output}{const char *name} = cg\_GoverningEquationsTypeName(\textcolor{input}{GoverningEquationsType\_t type}); \\
\textcolor{output}{const char *name} = cg\_ModelTypeName(\textcolor{input}{ModelType\_t type}); \\
\textcolor{output}{const char *name} = cg\_BCTypeName(\textcolor{input}{BCType\_t type}); \\
\textcolor{output}{const char *name} = cg\_DataTypeName(\textcolor{input}{DataType\_t type}); \\
\textcolor{output}{const char *name} = cg\_ElementTypeName(\textcolor{input}{ElementType\_t type}); \\
\textcolor{output}{const char *name} = cg\_ZoneTypeName(\textcolor{input}{ZoneType\_t type}); \\
\textcolor{output}{const char *name} = cg\_RigidGridMotionTypeName(\textcolor{input}{RigidGridMotionType\_t type}); \\
\textcolor{output}{const char *name} = cg\_ArbitraryGridMotionTypeName(\textcolor{input}{ArbitraryGridMotionType\_t type}); \\
\textcolor{output}{const char *name} = cg\_SimulationTypeName(\textcolor{input}{SimulationType\_t type}); \\
\textcolor{output}{const char *name} = cg\_WallFunctionTypeName(\textcolor{input}{WallFunctionType\_t type}); \\
\textcolor{output}{const char *name} = cg\_AreaTypeName(\textcolor{input}{AreaType\_t type}); \\
\textcolor{output}{const char *name} = cg\_AverageInterfaceTypeName(\textcolor{input}{AverageInterfaceType\_t type}); \\
\hline
\end{tabular}
\end{center}
\end{table}
 
\clearpage

\subsection{64-bit C Portability and Issues}

If you use the \texttt{cgsize\_t} data type in new code,
it will work in both 32 and 64-bit compilation modes.
In order to support CGNS versions prior to 3.1, you may also
want to add something like this to your code:

\noindent \texttt{\#if CGNS\_VERSION < 3100}  \\
\noindent \texttt{\#define cgsize\_t int}  \\
\noindent \texttt{\#endif}

Existing code that uses \texttt{int} will not work with a
CGNS 3.1 library compiled in 64-bit mode. You may want
to add something like this to your code:

\noindent \texttt{\#if CGNS\_VERSION >= 3100 \&\& CG\_BUILD\_64BIT} \\
\noindent \texttt{\#error does not work in 64 bit mode} \\
\noindent \texttt{\#endif}

\noindent or modify your code to use \texttt{cgsize\_t}.

\subsection{64-bit Fortran Portability and Issues}

All integer arguments in the Fortran interface are taken to be
\texttt{integer*4} in 32-bit mode and \texttt{integer*8}
in 64-bit mode. If you have used default or implicit integers in
your Fortran code, it should port to 64-bit mode in most cases by
simply turning on your compiler option that promotes implicit
integers to \texttt{integer*8}. If you have explicitly defined your
integers as \texttt{integer*4}, your code will not work in 64-bit mode.
In that case, you will either need to change them to \texttt{integer}
(recommended for portability) or \texttt{integer*8}.

A new integer parameter has been added to the \textit{cgnslib\_f.h}
header, \texttt{CG\_BUILD\_64BIT}, which will be set
to 1 in 64-bit mode and 0 otherwise. You may use this parameter
to check at run time if the CGNS library has been compiled in 64-bit
mode or not, as in:

\noindent \texttt{if (CG\_BUILD\_64BIT .ne. 0) then} \\
\indent \texttt{print \*,'will not work in 64-bit mode'} \\
\indent \texttt{stop} \\
\noindent \texttt{endif}

If you are using a CGNS library prior to version 3.1, this parameter will
not be defined and you will need to rely on your compiler initializing all
undefined values to 0 (not always the case) for this test to work.

If your compiler supports automatic promotion of integers, and you use
implicit integers, your code should port to 64-bit with the following exception.

If you use an \texttt{Integer} data type in any routine that takes
a data type specification, and an implicit integer for the data, the code
will fail when compiled in 64-bit mode with automatic integer promotion.
An example of this would be:

\noindent \texttt{integer dim} \\
\noindent \texttt{integer data(dim)} \\
\noindent \texttt{call cg\_array\_write\_f('array',Integer,1,dim,data)} \\

This is because the MLL interprets the \texttt{Integer} data type
as \texttt{integer*4} regardless of the compilation mode. The compiler,
however, has automatically promoted \texttt{data} to be
\texttt{integer*8}. What you will need to do to prevent this problem,
is to either explicitly define \texttt{data} as in:

\noindent \texttt{integer dim} \\
\noindent \texttt{integer*4 data(dim)} \\
\noindent \texttt{call cg\_array\_write\_f('array',Integer,1,dim,data)} \\

\noindent or

\noindent \texttt{integer dim} \\
\noindent \texttt{integer*8 data(dim)} \\
\noindent \texttt{call cg\_array\_write\_f('array',LongInteger,1,dim,data)} \\

\noindent or test on \texttt{CG\_BUILD\_64BIT} as in:

\noindent \texttt{integer dim} \\
\noindent \texttt{integer data(dim)} \\
\noindent \texttt{if (CG\_BUILD\_64BIT .eq. 0) then} \\
\indent  \texttt{call cg\_array\_write\_f('array',Integer,1,dim,data)} \\
\noindent \texttt{else} \\
\indent  \texttt{call cg\_array\_write\_f('array',LongInteger,1,dim,data)} \\
\noindent \texttt{endif}

The last 2 options will only work with CGNS Version 3.1, since
\texttt{LongInteger} and \texttt{CG\_BUILD\_64BIT} are
not defined in previous versions.

You may also need to be careful when using integer constants as arguments
in 64-bit mode. If your compiler automatically promotes integer constants
to \texttt{integer*8}, then there is no problem. This is probably
the case if your compiler supports implicit integer promotion. If not,
then the constants will be \texttt{integer*4}, and your code
will not work in 64-bit mode. In that case you will need to do something like:

\noindent \texttt{integer*8 one,dim} \\
\noindent \texttt{integer*4 data(dim)} \\
\noindent \texttt{one = 1} \\
\noindent \texttt{call cg\_array\_write\_f('array',Integer,one,dim,data)}


\clearemptydoublepage
\section{File Operations}
\label{s:fileops}
\thispagestyle{plain}

\subsection{Opening and Closing a CGNS File}
\label{s:openclose}

\begin{fctbox}
\textcolor{output}{\textit{ier}} = cg\_open(\textcolor{input}{char *filename}, \textcolor{input}{int mode}, \textcolor{output}{\textit{int *fn}}); & r w m \\
\textcolor{output}{\textit{ier}} = cg\_version(\textcolor{input}{int fn}, \textcolor{output}{\textit{float *version}});         & r w m \\
\textcolor{output}{\textit{ier}} = cg\_close(\textcolor{input}{int fn});                           & r w m \\
\textcolor{output}{\textit{ier}} = cg\_is\_cgns(\textcolor{input}{const char *filename}, \textcolor{output}{\textit{int *file\_type}});                           & r w m \\
\textcolor{output}{\textit{ier}} = cg\_save\_as(\textcolor{input}{int fn}, \textcolor{input}{const char *filename}, \textcolor{input}{int file\_type}, \textcolor{input}{int follow\_links});                           & r w m \\
\textcolor{output}{\textit{ier}} = cg\_set\_file\_type(\textcolor{input}{int file\_type});                           & r w m \\
\textcolor{output}{\textit{ier}} = cg\_get\_file\_type(\textcolor{input}{int fn}, \textcolor{output}{\textit{int *file\_type}});                           & r w m \\
\hline
\hline
call cg\_open\_f(\textcolor{input}{filename}, \textcolor{input}{mode}, \textcolor{output}{\textit{fn}}, \textcolor{output}{\textit{ier}})          & r w m \\
call cg\_version\_f(\textcolor{input}{fn}, \textcolor{output}{\textit{version}}, \textcolor{output}{\textit{ier}})              & r w m \\
call cg\_close\_f(\textcolor{input}{fn}, \textcolor{output}{\textit{ier}})                         & r w m \\
call cg\_is\_cgns\_f(\textcolor{input}{filename}, \textcolor{output}{\textit{file\_type}}, \textcolor{output}{\textit{ier}})                         & r w m \\
call cg\_save\_as\_f(\textcolor{input}{fn}, \textcolor{input}{filename}, \textcolor{input}{file\_type}, \textcolor{input}{follow\_links}, \textcolor{output}{\textit{ier}})                         & r w m \\
call cg\_set\_file\_type\_f(\textcolor{input}{file\_type}, \textcolor{output}{\textit{ier}})                         & r w m \\
call cg\_get\_file\_type\_f(\textcolor{input}{fn}, \textcolor{output}{\textit{file\_type}}, \textcolor{output}{\textit{ier}})                         & r w m \\
\end{fctbox}

\noindent
\textbf{\textcolor{input}{Input}/\textcolor{output}{\textit{Output}}}

\begin{Ventryi}{\texttt{filename}}\raggedright
\item [\texttt{filename}]
      Name of the CGNS file, including path name if necessary.
      There is no limit on the length of this character variable.
      (\textcolor{input}{Input})
\item [\texttt{mode}]
      Mode used for opening the file.
      The modes currently supported are \texttt{CG\_MODE\_READ},
      \texttt{CG\_MODE\_WRITE}, and \texttt{CG\_MODE\_MODIFY}.
      (\textcolor{input}{Input})
\item [\texttt{fn}]
      CGNS file index number.
      (\textcolor{input}{Input} for \texttt{cg\_version} and
      \texttt{cg\_close}; \textcolor{output}{\textit{output}} for
      \texttt{cg\_open})
\item [\texttt{version}]
      CGNS version number.
      (\textcolor{output}{\textit{Output}})
\item [\texttt{file\_type}]
      Type of CGNS file.  This will typically be either \texttt{CG\_FILE\_ADF} or
      \texttt{CG\_FILE\_HDF5} depending on the underlying file format.
      (\textcolor{input}{Input} for \texttt{cg\_save\_as} and
      \texttt{cg\_set\_file\_type}; \textcolor{output}{\textit{output}} for
      \texttt{cg\_get\_file\_type})  \\
      However, note that when built in 32-bit, there is also an option
      to create a Version 2.5 CGNS file by setting the file type to
      \texttt{CG\_FILE\_ADF2}.
\item [\texttt{follow\_links}]
      This flag determines whether links are left intact when saving a CGNS 
      file. If non-zero, then the links will be removed and the data associated 
      with the linked files copied to the new file.
      (\textcolor{input}{Input})
\item [\texttt{ier}]
      Error status.
      (\textcolor{output}{\textit{Output}})
\end{Ventryi}

The function \texttt{cg\_open} must always be the first one called.
It opens a CGNS file for reading and/or writing and returns an index
number \texttt{fn}.
The index number serves to identify the CGNS file in subsequent
function calls.
Several CGNS files can be opened simultaneously.
The current limit on the number of files opened at once depends on the
platform.
On an SGI workstation, this limit is set at 100 (parameter
\texttt{FOPEN\_MAX} in \textit{stdio.h}).

The file can be opened in one of the following modes:

\begin{Ventryic}{\texttt{CG\_MODE\_MODIFY}}
\item [\texttt{CG\_MODE\_READ}]
      Read only mode.
\item [\texttt{CG\_MODE\_WRITE}]
      Write only mode.
\item [\texttt{CG\_MODE\_MODIFY}]
      Reading and/or writing is allowed.
\end{Ventryic}

When the file is opened, if no \texttt{CGNSLibraryVersion\_t} node is
found, a default value of 1.05 is assumed for the CGNS version number.
Note that this corresponds to an old version of the CGNS standard, that
doesn't include many data structures supported by the current standard.

The function \texttt{cg\_close} must always be the last one called.
It closes the CGNS file designated by the index number \texttt{fn} and
frees the memory where the CGNS data was kept.
When a file is opened for writing, \texttt{cg\_close} writes all the
CGNS data in memory onto disk prior to closing the file.
Consequently, if is omitted, the CGNS file is not written properly.

In order to reduce memory usage and improve execution speed,
large arrays such as grid coordinates or flow solutions are not actually
stored in memory.
Instead, only their ADF (or HDF) ID numbers are kept in memory for future
reference.
An attempt is also made to do the same with unstructured mesh element data.

The function \texttt{cg\_is\_cgns} may be used to determine if a file is a 
CGNS file or not, and the type of file (\texttt{CG\_FILE\_ADF} or \texttt{CG\_FILE\_HDF5}). 
If the file is a CGNS file, \texttt{cg\_is\_cgns} returns \texttt{CG\_OK}, 
otherwise \texttt{CG\_ERROR} is returned and \texttt{file\_type} is set to 
\texttt{CG\_FILE\_NONE}. 

The CGNS file identified by \texttt{fn} may be saved to a different filename 
and type using \texttt{cg\_save\_as}. In order to save as an HDF5 file, the 
library must have been built with HDF5 support. ADF support is always built.
The function 
\texttt{cg\_set\_file\_type} changes this default. The function 
\texttt{cg\_get\_file\_type} returns the file type for the CGNS file identified by 
\texttt{fn}. 

\subsection{Configuring CGNS Internals}
\label{s:configure}

\begin{fctbox}
\textcolor{output}{\textit{ier}} = cg\_configure(\textcolor{input}{int option}, \textcolor{input}{void *value}); & r w m \\
\textcolor{output}{\textit{ier}} = cg\_error\_handler(\textcolor{input}{void (*)(int, char *)}); & r w m \\
\textcolor{output}{\textit{ier}} = cg\_set\_compress(\textcolor{input}{int compress}); & r w m \\
\textcolor{output}{\textit{ier}} = cg\_get\_compress(\textcolor{output}{\textit{int *compress}}); & r w m \\
\textcolor{output}{\textit{ier}} = cg\_set\_path(\textcolor{input}{const char *path}); & r w m \\
\textcolor{output}{\textit{ier}} = cg\_add\_path(\textcolor{input}{const char *path}); & r w m \\
\hline
\hline
call cg\_set\_compress\_f(\textcolor{input}{compress}, \textcolor{output}{\textit{ier}})          & r w m \\
call cg\_get\_compress\_f(\textcolor{output}{\textit{compress}}, \textcolor{output}{\textit{ier}})          & r w m \\
call cg\_set\_path\_f(\textcolor{input}{path}, \textcolor{output}{\textit{ier}})          & r w m \\
call cg\_add\_path\_f(\textcolor{input}{path}, \textcolor{output}{\textit{ier}})          & r w m \\
\end{fctbox}

\noindent
\textbf{\textcolor{input}{Input}/\textcolor{output}{\textit{Output}}}

\begin{Ventryi}{\texttt{option}}\raggedright
\item [\texttt{option}]
      The option to configure, currently one of
      \texttt{CG\_CONFIG\_ERROR}, \texttt{CG\_CONFIG\_COMPRESS},
      \texttt{CG\_CONFIG\_SET\_PATH}, \texttt{CG\_CONFIG\_ADD\_PATH},
      or \texttt{CG\_CONFIG\_HDF5\_COMPRESS}
      as defined in \textit{cgnslib.h}.
      (\textcolor{input}{Input})
\item [\texttt{value}]
      The value to set, type cast as \texttt{void *}.
      (\textcolor{input}{Input})
\item [\texttt{compress}]
      CGNS compress (rewrite) setting).
      (\textcolor{input}{Input} for \texttt{cg\_set\_compress}; \textcolor{output}{\textit{output}} for
      \texttt{cg\_get\_compress})
\item [\texttt{path}]
      Pathname to search for linked to files when opening a file with external links.
      (\textcolor{input}{Input})
\item [\texttt{ier}]
      Error status.
      (\textcolor{output}{\textit{Output}})
\end{Ventryi}

The function \texttt{cg\_configure} allows certain CGNS library internal
options to be configured.
The currently supported options and expected values are:

\begin{Ventryi}{\texttt{CG\_CONFIG\_SET\_PATH}}
\item [\texttt{CG\_CONFIG\_ERROR}]
      This allows an error call-back function to be defined by the
      user.
      The value should be a pointer to a function to receive the error.
      The function is defined as \texttt{void err\_callback(int is\_error,
      char *errmsg)}, and will be called for errors and warnings.
      The first argument, \texttt{is\_error}, will be 0 for warning
      messages, 1 for error messages, and $-1$ if the program is
      going to terminate (i.e., a call to \texttt{cg\_error\_exit()}).
      The second argument is the error or warning message.
      If this is defined, warning and error messages will go to the
      function, rather than the terminal.
      A \texttt{value} of \texttt{NULL} will remove the call-back
      function.
\item [\texttt{CG\_CONFIG\_COMPRESS}]
      When a CGNS file is closed after being opened in modify mode, the
      normal operation of the CGNS library is to rewrite the file if
      there is unused space.
      This happens when nodes have been rewritten or deleted.
      Setting \texttt{value} to 0 will prevent the library from
      rewriting the file, and setting it to 1 will force the rewrite.
      The default value is $-1$.
\item [\texttt{CG\_CONFIG\_SET\_PATH}]
      Sets the search path for locating linked-to files.
      The argument \texttt{value} should be a character string
      containing one or more directories, formatted the same as for the
      \texttt{PATH} environment variable.
      This will replace any current settings.
      Setting \texttt{value} to \texttt{NULL} will remove all paths.
\item [\texttt{CG\_CONFIG\_ADD\_PATH}]
      Adds a directory, or list of directories, to the linked-to file
      search path.
      This is the same as \texttt{CG\_CONFIG\_SET\_PATH}, but adds to
      the path instead of replacing it.
\item [\texttt{CG\_CONFIG\_HDF5\_COMPRESS}]
      Sets the compression level for data written from HDF5. The default 
      is no compression. Setting \texttt{value} to $-1$, will use the default 
      compression level of $6$. The acceptable values are $0$ to $9$, 
      corresponding to gzip compression levels.
\end{Ventryi}

The routines \texttt{cg\_error\_handler}, \texttt{cg\_set\_compress}, 
\texttt{cg\_set\_path}, and \texttt{cg\_add\_path} are convenience functions 
built on top of \texttt{cg\_configure}.

There is no Fortran counterpart to function \texttt{cg\_configure} or 
\texttt{cg\_error\_handler}. 


\clearemptydoublepage
\section{Navigating a CGNS File}
\label{s:navigating}
\thispagestyle{plain}

\subsection{Accessing a Node}
\label{s:goto}

\begin{fctbox}
\textcolor{output}{\textit{ier}} = cg\_goto(\textcolor{input}{int fn}, \textcolor{input}{int B}, ..., \textcolor{input}{"end"}); & r w m \\
\textcolor{output}{\textit{ier}} = cg\_gorel(\textcolor{input}{int fn}, ..., \textcolor{input}{"end"}); & r w m \\
\textcolor{output}{\textit{ier}} = cg\_gopath(\textcolor{input}{int fn}, \textcolor{input}{const char *path}); & r w m \\
\textcolor{output}{\textit{ier}} = cg\_golist(\textcolor{input}{int fn}, \textcolor{input}{int B}, \textcolor{input}{int depth}, \textcolor{input}{char **label}, \textcolor{input}{int *index}); & r w m \\
\textcolor{output}{\textit{ier}} = cg\_gowhere(\textcolor{input}{int *fn}, \textcolor{input}{int *B}, \textcolor{output}{\textit{int *depth}}, \textcolor{output}{\textit{char **label}}, & r w m \\
~~~~~~~~~~~~~~~~~\textcolor{output}{\textit{int *index}}); & \\
\hline
call cg\_goto\_f(\textcolor{input}{fn}, \textcolor{input}{B}, \textcolor{output}{\textit{ier}}, ..., \textcolor{input}{'end'}) & r w m \\
call cg\_gorel\_f(\textcolor{input}{fn}, \textcolor{output}{\textit{ier}}, ..., \textcolor{input}{'end'}) & r w m \\
call cg\_gopath\_f(\textcolor{input}{fn}, \textcolor{input}{path}, \textcolor{output}{\textit{ier}}) & r w m \\
\end{fctbox}

\noindent
\textbf{\textcolor{input}{Input}/\textcolor{output}{\textit{Output}}}

\noindent (Note that for Fortran calls, all integer arguments are integer*4 in 32-bit mode and integer*8 in 64-bit mode.
See ``64-bit Fortran Portability and Issues" section.)

\begin{Ventryi}{\texttt{depth}}\raggedright
\item [\texttt{fn}]
      CGNS file index number.
      (\textcolor{input}{Input})
\item [\texttt{B}]
      Base index number, where $1 \leq \text{\texttt{B}} \leq \text{\texttt{nbases}}$.
      (\textcolor{input}{Input})
\item [\texttt{...}]
      Variable argument list used to specify the path to a node.
      It is composed of an unlimited list of pair-arguments
      identifying each node in the path.
      Nodes may be identified by their label or name.
      Thus, a pair-argument may be of the form
\begin{alltt}
   "CGNS\_NodeLabel", \textit{NodeIndex}
\end{alltt}
      where \texttt{CGNS\_NodeLabel} is the node label and
      \textit{NodeIndex} is the node index, or
\begin{alltt}
   "CGNS\_NodeName", 0
\end{alltt}
      where \texttt{CGNS\_NodeName} is the node name.
      The 0 in the second form is required, to indicate that a node
      name is being specified rather than a node label.
      In addition, a pair-argument may be specified as
\begin{alltt}
   "..", 0
\end{alltt}
      indicating the parent of the current node.
      The different pair-argument forms may be intermixed in the same
      function call.

      There is one exception to this rule.
      When accessing a \texttt{BCData\_t} node, the index must be set to
      either \texttt{Dirichlet} or \texttt{Neumann} since only these
      two types are allowed.
      (Note that \texttt{Dirichlet} and \texttt{Neumann} are defined in
      the include files \textit{cgnslib.h} and \textit{cgnslib\_f.h}).
      Since "\texttt{Dirichlet}" and "\texttt{Neuman}" are also the
      names for these nodes, you may also use the \texttt{"Dirichlet", 0}
      or \texttt{"Neuman", 0} to access the node.
      See the example below.
      (\textcolor{input}{Input})
\item [\texttt{end}]
      The character string \texttt{"end"} (or \texttt{'end'} for the Fortran
      function) must be the last argument.
      It is used to indicate the end of the argument list.
      You may also use the empty string, \texttt{""} (\texttt{{'}{'}} for
      Fortran), or the \texttt{NULL} string in C, to terminate the
      list.
      (\textcolor{input}{Input})
\item [\texttt{path}]
      The pathname for the node to go to.
      If a position has been already set, this may be a relative
      path, otherwise it is an absolute path name, starting with
      \texttt{"/Basename"}, where \texttt{Basename} is the base under
      which you wish to move.
      (\textcolor{input}{Input})
\item [\texttt{depth}]
      Depth of the path list.
      The maximum depth is defined in \textit{cgnslib.h} by
      \texttt{CG\_MAX\_GOTO\_DEPTH}, and is currently equal to 20.
      (\textcolor{input}{Input} for \texttt{cg\_golist};
      \textcolor{output}{\textit{output}} for \texttt{cg\_gowhere})
\item [\texttt{label}]
      Array of node labels for the path.
      This argument may be passed as \texttt{NULL} to
      \texttt{cg\_where()}, otherwise it must be dimensioned by the
      calling program.
      The maximum size required is \texttt{label[MAX\_GO\_TO\_DEPTH][33]}.
      You may call \texttt{cg\_where()} with both \texttt{label} and
      \texttt{index} set to \texttt{NULL} in order to get the current
      depth, then dimension to that value.
      (\textcolor{input}{Input} for \texttt{cg\_golist};
      \textcolor{output}{\textit{output}} for \texttt{cg\_gowhere})
\item [\texttt{index}]
      Array of node indices for the path.
      This argument may be passed as \texttt{NULL} to
      \texttt{cg\_where()}, otherwise it must be dimensioned by the
      calling program.
      The maximum size required is \texttt{index[MAX\_GO\_TO\_DEPTH]}.
      You may call \texttt{cg\_where()} with both \texttt{label} and
      \texttt{index} set to \texttt{NULL} in order to get the current
      depth, then dimension to that value.
      (\textcolor{input}{Input} for \texttt{cg\_golist};
      \textcolor{output}{\textit{output}} for \texttt{cg\_gowhere})
\item [\texttt{ier}]
      Error status.
      The possible values, with the corresponding C names (or Fortran
      parameters) defined in \textit{cgnslib.h} (or \textit{cgnslib\_f.h})
      are listed below.

      \begin{tabular}{@{}c l}
         \uline{\ital{Value}} & \uline{\ital{Name/Parameter}} \\
         0 & \texttt{CG\_OK} \\
         1 & \texttt{CG\_ERROR} \\
         2 & \texttt{CG\_NODE\_NOT\_FOUND} \\
         3 & \texttt{CG\_INCORRECT\_PATH}
      \end{tabular}

      For non-zero values, an error message may be printed using
      \texttt{cg\_error\_print()}, as described in \autoref{s:error}.
      (\textcolor{output}{\textit{Output}})
\end{Ventryi}

This function allows access to any parent-type nodes in a CGNS file.
A parent-type node is one that can have children.
Nodes that cannot have children, like \texttt{Descriptor\_t}, are
not supported by this function.

\noindent
\uline{\textit{Examples}}

To illustrate the use of the above routines, assume you have a file with
CGNS index number \texttt{filenum}, a base node named \texttt{Base}
with index number \texttt{basenum}, 2 zones (named \texttt{Zone1}
and \texttt{Zone2}, with indices 1 and 2), and user-defined data
(\texttt{User}, index 1) below each zone.
To move to the user-defined data node under zone 1, you may use any of
the following:
\begin{alltt}
   cg\_goto(filenum, basenum, "Zone\_t", 1, "UserDefinedData\_t", 1, NULL);
   cg\_goto(filenum, basenum "Zone1", 0, "UserDefinedData\_t", 1, NULL);
   cg\_goto(filenum, basenum, "Zone\_t", 1, "User", 0, NULL);
   cg\_goto(filenum, basenum, "Zone1", 0, "User", 0, NULL);
   cg\_gopath(filenum, "/Base/Zone1/User");
\end{alltt}
Now, to change to the user-defined data node under zone 2, you may use
the full path specification as above, or else a relative path, using one
of the following:
\begin{alltt}
   cg\_gorel(filenum, "..", 0, "..", 0, "Zone\_t", 2, "UserDefinedData\_t", 1, NULL);
   cg\_gorel(filenum, "..", 0, "..", 0, "Zone2", 0, "UserDefinedData\_t", 1, NULL);
   cg\_gorel(filenum, "..", 0, "..", 0, "Zone\_t", 2, "User", 0, NULL);
   cg\_gorel(filenum, "..", 0, "..", 0, "Zone2", 0, "User", 0, NULL);
   cg\_gopath(filenum, "../../Zone2/User");
\end{alltt}

Shown below are some additional examples of various uses of these
routines, in both C and Fortran, where \texttt{fn}, \texttt{B},
\texttt{Z}, etc., are index numbers.

\begin{alltt}
   ier = cg\_goto(fn, B, "Zone\_t", Z, "FlowSolution\_t", F, "..", 0, "MySolution",
                 0, "end");

   ier = cg\_gorel(fn, "..", 0, "FlowSolution\_t", F, NULL);

   ier = cg\_gopath(fn, "/MyBase/MyZone/MySolution");

   ier = cg\_gopath(fn, "../../MyZoneBC");

   call cg\_goto\_f(fn, B, ier, 'Zone\_t', Z, 'GasModel\_t', 1, 'DataArray\_t',
                  A, 'end')

   call cg\_goto\_f(fn, B, ier, 'Zone\_t', Z, 'ZoneBC\_t', 1, 'BC\_t', BC,
                  'BCDataSet\_t', S, 'BCData\_t', Dirichlet, 'end')

   call cg\_gorel\_f(fn, ier, '..', 0, 'Neumann', 0, '')

   call cg\_gopath\_f(fn, '../../MyZoneBC', ier)
\end{alltt}

\subsection{Deleting a Node}
\label{s:delete}

\begin{fctbox}
\textcolor{output}{\textit{ier}} = cg\_delete\_node(\textcolor{input}{char *NodeName}); & - - m \\
\hline
call cg\_delete\_node\_f(\textcolor{input}{NodeName}, \textcolor{output}{\textit{ier}}) & - - m \\
\end{fctbox}

\noindent
\textbf{\textcolor{input}{Input}/\textcolor{output}{\textit{Output}}}

\noindent (Note that for Fortran calls, all integer arguments are integer*4 in 32-bit mode and integer*8 in 64-bit mode.
See ``64-bit Fortran Portability and Issues" section.)

\begin{Ventryi}{\texttt{NodeName}}\raggedright
\item [\texttt{NodeName}]
      Name of the child to be deleted.
      (\textcolor{input}{Input})
\item [\texttt{ier}]
      Error status.
      (\textcolor{output}{\textit{Output}})
\end{Ventryi}

The function \texttt{cg\_delete\_node} is used is conjunction with
\texttt{cg\_goto}.
Once positioned at a parent node with \texttt{cg\_goto}, a child of
this node can be deleted with \texttt{cg\_delete\_node}.
This function requires a single argument, \texttt{NodeName}, which is
the name of the child to be deleted.

Since the highest level that can be pointed to with \texttt{cg\_goto} is
a base node for a CGNS database
(\texttt{CGNSBase\_t}), the
highest-level nodes that can be deleted are the children of a
\texttt{CGNSBase\_t} node.
In other words, nodes located directly under the ADF (or HDF) root node
(\texttt{CGNSBase\_t} and
\texttt{CGNSLibraryVersion\_t})
can not be deleted with \texttt{cg\_delete}.

A few other nodes are not allowed to be deleted from the database
because these are required nodes as defined by the SIDS, and deleting
them would make the file non-CGNS compliant.
These are:
\begin{itemize*}
\item Under \texttt{Zone\_t}:
      \texttt{ZoneType} 
\item Under \texttt{GridConnectivity1to1\_t}:
      \texttt{PointRange, PointRangeDonor, Transform}
\item Under \texttt{OversetHoles\_t}:
      \texttt{PointList} and any \texttt{IndexRange\_t}
\item Under \texttt{GridConnectivity\_t}:
      \texttt{PointRange, PointList, CellListDonor, PointListDonor}
\item Under \texttt{BC\_t}:
      \texttt{PointList, PointRange}
\item Under \texttt{GeometryReference\_t}:
      \texttt{GeometryFile, GeometryFormat}
\item Under \texttt{Elements\_t}:
      \texttt{ElementRange, ElementConnectivity}
\item Under \texttt{Gravity\_t}:
      \texttt{GravityVector}
\item Under \texttt{Axisymmetry\_t}:
      \texttt{AxisymmetryReferencePoint, AxisymmetryAxisVector}
\item Under \texttt{RotatingCoordinates\_t}:
      \texttt{RotationCenter, RotationRateVector}
\item Under \texttt{Periodic\_t}:
      \texttt{RotationCenter, RotationAngle, Translation}
\item Under \texttt{AverageInterface\_t}:
      \texttt{AverageInterfaceType}
\item Under \texttt{WallFunction\_t}:
      \texttt{WallFunctionType}
\item Under \texttt{Area\_t}:
      \texttt{AreaType, SurfaceArea, RegionName}
\end{itemize*}

When a child node is deleted, both the database and the file on disk are
updated to remove the node.
One must be careful not to delete a node from within a loop of that node
type.
For example, if the number of zones below a \texttt{CGNSBase\_t} node is
\texttt{nzones}, a zone should never be deleted from within a zone loop!
By deleting a zone, the total number of zones (\texttt{nzones}) changes,
as well as the zone indexing.
Suppose for example that \texttt{nzones} is 5, and that the third zone
is deleted.
After calling \texttt{cg\_delete\_node}, \texttt{nzones} is changed to 4,
and the zones originally indexed 4 and 5 are now indexed 3 and 4.


\clearemptydoublepage
\section{Error Handling}
\label{s:error}
\thispagestyle{plain}

\begin{fctbox}
\textcolor{output}{\textit{error\_message}} = const char *cg\_get\_error(); & r w m \\
void cg\_error\_exit(); & r w m \\
void cg\_error\_print(); & r w m \\
\hline
call cg\_get\_error\_f(\textcolor{output}{\textit{error\_message}}) & r w m \\
call cg\_error\_exit\_f() & r w m \\
call cg\_error\_print\_f() & r w m \\
\end{fctbox}

If an error occurs during the execution of a CGNS library function,
signified by a non-zero value of the error status variable
\texttt{ier}, an error message may be retrieved using the function
\texttt{cg\_get\_error}.
The function \texttt{cg\_error\_exit} may then be used to print the
error message and stop the execution of the program.
Alternatively, \texttt{cg\_error\_print} may be used to print the
error message and continue execution of the program.

In C, you may define a function to be called automatically in the case
of a warning or error using the \texttt{cg\_configure} routine.
The function is of the form \texttt{void err\_func(int is\_error, char
*errmsg)}, and will be called whenever an error or warning occurs.
The first argument, \texttt{is\_error}, will be 0 for warning messages,
1 for error messages, and $-1$ if the program is going to terminate
(i.e., a call to \texttt{cg\_error\_exit}).
The second argument is the error or warning message.


\clearemptydoublepage
\section{Structural Nodes}
\label{s:structural}
\thispagestyle{plain}

\subsection{CGNS Base Information}
\label{s:base}

\noindent
\textit{Node}: \fort{CGNSBase\_t}

\begin{fctbox}
\textcolor{output}{\textit{ier}} = cg\_base\_write(\textcolor{input}{int fn}, \textcolor{input}{char *basename}, \textcolor{input}{int cell\_dim}, \textcolor{input}{int phys\_dim}, & - w m \\
~~~~~~\textcolor{output}{\textit{int *B}}); & \\
\textcolor{output}{\textit{ier}} = cg\_nbases(\textcolor{input}{int fn}, \textcolor{output}{\textit{int *nbases}}); & r - m \\
\textcolor{output}{\textit{ier}} = cg\_base\_read(\textcolor{input}{int fn}, \textcolor{input}{int B}, \textcolor{output}{\textit{char *basename}}, \textcolor{output}{\textit{int *cell\_dim}}, & r - m \\
~~~~~~\textcolor{output}{\textit{int *phys\_dim}}); & \\
\hline
call cg\_base\_write\_f(\textcolor{input}{fn}, \textcolor{input}{basename}, \textcolor{input}{cell\_dim}, \textcolor{input}{phys\_dim}, \textcolor{output}{\textit{B}}, \textcolor{output}{\textit{ier}}) & - w m \\
call cg\_nbases\_f(\textcolor{input}{fn}, \textcolor{output}{\textit{nbases}}, \textcolor{output}{\textit{ier}}) & r - m \\
call cg\_base\_read\_f(\textcolor{input}{fn}, \textcolor{input}{B}, \textcolor{output}{\textit{basename}}, \textcolor{output}{\textit{cell\_dim}}, \textcolor{output}{\textit{phys\_dim}}, \textcolor{output}{\textit{ier}}) & r - m \\
\end{fctbox}

\noindent
\textbf{\textcolor{input}{Input}/\textcolor{output}{\textit{Output}}}

\begin{Ventryi}{\fort{basename}}\raggedright
\item [\fort{fn}]
      CGNS file index number.
      (\textcolor{input}{Input})
\item [\fort{B}]
      Base index number, where $1 \leq \text{\fort{B}} \leq \text{\fort{nbases}}$.
      (\textcolor{input}{Input} for \fort{cg\_base\_read};
      \textcolor{output}{\textit{output}} for \fort{cg\_base\_write})
\item [\fort{nbases}]
      Number of bases present in the CGNS file \fort{fn}.
      (\textcolor{output}{\textit{Output}})
\item [\fort{basename}]
      Name of the base.
      (\textcolor{input}{Input} for \fort{cg\_base\_write};
      \textcolor{output}{\textit{output}} for \fort{cg\_base\_read})
\item [\fort{cell\_dim}]
      Dimension of the cells; 3 for volume cells, 2 for surface cells.
      (\textcolor{input}{Input} for \fort{cg\_base\_write};
      \textcolor{output}{\textit{output}} for \fort{cg\_base\_read})
\item [\fort{phys\_dim}]
      Number of coordinates required to define a vector in the field.
      (\textcolor{input}{Input} for \fort{cg\_base\_write};
      \textcolor{output}{\textit{output}} for \fort{cg\_base\_read})
\item [\fort{ier}]
      Error status.
      (\textcolor{output}{\textit{Output}})
\end{Ventryi}

\subsection{Zone Information}
\label{s:zone}

\noindent
\textit{Node}: \fort{Zone\_t}

\begin{fctbox}
\textcolor{output}{\textit{ier}} = cg\_zone\_write(\textcolor{input}{int fn}, \textcolor{input}{int B}, \textcolor{input}{char *zonename}, \textcolor{input}{int *size}, & - w m \\
~~~~~~\textcolor{input}{ZoneType\_t zonetype}, \textcolor{output}{\textit{int *Z}}); & \\
\textcolor{output}{\textit{ier}} = cg\_nzones(\textcolor{input}{int fn}, \textcolor{input}{int B}, \textcolor{output}{\textit{int *nzones}}); & r - m \\
\textcolor{output}{\textit{ier}} = cg\_zone\_read(\textcolor{input}{int fn}, \textcolor{input}{int B}, \textcolor{input}{int Z}, \textcolor{output}{\textit{char *zonename}}, \textcolor{output}{\textit{int *size}}); & r - m \\
\textcolor{output}{\textit{ier}} = cg\_zone\_type(\textcolor{input}{int fn}, \textcolor{input}{int B}, \textcolor{input}{int Z}, \textcolor{output}{\textit{ZoneType\_t *zonetype}}); & r - m \\
\hline
call cg\_zone\_write\_f(\textcolor{input}{fn}, \textcolor{input}{B}, \textcolor{input}{zonename}, \textcolor{input}{size}, \textcolor{input}{zonetype}, \textcolor{output}{\textit{Z}}, \textcolor{output}{\textit{ier}}) & - w m \\
call cg\_nzones\_f(\textcolor{input}{fn}, \textcolor{input}{B}, \textcolor{output}{\textit{nzones}}, \textcolor{output}{\textit{ier}}) & r - m \\
call cg\_zone\_read\_f(\textcolor{input}{fn}, \textcolor{input}{B}, \textcolor{input}{Z}, \textcolor{output}{\textit{zonename}}, \textcolor{output}{\textit{size}}, \textcolor{output}{\textit{ier}}) & r - m \\
call cg\_zone\_type\_f(\textcolor{input}{fn}, \textcolor{input}{B}, \textcolor{input}{Z}, \textcolor{output}{\textit{zonetype}}, \textcolor{output}{\textit{ier}}) & r - m \\
\end{fctbox}

\noindent
\textbf{\textcolor{input}{Input}/\textcolor{output}{\textit{Output}}}

\begin{Ventryi}{\fort{zonename}}\raggedright
\item [\fort{fn}]
      CGNS file index number.
      (\textcolor{input}{Input})
\item [\fort{B}]
      Base index number, where $1 \leq \text{\fort{B}} \leq \text{\fort{nbases}}$.
      (\textcolor{input}{Input})
\item [\fort{Z}]
      Zone index number, where $1 \leq \text{\fort{Z}} \leq \text{\fort{nzones}}$.
      (\textcolor{input}{Input} for \fort{cg\_zone\_read},
      \fort{cg\_zone\_type}; \textcolor{output}{\textit{output}} for
      \fort{cg\_zone\_write})
\item [\fort{nzones}]
      Number of zones present in base \fort{B}.
      (\textcolor{output}{\textit{Output}})
\item [\fort{zonename}]
      Name of the zone.
      (\textcolor{input}{Input} for \fort{cg\_zone\_write};
      \textcolor{output}{\textit{output}} for \fort{cg\_zone\_read})
\item [\fort{size}]
      Number of vertices, cells, and boundary vertices in each
      (\textit{index})-dimension.

      Note that for unstructured grids, the number of cells is the
      number of highest order elements.
      Thus, in three dimensions it's the number of 3-D cells, and in
      two dimensions it's the number of 2-D cells.

      Also for unstructured grids, if the nodes are sorted between
      internal nodes and boundary nodes, the optional parameter
      \texttt{NBoundVertex} must be set equal to the number of boundary
      nodes.
      By default, \texttt{NBoundVertex} equals zero, meaning that the
      nodes are unsorted.

      Note that a non-zero value for \texttt{NBoundVertex} only applies
      to unstructured grids.
      For structured grids, the \texttt{NBoundVertex} parameter always
      equals 0 in all directions.

      \setlength{\tmplength}{\extrarowheight}	% temporarily decrease row heights
      \setlength{\extrarowheight}{0pt}
      \begin{tabular}{@{}l >{\ttfamily}l}
         \uline{\ital{Mesh Type}} & \uline{\ital{Size}} \\
         3D structured     & NVertexI, NVertexJ, NVertexK \\
                           & NCellI, NCellJ, NCellK \\
                           & NBoundVertexI $= 0$, NBoundVertexJ $= 0$, \\
                           & NBoundVertexK \\[0.75\baselineskip]
         2D structured     & NVertexI, NVertexJ \\
                           & NCellI, NCellJ \\
                           & NBoundVertexI $= 0$, NBoundVertexJ $= 0$ \\[0.75\baselineskip]
         3D unstructured   & NVertex, NCell3D, NBoundVertex \\[0.75\baselineskip]
         2D unstructured   & NVertex, NCell2D, NBoundVertex
      \end{tabular}
      \setlength{\extrarowheight}{\tmplength}

      (\textcolor{input}{Input} for \fort{cg\_zone\_write};
      \textcolor{output}{\textit{output}} for \fort{cg\_zone\_read})
\item [\fort{zonetype}]
      Type of the zone.
      The admissible types are \fort{Structured} and \fort{Unstructured}.
      (\textcolor{input}{Input} for \fort{cg\_zone\_write};
      \textcolor{output}{\textit{output}} for \fort{cg\_zone\_type})
\item [\fort{ier}]
      Error status.
      (\textcolor{output}{\textit{Output}})
\end{Ventryi}

Note that the zones are sorted alphanumerically to insure that
they can always be retrieved in the same order (for the same
model).

\subsection{Simulation Type}
\label{s:simulationtype}

\noindent
\textit{Node}: \fort{SimulationType\_t}

\begin{fctbox}
\textcolor{output}{\textit{ier}} = cg\_simulation\_type\_write(\textcolor{input}{int fn}, \textcolor{input}{int B}, & - w m \\
~~~~~~\textcolor{input}{SimulationType\_t SimulationType}); & \\
\textcolor{output}{\textit{ier}} = cg\_simulation\_type\_read(\textcolor{input}{int fn}, \textcolor{input}{int B}, & r - m \\
~~~~~~\textcolor{output}{\textit{SimulationType\_t SimulationType}}); & \\
\hline
call cg\_simulation\_type\_write\_f(\textcolor{input}{fn}, \textcolor{input}{B}, \textcolor{input}{SimulationType}, \textcolor{output}{\textit{ier}}) & - w m \\
call cg\_simulation\_type\_read\_f(\textcolor{input}{fn}, \textcolor{input}{B}, \textcolor{output}{\textit{SimulationType}}, \textcolor{output}{\textit{ier}}) & r - m \\
\end{fctbox}

\noindent
\textbf{\textcolor{input}{Input}/\textcolor{output}{\textit{Output}}}

\begin{Ventryi}{\fort{SimulationType}}\raggedright
\item [\fort{fn}]
      CGNS file index number.
      (\textcolor{input}{Input})
\item [\fort{B}]
      Base index number, where $1 \leq \text{\fort{B}} \leq \text{\fort{nbases}}$.
      (\textcolor{input}{Input})
\item [\fort{SimulationType}]
      Type of simulation.
      Valid types are \fort{Null}, \fort{UserDefined},
      \fort{TimeAccurate}, and \fort{NonTimeAccurate}.
      (\textcolor{input}{Input} for \fort{cg\_simulation\_type\_write};
      \textcolor{output}{\textit{output}} for \fort{cg\_simulation\_type\_read})
\item [\fort{ier}]
      Error status.
      (\textcolor{output}{\textit{Output}})
\end{Ventryi}


\clearemptydoublepage
\section{Descriptors}
\label{s:descriptors}
\thispagestyle{plain}

\subsection{Descriptive Text}
\label{s:descriptor}

\noindent
\textit{Node}: \fort{Descriptor\_t}

\begin{fctbox}
\textcolor{output}{\textit{ier}} = cg\_descriptor\_write(\textcolor{input}{char *name}, \textcolor{input}{char *text}); & - w m \\
\textcolor{output}{\textit{ier}} = cg\_ndescriptors(\textcolor{output}{\textit{int *ndescriptors}}); & r - m \\
\textcolor{output}{\textit{ier}} = cg\_descriptor\_read(\textcolor{input}{int D}, \textcolor{output}{\textit{char *name}}, \textcolor{output}{\textit{char **text}}); & r - m \\
\hline
call cg\_descriptor\_write\_f(\textcolor{input}{name}, \textcolor{input}{text}, \textcolor{output}{\textit{ier}}) & - w m \\
call cg\_ndescriptors\_f(\textcolor{output}{\textit{ndescriptors}}, \textcolor{output}{\textit{ier}}) & r - m \\
call cg\_descriptor\_size\_f(\textcolor{input}{D}, \textcolor{output}{\textit{size}}, \textcolor{output}{\textit{ier}}) & r - m \\
call cg\_descriptor\_read\_f(\textcolor{input}{D}, \textcolor{output}{\textit{name}}, \textcolor{output}{\textit{text}}, \textcolor{output}{\textit{ier}}) & r - m \\
\end{fctbox}

\noindent
\textbf{\textcolor{input}{Input}/\textcolor{output}{\textit{Output}}}

\begin{Ventryi}{\fort{ndescriptors}}\raggedright
\item [\fort{ndescriptors}]
      Number of \fort{Descriptor\_t} nodes under the current node.
      (\textcolor{output}{Output})
\item [\fort{D}]
      Descriptor index number, where $1 \leq \text{\fort{D}} \leq \text{\fort{ndescriptors}}$.
      (\textcolor{input}{Input})
\item [\fort{name}]
      Name of the \fort{Descriptor\_t} node.
      (\textcolor{input}{Input} for \fort{cg\_descriptor\_write};
      \textcolor{output}{\textit{output}} for \fort{cg\_descriptor\_read})
\item [\fort{text}]
      Description held in the \fort{Descriptor\_t} node.
      (\textcolor{input}{Input} for \fort{cg\_descriptor\_write};
      \textcolor{output}{\textit{output}} for \fort{cg\_descriptor\_read})
\item [\fort{size}]
      Size of the descriptor data (Fortran interface only).
      (\textcolor{output}{\textit{Output}})
\item [\fort{ier}]
      Error status.
      (\textcolor{output}{\textit{Output}})
\end{Ventryi}

Note that with \texttt{cg\_descriptor\_read} the memory for the descriptor
character string, \texttt{text}, will be allocated by the Mid-Level
Library.
The application code is responsible for releasing this memory when it is
no longer needed by calling \texttt{cg\_free(text)}, described in
\autoref{s:free}.

\subsection{Ordinal Value}
\label{s:ordinal}

\noindent
\textit{Node}: \fort{Ordinal\_t}

\begin{fctbox}
\textcolor{output}{\textit{ier}} = cg\_ordinal\_write(\textcolor{input}{int Ordinal}); & - w m \\
\textcolor{output}{\textit{ier}} = cg\_ordinal\_read(\textcolor{output}{\textit{int *Ordinal}}); & r - m \\
\hline
call cg\_ordinal\_write\_f(\textcolor{input}{Ordinal}, \textcolor{output}{\textit{ier}}) & - w m \\
call cg\_ordinal\_read\_f(\textcolor{output}{\textit{Ordinal}}, \textcolor{output}{\textit{ier}}) & r - m \\
\end{fctbox}

\noindent
\textbf{\textcolor{input}{Input}/\textcolor{output}{\textit{Output}}}

\begin{Ventryi}{\fort{Ordinal}}\raggedright
\item [\fort{Ordinal}]
      Any integer value.
      (\textcolor{input}{Input} for \fort{cg\_ordinal\_write};
      \textcolor{output}{\textit{output}} for \fort{cg\_ordinal\_read})
\item [\fort{ier}]
      Error status.
      (\textcolor{output}{\textit{Output}})
\end{Ventryi}


\clearemptydoublepage
\section{Physical Data}
\label{s:physical}
\thispagestyle{plain}

\subsection{Data Arrays}
\label{s:dataarray}

\noindent
\textit{Node}: \texttt{DataArray\_t}

\begin{fctbox}
\textcolor{output}{\textit{ier}} = cg\_array\_write(\textcolor{input}{char *ArrayName}, \textcolor{input}{DataType\_t DataType}, & - w m \\
~~~~~~\textcolor{input}{int DataDimension}, \textcolor{input}{cgsize\_t *DimensionVector}, \textcolor{input}{void *Data}); & \\
\textcolor{output}{\textit{ier}} = cg\_narrays(\textcolor{output}{\textit{int *narrays}}); & r - m \\
\textcolor{output}{\textit{ier}} = cg\_array\_info(\textcolor{input}{int A}, \textcolor{output}{\textit{char *ArrayName}}, \textcolor{output}{\textit{DataType\_t *DataType}}, & r - m \\
~~~~~~\textcolor{output}{\textit{int *DataDimension}}, \textcolor{output}{\textit{cgsize\_t *DimensionVector}}); & \\
\textcolor{output}{\textit{ier}} = cg\_array\_read(\textcolor{input}{int A}, \textcolor{output}{\textit{void *Data}}); & r - m \\
\textcolor{output}{\textit{ier}} = cg\_array\_read\_as(\textcolor{input}{int A}, \textcolor{input}{DataType\_t DataType}, \textcolor{output}{\textit{void *Data}}); & r - m \\
\hline
call cg\_array\_write\_f(\textcolor{input}{ArrayName}, \textcolor{input}{DataType}, \textcolor{input}{DataDimension}, & - w m \\
~~~~~\textcolor{input}{DimensionVector}, \textcolor{input}{Data}, \textcolor{output}{\textit{ier}}) & \\
call cg\_narrays\_f(\textcolor{output}{\textit{narrays}}, \textcolor{output}{\textit{ier}}) & r - m \\
call cg\_array\_info\_f(\textcolor{input}{A}, \textcolor{output}{\textit{ArrayName}}, \textcolor{output}{\textit{\textit{DataType}}}, \textcolor{output}{\textit{DataDimension}}, & r - m \\
~~~~~\textcolor{output}{\textit{DimensionVector}}, \textcolor{output}{\textit{ier}}) & \\
call cg\_array\_read\_f(\textcolor{input}{A}, \textcolor{output}{\textit{Data}}, \textcolor{output}{\textit{ier}}) & r - m \\
call cg\_array\_read\_as\_f(\textcolor{input}{A}, \textcolor{input}{DataType}, \textcolor{output}{\textit{Data}}, \textcolor{output}{\textit{ier}}) & r - m \\
\end{fctbox}

\noindent
\textbf{\textcolor{input}{Input}/\textcolor{output}{\textit{Output}}}

\noindent (Note that for Fortran calls, all integer arguments are integer*4 in 32-bit mode and integer*8 in 64-bit mode.
See ``64-bit Fortran Portability and Issues" section.)

\begin{Ventryi}{\texttt{DimensionVector}}\raggedright
\item [\texttt{narrays}]
      Number of \texttt{DataArray\_t} nodes under the current node.
      (\textcolor{output}{\textit{Output}})
\item [\texttt{A}]
      Data array index, where $1 \leq \text{\texttt{A}} \leq \text{\texttt{narrays}}$.
      (\textcolor{input}{Input})
\item [\texttt{ArrayName}]
      Name of the \texttt{DataArray\_t} node.
      (\textcolor{input}{Input} for \texttt{cg\_array\_write};
      \textcolor{output}{\textit{output}} for \texttt{cg\_array\_info})
\item [\texttt{DataType}]
      Type of data held in the \texttt{DataArray\_t} node.
      The admissible types are \texttt{Integer}, \texttt{RealSingle},
      \texttt{RealDouble}, and \texttt{Character}.
      (\textcolor{input}{Input} for \texttt{cg\_array\_write},
      \texttt{cg\_array\_read\_as}; \textcolor{output}{\textit{output}} for
      \texttt{cg\_array\_info})
\item [\texttt{DataDimension}]
      Number of dimensions.
      (\textcolor{input}{Input} for \texttt{cg\_array\_write};
      \textcolor{output}{\textit{output}} for \texttt{cg\_array\_info})
\item [\texttt{DimensionVector}]
      Number of data elements in each dimension.
      (\textcolor{input}{Input} for \texttt{cg\_array\_write};
      \textcolor{output}{\textit{output}} for \texttt{cg\_array\_info})
\item [\texttt{Data}]
      The data array.
      (\textcolor{input}{Input} for \texttt{cg\_array\_write};
      \textcolor{output}{\textit{output}} for \texttt{cg\_array\_read},
      \texttt{cg\_array\_read\_as})
\item [\texttt{ier}]
      Error status.
      (\textcolor{output}{\textit{Output}})
\end{Ventryi}

\newpage
\subsection{Data Class}
\label{s:dataclass}

\noindent
\textit{Node}: \texttt{DataClass\_t}

\begin{fctbox}
\textcolor{output}{\textit{ier}} = cg\_dataclass\_write(\textcolor{input}{DataClass\_t dataclass}); & - w m \\
\textcolor{output}{\textit{ier}} = cg\_dataclass\_read(\textcolor{output}{\textit{DataClass\_t *dataclass}}); & r - m \\
\hline
call cg\_dataclass\_write\_f(\textcolor{input}{dataclass}, \textcolor{output}{\textit{ier}}) & - w m \\
call cg\_dataclass\_read\_f(\textcolor{output}{\textit{dataclass}}, \textcolor{output}{\textit{ier}}) & r - m \\
\end{fctbox}

\noindent
\textbf{\textcolor{input}{Input}/\textcolor{output}{\textit{Output}}}

\noindent (Note that for Fortran calls, all integer arguments are integer*4 in 32-bit mode and integer*8 in 64-bit mode.
See ``64-bit Fortran Portability and Issues" section.)

\begin{Ventryi}{\texttt{dataclass}}\raggedright
\item [\texttt{dataclass}]
      Data class for the nodes at this level.
      See below for the data classes currently supported in CGNS.
      (\textcolor{input}{Input} for \texttt{cg\_dataclass\_write};
      \textcolor{output}{\textit{output}} for \texttt{cg\_dataclass\_read})
\item [\texttt{ier}]
      Error status.
      (\textcolor{output}{\textit{Output}})
\end{Ventryi}
The data classes currently supported in CGNS are:

\begin{Ventryi}{\texttt{NormalizedByUnknownDimensional}}\raggedright
\item [\texttt{Dimensional}]
      Regular dimensional data.
\item [\texttt{NormalizedByDimensional}]
      Nondimensional data that is normalized by dimensional reference
      quantities.
\item [\texttt{NormalizedByUnknownDimensional}]
      All fields and reference data are nondimensional.
\item [\texttt{NondimensionalParameter}]
      Nondimensional parameters such as Mach number and lift coefficient.
\item [\texttt{DimensionlessConstant}]
      Constant such as $\pi$.
\end{Ventryi}
These classes are declared within typedef \texttt{DataClass\_t} in
\textit{cgnslib.h}, and as parameters in \textit{cgnslib\_f.h}.

\subsection{Data Conversion Factors}
\label{s:dataconversion}

\noindent
\textit{Node}: \texttt{DataConversion\_t}

\begin{fctbox}
\textcolor{output}{\textit{ier}} = cg\_conversion\_write(\textcolor{input}{DataType\_t DataType}, & - w m \\
~~~~~~\textcolor{input}{void *ConversionFactors}); & \\
\textcolor{output}{\textit{ier}} = cg\_conversion\_info(\textcolor{output}{\textit{DataType\_t *DataType}}); & r - m \\
\textcolor{output}{\textit{ier}} = cg\_conversion\_read(\textcolor{output}{\textit{void *ConversionFactors}}); & r - m \\
\hline
call cg\_conversion\_write\_f(\textcolor{input}{DataType}, \textcolor{input}{ConversionFactors}, \textcolor{output}{\textit{ier}}) & - w m \\
call cg\_conversion\_info\_f(\textcolor{output}{\textit{DataType}}, \textcolor{output}{\textit{ier}}) & r - m \\
call cg\_conversion\_read\_f(\textcolor{output}{\textit{ConversionFactors}}, \textcolor{output}{\textit{ier}}) & r - m \\
\end{fctbox}

\noindent
\textbf{\textcolor{input}{Input}/\textcolor{output}{\textit{Output}}}

\noindent (Note that for Fortran calls, all integer arguments are integer*4 in 32-bit mode and integer*8 in 64-bit mode.
See ``64-bit Fortran Portability and Issues" section.)

\begin{Ventryi}{\texttt{ConversionOffset}}\raggedright
\item [\texttt{DataType}]
      Data type in which the conversion factors are recorded.
      Admissible data types for conversion factors are \texttt{RealSingle}
      and \texttt{RealDouble}.
      (\textcolor{input}{Input} for \texttt{cg\_conversion\_write};
      \textcolor{output}{\textit{output}} for \texttt{cg\_conversion\_info})
\item [\texttt{ConversionFactors}]
      Two-element array containing the scaling and offset factors.
      (\textcolor{input}{Input} for \texttt{cg\_conversion\_write};
      \textcolor{output}{\textit{output}} for \texttt{cg\_conversion\_read})
\item [\texttt{ier}]
      Error status.
      (\textcolor{output}{\textit{Output}})
\end{Ventryi}

The \texttt{DataConversion\_t} data structure contains factors to convert
the nondimensional data to ``raw'' dimensional data.
The scaling and offset factors are contained in the two-element array
\texttt{ConversionFactors}.
In pseudo-Fortran, the conversion process is as follows:

\begin{indlefttt}
ConversionScale  = ConversionFactors(1)
ConversionOffset = ConversionFactors(2)
Data(\textit{raw}) = Data(\textit{nondimensional})*ConversionScale + ConversionOffset
\end{indlefttt}

\subsection{Dimensional Units}
\label{s:dimensionalunits}

\noindent
\textit{Node}: \texttt{DimensionalUnits\_t}

\begin{fctbox}
\textcolor{output}{\textit{ier}} = cg\_units\_write(\textcolor{input}{MassUnits\_t mass}, \textcolor{input}{LengthUnits\_t length}, & - w m \\
~~~~~~\textcolor{input}{TimeUnits\_t time}, \textcolor{input}{TemperatureUnits\_t temperature}, & \\
~~~~~~\textcolor{input}{AngleUnits\_t angle}); & \\
\textcolor{output}{\textit{ier}} = cg\_unitsfull\_write(\textcolor{input}{MassUnits\_t mass}, \textcolor{input}{LengthUnits\_t length}, & - w m \\
~~~~~~\textcolor{input}{TimeUnits\_t time}, \textcolor{input}{TemperatureUnits\_t temperature}, & \\
~~~~~~\textcolor{input}{AngleUnits\_t angle}, \textcolor{input}{ElectricCurrentUnits\_t current}, & \\
~~~~~~\textcolor{input}{SubstanceAmountUnits\_t amount}, & \\
~~~~~~\textcolor{input}{LuminousIntensityUnits\_t intensity}); & \\
\textcolor{output}{\textit{ier}} = cg\_nunits(\textcolor{output}{\textit{int *nunits}}); & r - m \\
\textcolor{output}{\textit{ier}} = cg\_units\_read(\textcolor{output}{\textit{MassUnits\_t *mass}}, \textcolor{output}{\textit{LengthUnits\_t *length}}, & r - m \\
~~~~~~\textcolor{output}{\textit{TimeUnits\_t *time}}, \textcolor{output}{\textit{TemperatureUnits\_t *temperature}}, & \\
~~~~~~\textcolor{output}{\textit{AngleUnits\_t *angle}}); & \\
\textcolor{output}{\textit{ier}} = cg\_unitsfull\_read(\textcolor{output}{\textit{MassUnits\_t *mass}}, \textcolor{output}{\textit{LengthUnits\_t *length}}, & r - m \\
~~~~~~\textcolor{output}{\textit{TimeUnits\_t *time}}, \textcolor{output}{\textit{TemperatureUnits\_t *temperature}}, & \\
~~~~~~\textcolor{output}{\textit{AngleUnits\_t *angle}}, \textcolor{output}{\textit{ElectricCurrentUnits\_t *current}}, & \\
~~~~~~\textcolor{output}{\textit{SubstanceAmountUnits\_t *amount}}, & \\
~~~~~~\textcolor{output}{\textit{LuminousIntensityUnits\_t *intensity}}); & \\
\hline
call cg\_units\_write\_f(\textcolor{input}{mass}, \textcolor{input}{length}, \textcolor{input}{time}, \textcolor{input}{temperature}, \textcolor{input}{angle}, \textcolor{output}{\textit{ier}}) & - w m \\
call cg\_unitsfull\_write\_f(\textcolor{input}{mass}, \textcolor{input}{length}, \textcolor{input}{time}, \textcolor{input}{temperature}, \textcolor{input}{angle}, & - w m \\
~~~~~\textcolor{input}{current}, \textcolor{input}{amount}, \textcolor{input}{intensity}, \textcolor{output}{\textit{ier}}) & \\
call cg\_nunits\_f(\textcolor{output}{\textit{int *nunits}}) & r - m \\
call cg\_units\_read(\textcolor{output}{\textit{mass}}, \textcolor{output}{\textit{length}}, \textcolor{output}{\textit{time}}, \textcolor{output}{\textit{temperature}}, \textcolor{output}{\textit{angle}}, \textcolor{output}{\textit{ier}}) & r - m \\
call cg\_unitsfull\_read\_f(\textcolor{output}{\textit{mass}}, \textcolor{output}{\textit{length}}, \textcolor{output}{\textit{time}}, \textcolor{output}{\textit{temperature}}, \textcolor{output}{\textit{angle}}, & r - m \\
~~~~~\textcolor{output}{\textit{current}}, \textcolor{output}{\textit{amount}}, \textcolor{output}{\textit{intensity}}, \textcolor{output}{\textit{ier}}) & \\
\end{fctbox}

\noindent
\textbf{\textcolor{input}{Input}/\textcolor{output}{\textit{Output}}}

\noindent (Note that for Fortran calls, all integer arguments are integer*4 in 32-bit mode and integer*8 in 64-bit mode.
See ``64-bit Fortran Portability and Issues" section.)

\begin{Ventryi}{\texttt{temperature}}\raggedright
\item [\texttt{mass}]
      Mass units.
      Admissible values are \texttt{CG\_Null}, \texttt{CG\_UserDefined},
      \texttt{Kilogram}, \texttt{Gram}, \texttt{Slug}, and
      \texttt{PoundMass}.
      (\textcolor{input}{Input} for \texttt{cg\_units\_write},
      \texttt{cg\_unitsfull\_write};
      \textcolor{output}{\textit{output}} for \texttt{cg\_units\_read},
      \texttt{cg\_unitsfull\_read})
\item [\texttt{length}]
      Length units.
      Admissible values are \texttt{CG\_Null}, \texttt{CG\_UserDefined},
      \texttt{Meter}, \texttt{Centimeter}, \texttt{Millimeter},
      \texttt{Foot}, and \texttt{Inch}.
      (\textcolor{input}{Input} for \texttt{cg\_units\_write},
      \texttt{cg\_unitsfull\_write};
      \textcolor{output}{\textit{output}} for \texttt{cg\_units\_read},
      \texttt{cg\_unitsfull\_read})
\item [\texttt{time}]
      Time units.
      Admissible values are \texttt{CG\_Null}, \texttt{CG\_UserDefined}, and
      \texttt{Second}.
      (\textcolor{input}{Input} for \texttt{cg\_units\_write},
      \texttt{cg\_unitsfull\_write};
      \textcolor{output}{\textit{output}} for \texttt{cg\_units\_read},
      \texttt{cg\_unitsfull\_read})
\item [\texttt{temperature}]
      Temperature units.
      Admissible values are \texttt{CG\_Null}, \texttt{CG\_UserDefined},
      \texttt{Kelvin}, \texttt{Celsius}, \texttt{Rankine}, and
      \texttt{Fahrenheit}.
      (\textcolor{input}{Input} for \texttt{cg\_units\_write},
      \texttt{cg\_unitsfull\_write};
      \textcolor{output}{\textit{output}} for \texttt{cg\_units\_read},
      \texttt{cg\_unitsfull\_read})
\item [\texttt{angle}]
      Angle units.
      Admissible values are \texttt{CG\_Null}, \texttt{CG\_UserDefined},
      \texttt{Degree}, and \texttt{Radian}.
      (\textcolor{input}{Input} for \texttt{cg\_units\_write},
      \texttt{cg\_unitsfull\_write};
      \textcolor{output}{\textit{output}} for \texttt{cg\_units\_read},
      \texttt{cg\_unitsfull\_read})
\item [\texttt{current}]
      Electric current units.
      Admissible values are \texttt{CG\_Null}, \texttt{CG\_UserDefined},
      \texttt{Ampere}, \texttt{Abampere}, \texttt{Statampere},
      \texttt{Edison}, and \texttt{auCurrent}.
      (\textcolor{input}{Input} for \texttt{cg\_unitsfull\_write};
      \textcolor{output}{\textit{output}} for \texttt{cg\_unitsfull\_read})
\item [\texttt{amount}]
      Substance amount units.
      Admissible values are \texttt{CG\_Null}, \texttt{CG\_UserDefined},
      \texttt{Mole}, \texttt{Entities}, \texttt{StandardCubicFoot}, and
      \texttt{StandardCubicMeter}.
      (\textcolor{input}{Input} for \texttt{cg\_unitsfull\_write};
      \textcolor{output}{\textit{output}} for \texttt{cg\_unitsfull\_read})
\item [\texttt{intensity}]
      Luminous intensity units.
      Admissible values are \texttt{CG\_Null}, \texttt{CG\_UserDefined},
      \texttt{Candela}, \texttt{Candle}, \texttt{Carcel},
      \texttt{Hefner}, and \texttt{Violle}.
      (\textcolor{input}{Input} for \texttt{cg\_unitsfull\_write};
      \textcolor{output}{\textit{output}} for \texttt{cg\_unitsfull\_read})
\item [\texttt{nunits}]
      Number of units used in the file (i.e., either 5 or 8).
      (\textcolor{output}{\textit{Output}})
\item [\texttt{ier}]
      Error status.
      (\textcolor{output}{\textit{Output}})
\end{Ventryi}

The supported units are declared within typedefs in \textit{cgnslib.h}
and as parameters in \textit{cgnslib\_f.h}.

When reading units data, either \texttt{cg\_units\_read} or
\texttt{cg\_unitsfull\_read} may be used, regardless of the number of
units used in the file.
If \texttt{cg\_unitsfull\_read} is used, but only five units are used
in the file, the returned values of \texttt{current}, \texttt{amount},
and \texttt{intensity} will be \texttt{CG\_Null}.

\newpage
\subsection{Dimensional Exponents}
\label{s:dimensionalexponents}

\noindent
\textit{Node}: \texttt{DimensionalExponents\_t}

\begin{fctbox}
\textcolor{output}{\textit{ier}} = cg\_exponents\_write(\textcolor{input}{DataType\_t DataType}, \textcolor{input}{void *exponents}); & - w m \\
\textcolor{output}{\textit{ier}} = cg\_expfull\_write(\textcolor{input}{DataType\_t DataType}, \textcolor{input}{void *exponents}); & - w m \\
\textcolor{output}{\textit{ier}} = cg\_nexponents(\textcolor{output}{\textit{int *nexponents}}); & r - m \\
\textcolor{output}{\textit{ier}} = cg\_exponents\_info(\textcolor{output}{\textit{DataType\_t *DataType}}); & r - m \\
\textcolor{output}{\textit{ier}} = cg\_exponents\_read(\textcolor{output}{\textit{void *exponents}}); & r - m \\
\textcolor{output}{\textit{ier}} = cg\_expfull\_read(\textcolor{output}{\textit{void *exponents}}); & r - m \\
\hline
call cg\_exponents\_write\_f(\textcolor{input}{DataType}, \textcolor{input}{exponents}, \textcolor{output}{\textit{ier}}) & - w m \\
call cg\_expfull\_write\_f(\textcolor{input}{DataType}, \textcolor{input}{exponents}, \textcolor{output}{\textit{ier}}) & - w m \\
call cg\_nexponents\_f(\textcolor{output}{\textit{nexponents}}, \textcolor{output}{\textit{ier}} ) & r - m \\
call cg\_exponents\_info\_f(\textcolor{output}{\textit{DataType}}, \textcolor{output}{\textit{ier}} ) & r - m \\
call cg\_exponents\_read\_f(\textcolor{output}{\textit{exponents}}, \textcolor{output}{\textit{ier}}) & r - m \\
call cg\_expfull\_read\_f(\textcolor{output}{\textit{exponents}}, \textcolor{output}{\textit{ier}}) & r - m \\
\end{fctbox}

\noindent
\textbf{\textcolor{input}{Input}/\textcolor{output}{\textit{Output}}}

\noindent (Note that for Fortran calls, all integer arguments are integer*4 in 32-bit mode and integer*8 in 64-bit mode.
See ``64-bit Fortran Portability and Issues" section.)

\begin{Ventryi}{\texttt{exponents}}\raggedright
\item [\texttt{DataType}]
      Data type in which the exponents are recorded.
      Admissible data types for the exponents are \texttt{RealSingle}
      and \texttt{RealDouble}.
      (\textcolor{input}{Input} for \texttt{cg\_exponents\_write};
      \textcolor{output}{\textit{output}} for \texttt{cg\_exponents\_info})
\item [\texttt{exponents}]
      Exponents for the dimensional units for mass, length, time,
      temperature, angle, electric current, substance amount, and
      luminous intensity, in that order.
      (\textcolor{input}{Input} for \texttt{cg\_exponents\_write},
      \texttt{cg\_expfull\_write};
      \textcolor{output}{\textit{output}} for \texttt{cg\_exponents\_read},
      \texttt{cg\_expfull\_read})
\item [\texttt{nexponents}]
      Number of exponents used in the file (i.e., either 5 or 8).
      (\textcolor{output}{\textit{Output}})
\item [\texttt{ier}]
      Error status.
      (\textcolor{output}{\textit{Output}})
\end{Ventryi}

When reading exponent data, either \texttt{cg\_exponents\_read} or
\texttt{cg\_expfull\_read} may be used, regardless of the number of
exponents used in the file.
If \texttt{cg\_exponents\_read} is used, but all eight exponents are used
in the file, only the first five exponents are returned.
If \texttt{cg\_expfull\_read} is used, but only five exponents are used
in the file, the returned values of the exponents for electric current,
substance amount, and luminous intensity will be zero.


\clearemptydoublepage
\section{Location and Position}
\label{s:location}
\thispagestyle{plain}

\subsection{Grid Location}
\label{s:gridlocation}

\noindent
\textit{Node}: \fort{GridLocation\_t}

\begin{fctbox}
\textcolor{output}{\textit{ier}} = cg\_gridlocation\_write(\textcolor{input}{GridLocation\_t GridLocation}); & - w m \\
\textcolor{output}{\textit{ier}} = cg\_gridlocation\_read(\textcolor{output}{\textit{GridLocation\_t *GridLocation}}); & r - m \\
\hline
call cg\_gridlocation\_write\_f(\textcolor{input}{GridLocation}, \textcolor{output}{\textit{ier}}) & - w m \\
call cg\_gridlocation\_read\_f(\textcolor{output}{\textit{GridLocation}}, \textcolor{output}{\textit{ier}}) & r - m \\
\end{fctbox}

\noindent
\textbf{\textcolor{input}{Input}/\textcolor{output}{\textit{Output}}}

\begin{Ventryi}{\fort{GridLocation}}\raggedright
\item [\fort{GridLocation}]
      Location in the grid.
      The admissible locations are \fort{Null}, \fort{UserDefined},
      \fort{Vertex}, \fort{CellCenter}, \fort{FaceCenter},
      \fort{IFaceCenter}, \fort{JFaceCenter}, \fort{KFaceCenter},
      and \fort{EdgeCenter}.
      (\textcolor{input}{Input} for \fort{cg\_gridlocation\_write};
      \textcolor{output}{\textit{output}} for \fort{cg\_gridlocation\_read})
\item [\fort{ier}]
      Error status.
      (\textcolor{output}{\textit{Output}})
\end{Ventryi}

\subsection{Point Sets}
\label{s:ptset}

\noindent
\textit{Node}: \fort{IndexArray\_t}, \fort{IndexRange\_t}

\begin{fctbox}
\textcolor{output}{\textit{ier}} = cg\_ptset\_write(\textcolor{input}{PointSetType\_t *ptset\_type}, \textcolor{input}{int npnts}, \textcolor{input}{int *pnts}); & - w m \\
\textcolor{output}{\textit{ier}} = cg\_ptset\_info(\textcolor{output}{\textit{PointSetType\_t *ptset\_type}}, \textcolor{output}{\textit{int *npnts}}); & r - m \\
\textcolor{output}{\textit{ier}} = cg\_ptset\_read(\textcolor{output}{\textit{int *pnts}}); & r - m \\
\hline
call cg\_ptset\_write\_f(\textcolor{input}{ptset\_type}, \textcolor{input}{npnts}, \textcolor{input}{pnts}, \textcolor{output}{\textit{ier}}) & - w m \\
call cg\_ptset\_info\_f(\textcolor{output}{\textit{ptset\_type}}, \textcolor{output}{\textit{npnts}}, \textcolor{output}{\textit{ier}}) & r - m \\
call cg\_ptset\_read\_f(\textcolor{output}{\textit{pnts}}, \textcolor{output}{\textit{ier}}) & r - m \\
\end{fctbox}

\noindent
\textbf{\textcolor{input}{Input}/\textcolor{output}{\textit{Output}}}

\begin{Ventryi}{\fort{ptset\_type}}\raggedright
\item [\fort{ptset\_type}]
      The point set type; either \fort{PointRange} for a range of
      points or cells, or \fort{PointList} for a list of discrete
      points or cells.
      (\textcolor{input}{Input} for \fort{cg\_ptset\_write};
      \textcolor{output}{\textit{output}} for \fort{cg\_ptset\_info})
\item [\fort{npnts}]
      The number of points or cells in the point set.
      For a point set type of \fort{PointRange}, \fort{npnts} is
      always two.
      For a point set type of \fort{PointList}, \fort{npnts} is
      the number of points or cells in the list.
      (\textcolor{input}{Input} for \fort{cg\_ptset\_write};
      \textcolor{output}{\textit{output}} for \fort{cg\_ptset\_info})
\item [\fort{pnts}]
      The array of point or cell indices defining the point set.
      There should be \fort{npnts} values, each of dimension
      \fort{IndexDimension}
      (i.e., 1 for unstructured grids, and 2 or 3 for structured grids
      with 2-D or 3-D elements, respectively).
      (\textcolor{input}{Input} for \fort{cg\_ptset\_write};
      \textcolor{output}{\textit{output}} for \fort{cg\_ptset\_read})
\item [\fort{ier}]
      Error status.
      (\textcolor{output}{\textit{Output}})
\end{Ventryi}

These functions may be used to write and read point set data (i.e., an
\fort{IndexArray\_t} node named \fort{PointList}, or an
\fort{IndexRange\_t} node named \fort{PointRange}).
They are only applicable at nodes that are descendents of a
\fort{Zone\_t} node.

\subsection{Rind Layers}
\label{s:rind}

\noindent
\textit{Node}: \fort{Rind\_t}

\begin{fctbox}
\textcolor{output}{\textit{ier}} = cg\_rind\_write(\textcolor{input}{int *RindData}); & - w m \\
\textcolor{output}{\textit{ier}} = cg\_rind\_read(\textcolor{output}{\textit{int *RindData}}); & r - m \\
\hline
call cg\_rind\_write\_f(\textcolor{input}{RindData}, \textcolor{output}{\textit{ier}}) & - w m \\
call cg\_rind\_read\_f(\textcolor{output}{\textit{RindData}}, \textcolor{output}{\textit{ier}}) & r - m \\
\end{fctbox}

\noindent
\textbf{\textcolor{input}{Input}/\textcolor{output}{\textit{Output}}}

\begin{Ventryi}{\fort{RindData}}\raggedright
\item [\fort{RindData}]
      Number of rind layers for each computational direction (structured
      grid) or number of rind points or elements (unstructured grid).
      (\textcolor{input}{Input} for \fort{cg\_rind\_write};
      \textcolor{output}{\textit{output}} for \fort{cg\_rind\_read})
\item [\fort{ier}]
      Error status.
      (\textcolor{output}{\textit{Output}})
\end{Ventryi}

When writing rind data for elements, \texttt{cg\_section\_write} must be
called first (see \autoref{s:elements}), followed by \texttt{cg\_goto}
(\autoref{s:navigating}) to access the \texttt{Elements\_t} node, and
then \texttt{cg\_rind\_write}.


\clearemptydoublepage
\section{Auxiliary Data}
\label{s:auxiliary}
\thispagestyle{plain}

\subsection{Reference State}
\label{s:refstate}

\noindent
\textit{Node}: \fort{ReferenceState\_t}

\begin{fctbox}
\textcolor{output}{\textit{ier}} = cg\_state\_write(\textcolor{input}{char *StateDescription}); & - w m \\
\textcolor{output}{\textit{ier}} = cg\_state\_read(\textcolor{output}{\textit{char **StateDescription}}); & r - m \\
\hline
call cg\_state\_write\_f(\textcolor{input}{StateDescription}, \textcolor{output}{\textit{ier}}) & - w m \\
call cg\_state\_size\_f(\textcolor{output}{\textit{Size}}, \textcolor{output}{\textit{ier}}) & r - m \\
call cg\_state\_read\_f(\textcolor{output}{\textit{StateDescription}}, \textcolor{output}{\textit{ier}}) & r - m \\
\end{fctbox}

\noindent
\textbf{\textcolor{input}{Input}/\textcolor{output}{\textit{Output}}}

\begin{Ventryi}{\fort{StateDescription}}\raggedright
\item [\fort{StateDescription}]
      Text description of reference state.
      (\textcolor{input}{Input} for \fort{cg\_state\_write};
      \textcolor{output}{\textit{output}} for \fort{cg\_state\_read})
\item [\fort{Size}]
      Number of characters in the \fort{StateDescription} string
      (Fortran interface only).
      (\textcolor{output}{\textit{Output}})
\item [\fort{ier}]
      Error status.
      (\textcolor{output}{\textit{Output}})
\end{Ventryi}

The function \fort{cg\_state\_write} creates the
\fort{ReferenceState\_t} node and must be called even if
\fort{StateDescription} is undefined (i.e., a blank string).
The descriptors, data arrays, data class, and dimensional units
characterizing the \fort{ReferenceState\_t} data structure may be added
to this data structure after its creation.

The function \fort{cg\_state\_read} reads the \fort{StateDescription}
of the local \fort{ReferenceState\_t} node.
If \fort{StateDescription} is undefined in the CGNS database, this
function returns a null string.
If \texttt{StateDescription} exists, the library will allocate the space
to store the description string, and return the description string to
the application.
It is the responsibility of the application to free this space when it
is no longer needed by a call to \texttt{cg\_free(StateDescription)},
described in \autoref{s:free}.

\subsection{Gravity}
\label{s:gravity}

\noindent
\textit{Node}: \fort{Gravity\_t}

\begin{fctbox}
\textcolor{output}{\textit{ier}} = cg\_gravity\_write(\textcolor{input}{int fn}, \textcolor{input}{int B}, \textcolor{input}{float *GravityVector}); & - w m \\
\textcolor{output}{\textit{ier}} = cg\_gravity\_read(\textcolor{input}{int fn}, \textcolor{input}{int B}, \textcolor{output}{\textit{float *GravityVector}}); & r - m \\
\hline
call cg\_gravity\_write\_f(\textcolor{input}{fn}, \textcolor{input}{B}, \textcolor{input}{GravityVector}, \textcolor{output}{\textit{ier}}) & - w m \\
call cg\_gravity\_read\_f(\textcolor{input}{fn}, \textcolor{input}{B}, \textcolor{output}{\textit{GravityVector}}, \textcolor{output}{\textit{ier}}) & r - m \\
\end{fctbox}

\noindent
\textbf{\textcolor{input}{Input}/\textcolor{output}{\textit{Output}}}

\begin{Ventryi}{\fort{GravityVector}}\raggedright
\item [\fort{fn}]
      CGNS file index number.
      (\textcolor{input}{\textit{Input}})
\item [\fort{B}]
      Base index number, where $1 \leq \text{\fort{B}} \leq \text{\fort{nbases}}$.
      (\textcolor{input}{\textit{Input}})
\item [\fort{GravityVector}]
      Description of the convergence information recorded in the data arrays.
      (\textcolor{input}{Input} for \fort{cg\_gravity\_write};
      \textcolor{output}{\textit{output}} for \fort{cg\_gravity\_read})
\item [\fort{ier}]
      Error status.
      (\textcolor{output}{\textit{Output}})
\end{Ventryi}

\subsection{Convergence History}
\label{s:convergence}

\noindent
\textit{Node}: \fort{ConvergenceHistory\_t}

\begin{fctbox}
\textcolor{output}{\textit{ier}} = cg\_convergence\_write(\textcolor{input}{int niterations}, \textcolor{input}{char *NormDefinitions}); & - w m \\
\textcolor{output}{\textit{ier}} = cg\_convergence\_read(\textcolor{output}{\textit{int *niterations}}, \textcolor{output}{\textit{char **NormDefinitions}}); & r - m \\
\hline
call cg\_convergence\_write\_f(\textcolor{input}{niterations}, \textcolor{input}{NormDefinitions}, \textcolor{output}{\textit{ier}}) & - w m \\
call cg\_convergence\_read\_f(\textcolor{output}{\textit{niterations}}, \textcolor{output}{\textit{NormDefinitions}}, \textcolor{output}{\textit{ier}}) & r - m \\
\end{fctbox}

\noindent
\textbf{\textcolor{input}{Input}/\textcolor{output}{\textit{Output}}}

\begin{Ventryi}{\fort{NormDefinitions}}\raggedright
\item [\fort{niterations}]
      Number of iterations for which convergence information is recorded.
      (\textcolor{input}{Input} for \fort{cg\_convergence\_write};
      \textcolor{output}{\textit{output}} for \fort{cg\_convergence\_read})
\item [\fort{NormDefinitions}]
      Description of the convergence information recorded in the data arrays.
      (\textcolor{input}{Input} for \fort{cg\_convergence\_write};
      \textcolor{output}{\textit{output}} for \fort{cg\_convergence\_read})
\item [\fort{ier}]
      Error status.
      (\textcolor{output}{\textit{Output}})
\end{Ventryi}

The function \fort{cg\_convergence\_write} creates a
\fort{ConvergenceHistory\_t} node.
It must be the first one called when recording convergence history data.
The \fort{NormDefinitions} may be left undefined (i.e., a blank string).
After creation of this node, the descriptors, data arrays, data class,
and dimensional units characterizing the \fort{ConvergenceHistory\_t}
data structure may be added.

The function \fort{cg\_convergence\_read} reads a
\fort{ConvergenceHistory\_t} node.
If \fort{NormDefinitions} is not defined in the CGNS database, this
function returns a null string.
If \texttt{NormDefinitions} exists, the library will allocate the space
to store the description string, and return the description string to
the application.
It is the responsibility of the application to free this space when it
is no longer needed by a call to \texttt{cg\_free(NormDefinitions)},
described in \autoref{s:free}.

\subsection{Integral Data}
\label{s:integral}

\noindent
\textit{Node}: \fort{IntegralData\_t}

\begin{fctbox}
\textcolor{output}{\textit{ier}} = cg\_integral\_write(\textcolor{input}{char *Name}); & - w m \\
\textcolor{output}{\textit{ier}} = cg\_nintegrals(\textcolor{output}{\textit{int *nintegrals}}); & r - m \\
\textcolor{output}{\textit{ier}} = cg\_integral\_read(\textcolor{input}{int Index}, \textcolor{output}{\textit{char *Name}}); & r - m \\
\hline
call cg\_integral\_write\_f(\textcolor{input}{Name}, \textcolor{output}{\textit{ier}}) & - w m \\
call cg\_nintegrals\_f(\textcolor{output}{\textit{nintegrals}}, \textcolor{output}{\textit{ier}}) & r - m \\
call cg\_integral\_read\_f(\textcolor{input}{Index}, \textcolor{output}{\textit{Name}}, \textcolor{output}{\textit{ier}}) & r - m \\
\end{fctbox}

\noindent
\textbf{\textcolor{input}{Input}/\textcolor{output}{\textit{Output}}}

\begin{Ventryi}{\fort{nintegrals}}\raggedright
\item [\fort{Name}]
      Name of the \fort{IntegralData\_t} data structure.
      (\textcolor{input}{Input} for \fort{cg\_integral\_write};
      \textcolor{output}{\textit{output}} for \fort{cg\_integral\_read})
\item [\fort{nintegrals}]
      Number of \fort{IntegralData\_t} nodes under current node.
      (\textcolor{output}{\textit{Output}})
\item [\fort{Index}]
      Integral data index number, where $1 \leq \text{\fort{Index}} \leq \text{\fort{nintegrals}}$.
      (\textcolor{input}{Input})
\item [\fort{ier}]
      Error status.
      (\textcolor{output}{\textit{Output}})
\end{Ventryi}

\subsection{User-Defined Data}
\label{s:userdefined}

\noindent
\textit{Node}: \fort{UserDefinedData\_t}

\begin{fctbox}
\textcolor{output}{\textit{ier}} = cg\_user\_data\_write(\textcolor{input}{char *Name}); & - w m \\
\textcolor{output}{\textit{ier}} = cg\_nuser\_data(\textcolor{output}{\textit{int *nuserdata}}); & r - m \\
\textcolor{output}{\textit{ier}} = cg\_user\_data\_read(\textcolor{input}{int Index}, \textcolor{output}{\textit{char *Name}}); & r - m \\
\hline
call cg\_user\_data\_write\_f(\textcolor{input}{Name}, \textcolor{output}{\textit{ier}}) & - w m \\
call cg\_nuser\_data\_f(\textcolor{output}{\textit{nuserdata}}, \textcolor{output}{\textit{ier}}) & r - m \\
call cg\_user\_data\_read\_f(\textcolor{input}{Index}, \textcolor{output}{\textit{Name}}, \textcolor{output}{\textit{ier}}) & r - m \\
\end{fctbox}

\noindent
\textbf{\textcolor{input}{Input}/\textcolor{output}{\textit{Output}}}

\begin{Ventryi}{\fort{nintegrals}}\raggedright
\item [\fort{nuserdata}]
      Number of \fort{UserDefinedData\_t} nodes under current node.
      (\textcolor{output}{\textit{Output}})
\item [\fort{Name}]
      Name of the \fort{UserDefinedData\_t} node.
      (\textcolor{input}{Input} for \fort{cg\_user\_data\_write};
      \textcolor{output}{\textit{output}} for \fort{cg\_user\_data\_read})
\item [\fort{Index}]
      User-defined data index number, where $1 \leq \text{\fort{Index}} \leq \text{\fort{nuserdata}}$.
      (\textcolor{input}{Input})
\item [\fort{ier}]
      Error status.
      (\textcolor{output}{\textit{Output}})
\end{Ventryi}

After accessing a particular \fort{UserDefinedData\_t} node using
\fort{cg\_goto}, the Point Set functions described in \autoref{s:ptset}
may be used to read or write point set information for the node.
The function \fort{cg\_gridlocation\_write} may also be used to specify
the location of the data with respect to the grid (e.g., \fort{Vertex}
or \fort{FaceCenter}).

Multiple levels of \fort{UserDefinedData\_t} nodes may be written and
retrieved by positioning via \fort{cg\_goto}.
E.g.,
\begin{alltt}
   ier = cg\_goto(fn, B, "Zone\_t", Z, "UserDefinedData\_t", ud1, 
                 "UserDefinedData\_t", ud2, "UserDefinedData\_t", ud3, "end");
\end{alltt}

\subsection{Freeing Memory}
\label{s:free}

\begin{fctbox}
\textcolor{output}{\textit{ier}} = cg\_free(\textcolor{input}{char *data}); & r w m \\
\end{fctbox}

\noindent
\textbf{\textcolor{input}{Input}/\textcolor{output}{\textit{Output}}}

\begin{Ventryi}{\fort{nintegrals}}\raggedright
\item [\fort{data}]
      Character data allocated by the Mid-Level Library.
      (\textcolor{input}{\textit{Input}})
\item [\fort{ier}]
      Error status.
      (\textcolor{output}{\textit{Output}})
\end{Ventryi}

This function does not affect the structure of a CGNS file; it is
provided as a convenience to free memory allocated by the Mid-Level
Library when using C.
This isn't necessary in Fortran, and thus an equivalent Fortran function
is not provided.

The functions that are used to allocate memory for return values are
\texttt{cg\_descriptor\_read}, \texttt{cg\_convergence\_read},
\texttt{cg\_geo\_read}, \texttt{cg\_link\_read}, and
\texttt{cg\_state\_read}.
Each of these may allocate space to contain the character data returned
to the application.
It is the responsibility of the application to free this data when it is
no longer needed.
Calling \texttt{cg\_free} is identical to calling the standard C function
\texttt{free}, however it is probably safer in that the memory is
freed in the same module in which it is created, particularly when the
Mid-Level Library is a shared library or DLL.
The routine checks for NULL data and will return \texttt{CG\_ERROR} in
this case, otherwise it returns \texttt{CG\_OK}.


\clearemptydoublepage
\section{Grid Specification}
\label{s:grid}
\thispagestyle{plain}

\subsection{Zone Grid Coordinates}
\label{s:gridcoordinates}

\noindent
\textit{Node}: \texttt{GridCoordinates\_t}

\texttt{GridCoordinates\_t} nodes are used to describe grids associated
with a particular zone.
The original grid must be described by a \texttt{GridCoordinates\_t} node
named \texttt{GridCoordinates}.
Additional \texttt{GridCoordinates\_t} nodes may be used, with user-defined
names, to store grids at multiple time steps or iterations.
In addition to the discussion of the \texttt{GridCoordinates\_t} node in
the \href{../sids/sids.pdf}{SIDS} and \href{../filemap/filemap.pdf}{File
Mapping} manuals, see the discussion of the \texttt{ZoneIterativeData\_t}
and \texttt{ArbitraryGridMotion\_t} nodes in the SIDS manual.

\begin{fctbox}
\textcolor{output}{\textit{ier}} = cg\_grid\_write(\textcolor{input}{int fn}, \textcolor{input}{int B}, \textcolor{input}{int Z}, \textcolor{input}{char *GridCoordName}, & - w m \\
~~~~~~\textcolor{output}{\textit{int *G}}); & \\
\textcolor{output}{\textit{ier}} = cg\_ngrids(\textcolor{input}{int fn}, \textcolor{input}{int B}, \textcolor{input}{int Z}, \textcolor{output}{\textit{int *ngrids}}); & - w m \\
\textcolor{output}{\textit{ier}} = cg\_grid\_read(\textcolor{input}{int fn}, \textcolor{input}{int B}, \textcolor{input}{int Z}, \textcolor{input}{int G}, \textcolor{output}{\textit{char *GridCoordName}}); & r - m \\
\hline
call cg\_grid\_write\_f(\textcolor{input}{fn}, \textcolor{input}{B}, \textcolor{input}{Z}, \textcolor{input}{GridCoordName}, \textcolor{output}{\textit{G}}, \textcolor{output}{\textit{ier}}) & - w m \\
call cg\_ngrids\_f(\textcolor{input}{fn}, \textcolor{input}{B}, \textcolor{input}{Z}, \textcolor{output}{\textit{ngrids}}, \textcolor{output}{\textit{ier}}) & - w m \\
call cg\_grid\_read\_f(\textcolor{input}{fn}, \textcolor{input}{B}, \textcolor{input}{Z}, \textcolor{input}{G}, \textcolor{output}{\textit{GridCoordName}}, \textcolor{output}{\textit{ier}}) & r - m \\
\end{fctbox}

\noindent
\textbf{\textcolor{input}{Input}/\textcolor{output}{Output}}

\noindent (Note that for Fortran calls, all integer arguments are integer*4 in 32-bit mode and integer*8 in 64-bit mode.
See ``64-bit Fortran Portability and Issues" section.)

\begin{Ventryi}{\texttt{GridCoordinateName}}\raggedright
\item [\texttt{fn}]
      CGNS file index number.
      (\textcolor{input}{Input})
\item [\texttt{B}]
      Base index number, where $1 \leq \text{\texttt{B}} \leq \text{\texttt{nbases}}$.
      (\textcolor{input}{Input})
\item [\texttt{Z}]
      Zone index number, where $1 \leq \text{\texttt{Z}} \leq \text{\texttt{nzones}}$.
      (\textcolor{input}{Input})
\item [\texttt{G}]
      Grid index number, where $1 \leq \text{\texttt{G}} \leq \text{\texttt{ngrids}}$.
      (\textcolor{input}{Input} for \texttt{cg\_grid\_read};
      \textcolor{output}{\textit{output}} for \texttt{cg\_grid\_write})
\item [\texttt{ngrids}]
      Number of \texttt{GridCoordinates\_t} nodes for zone \texttt{Z}.
      (\textcolor{output}{\textit{Output}})
\item [\texttt{GridCoordinateName}]
      Name of the \texttt{GridCoordinates\_t} node.
      Note that the name ``\texttt{GridCoordinates}'' is reserved for the
      original grid and must be the first \texttt{GridCoordinates\_t}
      node to be defined.
      (\textcolor{input}{Input} for \texttt{cg\_grid\_write};
      \textcolor{output}{\textit{output}} for \texttt{cg\_grid\_read})
\item [\texttt{ier}]
      Error status.
      (\textcolor{output}{\textit{Output}})
\end{Ventryi}

The above functions are applicable to any \texttt{GridCoordinates\_t} node.

\begin{fctbox}
\textcolor{output}{\textit{ier}} = cg\_coord\_write(\textcolor{input}{int fn}, \textcolor{input}{int B}, \textcolor{input}{int Z}, \textcolor{input}{DataType\_t datatype}, & - w m \\
~~~~~~\textcolor{input}{char *coordname}, \textcolor{input}{void *coord\_array}, \textcolor{output}{\textit{int *C}}); & \\
\textcolor{output}{\textit{ier}} = cg\_coord\_partial\_write(\textcolor{input}{int fn}, \textcolor{input}{int B}, \textcolor{input}{int Z}, & - w m \\
~~~~~~\textcolor{input}{DataType\_t datatype}, \textcolor{input}{char *coordname}, \textcolor{input}{cgsize\_t *range\_min}, & \\
~~~~~~\textcolor{input}{cgsize\_t *range\_max}, \textcolor{input}{void *coord\_array}, \textcolor{output}{\textit{int *C}}); & \\
\textcolor{output}{\textit{ier}} = cg\_ncoords(\textcolor{input}{int fn}, \textcolor{input}{int B}, \textcolor{input}{int Z}, \textcolor{output}{\textit{int *ncoords}}); & r - m \\
\textcolor{output}{\textit{ier}} = cg\_coord\_info(\textcolor{input}{int fn}, \textcolor{input}{int B}, \textcolor{input}{int Z}, \textcolor{input}{int C}, \textcolor{output}{\textit{DataType\_t *datatype}}, & r - m \\
~~~~~~\textcolor{output}{\textit{char *coordname}}); & \\
\textcolor{output}{\textit{ier}} = cg\_coord\_read(\textcolor{input}{int fn}, \textcolor{input}{int B}, \textcolor{input}{int Z}, \textcolor{input}{char *coordname}, & r - m \\
~~~~~~\textcolor{input}{DataType\_t datatype}, \textcolor{input}{cgsize\_t *range\_min}, \textcolor{input}{cgsize\_t *range\_max}, & \\
~~~~~~\textcolor{output}{\textit{void *coord\_array}}); & \\
\hline
call cg\_coord\_write\_f(\textcolor{input}{fn}, \textcolor{input}{B}, \textcolor{input}{Z}, \textcolor{input}{datatype}, \textcolor{input}{coordname}, \textcolor{input}{coord\_array}, \textcolor{output}{\textit{C}}, & - w m \\
~~~~~\textcolor{output}{\textit{ier}}) & \\
call cg\_coord\_partial\_write\_f(\textcolor{input}{fn}, \textcolor{input}{B}, \textcolor{input}{Z}, \textcolor{input}{datatype}, \textcolor{input}{coordname}, \textcolor{input}{range\_min}, & - w m \\
~~~~~\textcolor{input}{range\_max}, \textcolor{input}{coord\_array}, \textcolor{output}{\textit{C}}, \textcolor{output}{\textit{ier}}) & \\
call cg\_ncoords\_f(\textcolor{input}{fn}, \textcolor{input}{B}, \textcolor{input}{Z}, \textcolor{output}{\textit{ncoords}}, \textcolor{output}{\textit{ier}}) & r - m \\
call cg\_coord\_info\_f(\textcolor{input}{fn}, \textcolor{input}{B}, \textcolor{input}{Z}, \textcolor{input}{C}, \textcolor{output}{\textit{datatype}}, \textcolor{output}{\textit{coordname}}, \textcolor{output}{\textit{ier}}) & r - m \\
call cg\_coord\_read\_f(\textcolor{input}{fn}, \textcolor{input}{B}, \textcolor{input}{Z}, \textcolor{input}{coordname}, \textcolor{input}{datatype}, \textcolor{input}{range\_min}, & r - m \\
~~~~~\textcolor{input}{range\_max}, \textcolor{output}{\textit{coord\_array}}, \textcolor{output}{\textit{ier}}) & \\
\end{fctbox}

\noindent
\textbf{\textcolor{input}{Input}/\textcolor{output}{\textit{Output}}}

\noindent (Note that for Fortran calls, all integer arguments are integer*4 in 32-bit mode and integer*8 in 64-bit mode.
See ``64-bit Fortran Portability and Issues" section.)

\begin{Ventryi}{\texttt{coord\_array}}\raggedright
\item [\texttt{fn}]
      CGNS file index number.
      (\textcolor{input}{Input})
\item [\texttt{B}]
      Base index number, where $1 \leq \text{\texttt{B}} \leq \text{\texttt{nbases}}$.
      (\textcolor{input}{Input})
\item [\texttt{Z}]
      Zone index number, where $1 \leq \text{\texttt{Z}} \leq \text{\texttt{nzones}}$.
      (\textcolor{input}{Input})
\item [\texttt{C}]
      Coordinate array index number, where $1 \leq \text{\texttt{C}} \leq \text{\texttt{ncoords}}$.
      (\textcolor{input}{Input} for \texttt{cg\_coord\_info};
      \textcolor{output}{\textit{output}} for \texttt{cg\_coord\_write})
\item [\texttt{ncoords}]
      Number of coordinate arrays for zone \texttt{Z}.
      (\textcolor{output}{\textit{Output}})
\item [\texttt{datatype}]
      Data type in which the coordinate array is written.
      Admissible data types for a coordinate array are \texttt{RealSingle}
      and \texttt{RealDouble}.
      (\textcolor{input}{Input} for \texttt{cg\_coord\_write},
      \texttt{cg\_coord\_partial\_write}, \texttt{cg\_coord\_read};
      \textcolor{output}{\textit{output}} for \texttt{cg\_coord\_info})
\item [\texttt{coordname}]
      Name of the coordinate array.
      It is strongly advised to use the SIDS nomenclature conventions
      when naming the coordinate arrays to insure file compatibility.
      (\textcolor{input}{Input} for \texttt{cg\_coord\_write},
      \texttt{cg\_coord\_partial\_write}, \texttt{cg\_coord\_read};
      \textcolor{output}{\textit{output}} for \texttt{cg\_coord\_info})
\item [\texttt{range\_min}]
      Lower range index (eg., \texttt{imin, jmin, kmin}).
      (\textcolor{input}{Input})
\item [\texttt{range\_max}]
      Upper range index (eg., \texttt{imax, jmax, kmax}).
      (\textcolor{input}{Input})
\item [\texttt{coord\_array}]
      Array of coordinate values for the range prescribed.
      (\textcolor{input}{Input} for \texttt{cg\_coord\_write};
      \texttt{cg\_coord\_partial\_write}, \textcolor{output}{\textit{output}} for \texttt{cg\_coord\_read})
\item [\texttt{ier}]
      Error status.
      (\textcolor{output}{\textit{Output}})
\end{Ventryi}

The above functions are applicable \emph{only} to the 
\texttt{GridCoordinates\_t} node named \texttt{GridCoordinates}, used
for the original grid in a zone.
Coordinates for additional \texttt{GridCoordinates\_t} nodes in a zone
must be read and written using the \texttt{cg\_array\_\textit{xxx}} functions
described in \autoref{s:dataarray}.

When writing, the function \texttt{cg\_coord\_write} will
automatically write the full range of coordinates (i.e., the entire
\texttt{coord\_array}).
The function \texttt{cg\_coord\_partial\_write} may be used to write
only a subset of \texttt{coord\_array}.
When using the partial write, any existing data as defined by
\texttt{range\_min} and \texttt{range\_max} will be overwritten by the
new values.
All other values will not be affected.

The function \texttt{cg\_coord\_read} returns the coordinate array
\texttt{coord\_array}, for the range prescribed by \texttt{range\_min} and
\texttt{range\_max}.
The array is returned to the application in the data type requested in
\texttt{datatype}.
This data type does not need to be the same as the one in which the
coordinates are stored in the file.
A coordinate array stored as double precision in the CGNS file can be
returned to the application as single precision, or vice versa.

In Fortran, when using \texttt{cg\_coord\_read\_f} to read 2D or 3D
coordinates, the extent of each dimension of \texttt{coord\_array} must
be consistent with the requested range.
When reading a 1D solution, the declared size can be larger than the
requested range.
For example, for a 2D zone with $100 \times 50$ vertices, if
\texttt{range\_min} and \texttt{range\_max} are set to (11,11) and (20,20)
to read a subset of the coordinates, then \texttt{coord\_array} must be
dimensioned (10,10).
If \texttt{coord\_array} is declared larger (e.g., (100,50)) the
indices for the returned coordinates will be wrong.

\newpage
\subsection{Element Connectivity}
\label{s:elements}

\noindent
\textit{Node}: \texttt{Elements\_t}

\begin{fctbox}
\textcolor{output}{\textit{ier}} = cg\_section\_write(\textcolor{input}{int fn}, \textcolor{input}{int B}, \textcolor{input}{int Z}, \textcolor{input}{char *ElementSectionName}, & - w m \\
~~~~~~\textcolor{input}{ElementType\_t type}, \textcolor{input}{cgsize\_t start}, \textcolor{input}{cgsize\_t end}, \textcolor{input}{int nbndry}, & \\
~~~~~~\textcolor{input}{cgsize\_t *Elements}, \textcolor{output}{\textit{int *S}}); & \\
\textcolor{output}{\textit{ier}} = cg\_section\_partial\_write(\textcolor{input}{int fn}, \textcolor{input}{int B}, \textcolor{input}{int Z}, & - w m \\
~~~~~~\textcolor{input}{char *ElementSectionName}, \textcolor{input}{ElementType\_t type}, \textcolor{input}{cgsize\_t start}, & \\
~~~~~~\textcolor{input}{cgsize\_t end}, \textcolor{input}{int nbndry}, \textcolor{output}{\textit{int *S}}); & \\
\textcolor{output}{\textit{ier}} = cg\_elements\_partial\_write(\textcolor{input}{int fn}, \textcolor{input}{int B}, \textcolor{input}{int Z}, \textcolor{input}{int S}, & - w m \\
~~~~~~\textcolor{input}{cgsize\_t start}, \textcolor{input}{cgsize\_t end}, \textcolor{input}{cgsize\_t *Elements}); & \\
\textcolor{output}{\textit{ier}} = cg\_parent\_data\_write(\textcolor{input}{int fn}, \textcolor{input}{int B}, \textcolor{input}{int Z}, \textcolor{input}{int S}, & - w m \\
~~~~~~\textcolor{input}{cgsize\_t *ParentData}); & \\
\textcolor{output}{\textit{ier}} = cg\_parent\_data\_partial\_write(\textcolor{input}{int fn}, \textcolor{input}{int B}, \textcolor{input}{int Z}, \textcolor{input}{int S}, & - w m \\
~~~~~~\textcolor{input}{cgsize\_t start}, \textcolor{input}{cgsize\_t end}, \textcolor{input}{cgsize\_t *ParentData}); & \\
\textcolor{output}{\textit{ier}} = cg\_nsections(\textcolor{input}{int fn}, \textcolor{input}{int B}, \textcolor{input}{int Z}, \textcolor{output}{\textit{int *nsections}}); & r - m \\
\textcolor{output}{\textit{ier}} = cg\_section\_read(\textcolor{input}{int fn}, \textcolor{input}{int B}, \textcolor{input}{int Z}, \textcolor{input}{int S}, & r - m \\
~~~~~~\textcolor{output}{\textit{char *ElementSectionName}}, \textcolor{output}{\textit{ElementType\_t *type}}, \textcolor{output}{\textit{cgsize\_t *start}}, & \\
~~~~~~\textcolor{output}{\textit{cgsize\_t *end}}, \textcolor{output}{\textit{int *nbndry}}, \textcolor{output}{\textit{int *parent\_flag}}); & \\
\textcolor{output}{\textit{ier}} = cg\_ElementDataSize(\textcolor{input}{int fn}, \textcolor{input}{int B}, \textcolor{input}{int Z}, \textcolor{input}{int S}, & r - m \\
~~~~~~\textcolor{output}{\textit{cgsize\_t *ElementDataSize}}); & \\
\textcolor{output}{\textit{ier}} = cg\_ElementPartialSize(\textcolor{input}{int fn}, \textcolor{input}{int B}, \textcolor{input}{int Z}, \textcolor{input}{int S}, & r - m \\
~~~~~~\textcolor{input}{cgsize\_t start}, \textcolor{input}{cgsize\_t end}, \textcolor{output}{\textit{cgsize\_t *ElementDataSize}}); & \\
\textcolor{output}{\textit{ier}} = cg\_elements\_read(\textcolor{input}{int fn}, \textcolor{input}{int B}, \textcolor{input}{int Z}, \textcolor{input}{int S}, & r - m \\
~~~~~~\textcolor{output}{\textit{cgsize\_t *Elements}}, \textcolor{output}{\textit{cgsize\_t *ParentData}}); & \\
\textcolor{output}{\textit{ier}} = cg\_elements\_partial\_read(\textcolor{input}{int fn}, \textcolor{input}{int B}, \textcolor{input}{int Z}, \textcolor{input}{int S}, & r - m \\
~~~~~~\textcolor{input}{cgsize\_t start}, \textcolor{input}{cgsize\_t end}, \textcolor{output}{\textit{cgsize\_t *Elements}}, & \\
~~~~~~\textcolor{output}{\textit{cgsize\_t *ParentData}}); & \\
\textcolor{output}{\textit{ier}} = cg\_npe(\textcolor{input}{ElementType\_t type}, \textcolor{output}{\textit{int *npe}}); & r w m \\
\end{fctbox}

\begin{fctbox}
call cg\_section\_write\_f(\textcolor{input}{fn}, \textcolor{input}{B}, \textcolor{input}{Z}, \textcolor{input}{ElementSectionName}, \textcolor{input}{type}, \textcolor{input}{start}, \textcolor{input}{end}, & - w m \\
~~~~~\textcolor{input}{nbndry}, \textcolor{input}{Elements}, \textcolor{output}{\textit{S}}, \textcolor{output}{\textit{ier}}) & \\
call cg\_section\_partial\_write\_f(\textcolor{input}{fn}, \textcolor{input}{B}, \textcolor{input}{Z}, \textcolor{input}{ElementSectionName}, \textcolor{input}{type}, & - w m \\
~~~~~\textcolor{input}{start}, \textcolor{input}{end}, \textcolor{input}{nbndry}, \textcolor{output}{\textit{S}}, \textcolor{output}{\textit{ier}}) & \\
call cg\_elements\_partial\_write\_f(\textcolor{input}{fn}, \textcolor{input}{B}, \textcolor{input}{Z}, \textcolor{input}{S}, & - w m \\
~~~~~~\textcolor{input}{start}, \textcolor{input}{end}, \textcolor{input}{Elements}, \textcolor{output}{\textit{ier}}); & \\
call cg\_parent\_data\_write\_f(\textcolor{input}{fn}, \textcolor{input}{B}, \textcolor{input}{Z}, \textcolor{input}{S}, \textcolor{input}{ParentData}, \textcolor{output}{\textit{ier}}) & - w m \\
call cg\_parent\_data\_partial\_write\_f(\textcolor{input}{fn}, \textcolor{input}{B}, \textcolor{input}{Z}, \textcolor{input}{S}, \textcolor{input}{start}, & - w m \\
~~~~~\textcolor{input}{end}, \textcolor{input}{ParentData}, \textcolor{output}{\textit{ier}}) & \\
call cg\_nsections\_f(\textcolor{input}{fn}, \textcolor{input}{B}, \textcolor{input}{Z}, \textcolor{output}{\textit{nsections}}, \textcolor{output}{\textit{ier}}) & r - m \\
call cg\_section\_read\_f(\textcolor{input}{fn}, \textcolor{input}{B}, \textcolor{input}{Z}, \textcolor{input}{S}, \textcolor{output}{\textit{ElementSectionName}}, \textcolor{output}{\textit{type}}, & r - m \\
~~~~~\textcolor{output}{\textit{start}}, \textcolor{output}{\textit{end}}, \textcolor{output}{\textit{nbndry}}, \textcolor{output}{\textit{parent\_flag}}, \textcolor{output}{\textit{ier}}) & \\
call cg\_ElementDataSize\_f(\textcolor{input}{fn}, \textcolor{input}{B}, \textcolor{input}{Z}, \textcolor{input}{S}, \textcolor{output}{\textit{ElementDataSize}}, \textcolor{output}{\textit{ier}}) & r - m \\
call cg\_ElementPartialSize\_f(\textcolor{input}{fn}, \textcolor{input}{B}, \textcolor{input}{Z}, \textcolor{input}{S}, \textcolor{input}{start}, \textcolor{input}{end}, \textcolor{output}{\textit{ElementDataSize}}, & r - m \\
~~~~~\textcolor{output}{\textit{ier}}) & \\
call cg\_elements\_read\_f(\textcolor{input}{fn}, \textcolor{input}{B}, \textcolor{input}{Z}, \textcolor{input}{S}, \textcolor{output}{\textit{Elements}}, \textcolor{output}{\textit{ParentData}}, \textcolor{output}{\textit{ier}}) & r - m \\
call cg\_elements\_partial\_read\_f(\textcolor{input}{fn}, \textcolor{input}{B}, \textcolor{input}{Z}, \textcolor{input}{S}, \textcolor{input}{start}, \textcolor{input}{end}, \textcolor{output}{\textit{Elements}}, & r - m \\
~~~~~\textcolor{output}{\textit{ParentData}}, \textcolor{output}{\textit{ier}}) & \\
call cg\_npe\_f(\textcolor{input}{type}, \textcolor{output}{\textit{npe}}, \textcolor{output}{\textit{ier}}) & r w m \\
\end{fctbox}

\noindent
\textbf{\textcolor{input}{Input}/\textcolor{output}{\textit{Output}}}

\noindent (Note that for Fortran calls, all integer arguments are integer*4 in 32-bit mode and integer*8 in 64-bit mode.
See ``64-bit Fortran Portability and Issues" section.)

\begin{Ventryi}{\texttt{ElementSectionName}}\raggedright
\item [\texttt{fn}]
      CGNS file index number.
      (\textcolor{input}{Input})
\item [\texttt{B}]
      Base index number, where $1 \leq \text{\texttt{B}} \leq \text{\texttt{nbases}}$.
      (\textcolor{input}{Input})
\item [\texttt{Z}]
      Zone index number, where $1 \leq \text{\texttt{Z}} \leq \text{\texttt{nzones}}$.
      (\textcolor{input}{Input})
\item [\texttt{ElementSectionName}]
      Name of the \texttt{Elements\_t} node.
      (\textcolor{input}{Input} for \texttt{cg\_section\_write};
      \textcolor{output}{\textit{output}} for \texttt{cg\_section\_read})
\item [\texttt{type}]
      Type of element.
      See the eligible types for \texttt{ElementType\_t} in \autoref{s:typedefs}.
      (\textcolor{input}{Input} for \texttt{cg\_section\_write}, \texttt{cg\_npe};
      \textcolor{output}{\textit{output}} for \texttt{cg\_section\_read})
\item [\texttt{start}]
      Index of first element in the section.
      (\textcolor{input}{Input} for
      \texttt{cg\_section\_write},
      \texttt{cg\_section\_partial\_write},
      \texttt{cg\_parent\_data\_partial\_write},
      \texttt{cg\_ElementPartialSize},
      \texttt{cg\_elements\_partial\_write};
      \textcolor{output}{\textit{output}} for
      \texttt{cg\_section\_read})
\item [\texttt{end}]
      Index of last element in the section.
      (\textcolor{input}{Input} for
      \texttt{cg\_section\_write},
      \texttt{cg\_section\_partial\_write},
      \texttt{cg\_parent\_data\_partial\_write},
      \texttt{cg\_ElementPartialSize},
      \texttt{cg\_elements\_partial\_write};
      \textcolor{output}{\textit{output}} for
      \texttt{cg\_section\_read})
\item [\texttt{nbndry}]
      Index of last boundary element in the section.
      Set to zero if the elements are unsorted.
      (\textcolor{input}{Input} for \texttt{cg\_section\_write};
      \textcolor{output}{\textit{output}} for \texttt{cg\_section\_read})
\item [\texttt{nsections}]
      Lower range index (eg., \texttt{imin, jmin, kmin}).
      (\textcolor{output}{\textit{Output}})
\item [\texttt{S}]
      Element section index, where $1 \leq \text{\texttt{S}} \leq \text{\texttt{nsections}}$.
      (\textcolor{input}{Input} for \texttt{cg\_parent\_data\_write},
      \texttt{cg\_section\_read}, \texttt{cg\_ElementDataSize},
      \texttt{cg\_elements\_read};
      \textcolor{output}{\textit{output}} for \texttt{cg\_section\_write})
\item [\texttt{parent\_flag}]
      Flag indicating if the parent data are defined.
      If the parent data exist, \texttt{parent\_flag} is set to 1;
      otherwise it is set to 0.
      (\textcolor{output}{\textit{Output}})
\item [\texttt{ElementDataSize}]
      Number of element connectivity data values.
      (\textcolor{output}{\textit{Output}})
\item [\texttt{Elements}]
      Element connectivity data.
      (\textcolor{input}{Input} for \texttt{cg\_section\_write};
      \textcolor{output}{\textit{output}} for \texttt{cg\_elements\_read})
\item [\texttt{ParentData}]
      For boundary or interface elements, this array contains
      information on the cell(s) and cell face(s) sharing the element.
      If you do not need to read the \texttt{ParentData} when reading
      the \texttt{ElementData}, you may set the value to \texttt{NULL}.
      (\textcolor{output}{\textit{Output}})
\item [\texttt{npe}]
      Number of nodes for an element of type \texttt{type}.
      (\textcolor{output}{\textit{Output}})
\item [\texttt{ier}]
      Error status.
      (\textcolor{output}{\textit{Output}})
\end{Ventryi}

It is important to note that each element under a given \texttt{Zone\_t} -- including
all cells, faces, edges, boundary elements, etc. -- must have a unique element index
number. The numbering should be consecutive (i.e., no gaps).
This global numbering system insures that each and every element
within a zone is uniquely identified by its number.

If the specified \texttt{Elements\_t} node doesn't yet exist,
it may be created using either \texttt{cg\_sec\-tion\_write} or
\texttt{cg\_section\_partial\_write}.
The function \texttt{cg\_section\_write} writes the full range as
indicated by \texttt{start} and \texttt{end}
and supplied by the element connectivity array \texttt{Elements}. The 
\texttt{cg\_section\_partial\_write} function will create the element 
section data for the range start to end with the element data intialized 
to $0$. To add elements to the section, use \texttt{cg\_elements\_partial\_write} 
and parent data (it it exists) using \texttt{cg\_parent\_data\_partial\_write}. 
Both of these functions will replace the data for the range as indicated by 
\texttt{start} and \texttt{end} with the new values. In most cases, the data is 
not duplicated in the mid-level library, but written directly from the user data 
to disk. The exception to this is in the case of \texttt{MIXED}, 
\texttt{NGON\_n}, and \texttt{NFACE\_n} element sets. Since the size of the 
element connectivity array is not known directly, the MLL will keep a copy 
of the data in memory for the partial writes.

The function \texttt{cg\_elements\_read} returns all of the element
connectivity and parent data.
Specified subsets of the element connectivity and parent data may be
read using \texttt{cg\_elements\_partial\_read}.

\subsection{Axisymmetry}
\label{s:axisymmetry}

\noindent
\textit{Node}: \texttt{Axisymmetry\_t}

\begin{fctbox}
\textcolor{output}{\textit{ier}} = cg\_axisym\_write(\textcolor{input}{int fn}, \textcolor{input}{int B}, \textcolor{input}{float *ReferencePoint}, & - w m \\
~~~~~~\textcolor{input}{float *AxisVector}); & \\
\textcolor{output}{\textit{ier}} = cg\_axisym\_read(\textcolor{input}{int fn}, \textcolor{input}{int B}, \textcolor{output}{\textit{float *ReferencePoint}}, & r - m \\
~~~~~~\textcolor{output}{\textit{float *AxisVector}}); & \\
\hline
call cg\_axisym\_write\_f(\textcolor{input}{fn}, \textcolor{input}{B}, \textcolor{input}{ReferencePoint}, \textcolor{input}{AxisVector}, \textcolor{output}{\textit{ier}}) & - w m \\
call cg\_axisym\_read\_f(\textcolor{input}{fn}, \textcolor{input}{B}, \textcolor{output}{\textit{ReferencePoint}}, \textcolor{output}{\textit{AxisVector}}, \textcolor{output}{\textit{ier}}) & r - m \\
\end{fctbox}

\noindent
\textbf{\textcolor{input}{Input}/\textcolor{output}{\textit{Output}}}

\noindent (Note that for Fortran calls, all integer arguments are integer*4 in 32-bit mode and integer*8 in 64-bit mode.
See ``64-bit Fortran Portability and Issues" section.)

\begin{Ventryi}{\texttt{ReferencePoint}}\raggedright
\item [\texttt{fn}]
      CGNS file index number.
      (\textcolor{input}{\textit{Input}})
\item [\texttt{B}]
      Base index number, where $1 \leq \text{\texttt{B}} \leq \text{\texttt{nbases}}$.
      (\textcolor{input}{\textit{Input}})
\item [\texttt{ReferencePoint}]
      Origin used for defining the axis of rotation.
      (In Fortran, this is an array of Real*4 values.)
      (\textcolor{input}{Input} for \texttt{cg\_axisym\_write};
      \textcolor{output}{\textit{output}} for \texttt{cg\_axisym\_read})
\item [\texttt{AxisVector}]
      Direction cosines of the axis of rotation, through the reference
      point.
      (In Fortran, this is an array of Real*4 values.)
      (\textcolor{input}{Input} for \texttt{cg\_axisym\_write};
      \textcolor{output}{\textit{output}} for \texttt{cg\_axisym\_read})
\item [\texttt{ier}]
      Error status.
      (\textcolor{output}{\textit{Output}})
\end{Ventryi}

This node can only be used for a bi-dimensional model, i.e.,
\texttt{PhysicalDimension} must equal two.

\subsection{Rotating Coordinates}
\label{s:rotatingcoordinates}

\noindent
\textit{Node}: \texttt{RotatingCoordinates\_t}

\begin{fctbox}
\textcolor{output}{\textit{ier}} = cg\_rotating\_write(\textcolor{input}{float *RotationRateVector}, & - w m \\
~~~~~~\textcolor{input}{float *RotationCenter}); & \\
\textcolor{output}{\textit{ier}} = cg\_rotating\_read(\textcolor{output}{\textit{float *RotationRateVector}}, & r - m \\
~~~~~~\textcolor{output}{\textit{float *RotationCenter}}); & \\
\hline
call cg\_rotating\_write\_f(\textcolor{input}{RotationRateVector}, \textcolor{input}{RotationCenter}, \textcolor{output}{\textit{ier}}) & - w m \\
call cg\_rotating\_read\_f(\textcolor{output}{\textit{RotationRateVector}}, \textcolor{output}{\textit{RotationCenter}}, \textcolor{output}{\textit{ier}}) & r - m \\
\end{fctbox}

\noindent
\textbf{\textcolor{input}{Input}/\textcolor{output}{\textit{Output}}}

\noindent (Note that for Fortran calls, all integer arguments are integer*4 in 32-bit mode and integer*8 in 64-bit mode.
See ``64-bit Fortran Portability and Issues" section.)

\begin{Ventryi}{\texttt{RotationRateVector}}\raggedright
\item [\texttt{RotationRateVector}]
      Components of the angular velocity of the grid about the center
      of rotation.
      (In Fortran, this is an array of Real*4 values.)
      (\textcolor{input}{Input} for \texttt{cg\_rotating\_write};
      \textcolor{output}{\textit{output}} for \texttt{cg\_rotating\_read})
\item [\texttt{RotationCenter}]
      Coordinates of the center of rotation.
      (In Fortran, this is an array of Real*4 values.)
      (\textcolor{input}{Input} for \texttt{cg\_rotating\_write};
      \textcolor{output}{\textit{output}} for \texttt{cg\_rotating\_read})
\item [\texttt{ier}]
      Error status.
      (\textcolor{output}{\textit{Output}})
\end{Ventryi}


\clearemptydoublepage
\section{Solution Data}
\label{s:solution}
\thispagestyle{plain}

\subsection{Flow Solution}
\label{s:flowsolution}

\noindent
\textit{Node}: \texttt{FlowSolution\_t}

\begin{fctbox}
\textcolor{output}{\textit{ier}} = cg\_sol\_write(\textcolor{input}{int fn}, \textcolor{input}{int B}, \textcolor{input}{int Z}, \textcolor{input}{char *solname}, & - w m \\
~~~~~~\textcolor{input}{GridLocation\_t location}, \textcolor{output}{\textit{int *S}}); & \\
\textcolor{output}{\textit{ier}} = cg\_nsols(\textcolor{input}{int fn}, \textcolor{input}{int B}, \textcolor{input}{int Z}, \textcolor{output}{\textit{int *nsols}}); & r - m \\
\textcolor{output}{\textit{ier}} = cg\_sol\_info(\textcolor{input}{int fn}, \textcolor{input}{int B}, \textcolor{input}{int Z}, \textcolor{input}{int S}, \textcolor{output}{\textit{char *solname}}, & r - m \\
~~~~~~\textcolor{output}{\textit{GridLocation\_t *location}}); & \\
\textcolor{output}{\textit{ier}} = cg\_sol\_ptset\_write(\textcolor{input}{int fn}, \textcolor{input}{int B}, \textcolor{input}{int Z}, \textcolor{input}{char *solname}, & - w m \\
~~~~~~\textcolor{input}{GridLocation\_t location}, \textcolor{input}{PointSetType\_t ptset\_type}, & \\
~~~~~~\textcolor{input}{cgsize\_t npnts}, \textcolor{input}{cgsize\_t *pnts}, \textcolor{output}{\textit{int *S}}); & \\
\textcolor{output}{\textit{ier}} = cg\_sol\_ptset\_info(\textcolor{input}{int fn}, \textcolor{input}{int B}, \textcolor{input}{int Z}, \textcolor{input}{int S}, & r - m \\
~~~~~~\textcolor{output}{\textit{PointSetType\_t *ptset\_type}}, \textcolor{output}{\textit{cgsize\_t *npnts}}); & \\
\textcolor{output}{\textit{ier}} = cg\_sol\_ptset\_read(\textcolor{input}{int fn}, \textcolor{input}{int B}, \textcolor{input}{int Z}, \textcolor{input}{int S}, \textcolor{output}{\textit{cgsize\_t *pnts}}); & r - m \\
\textcolor{output}{\textit{ier}} = cg\_sol\_size(\textcolor{input}{int fn}, \textcolor{input}{int B}, \textcolor{input}{int Z}, \textcolor{input}{int S}, \textcolor{output}{\textit{int *data\_dim}}, & r - m \\
~~~~~~\textcolor{output}{\textit{cgsize\_t *dim\_vals}}); & \\
\hline
call cg\_sol\_write\_f(\textcolor{input}{fn}, \textcolor{input}{B}, \textcolor{input}{Z}, \textcolor{input}{solname}, \textcolor{input}{location}, \textcolor{output}{\textit{S}}, \textcolor{output}{\textit{ier}}) & - w m \\
call cg\_nsols\_f(\textcolor{input}{fn}, \textcolor{input}{B}, \textcolor{input}{Z}, \textcolor{output}{\textit{nsols}}, \textcolor{output}{\textit{ier}}) & r - m \\
call cg\_sol\_info\_f(\textcolor{input}{fn}, \textcolor{input}{B}, \textcolor{input}{Z}, \textcolor{input}{S}, \textcolor{output}{\textit{solname}}, \textcolor{output}{\textit{location}}, \textcolor{output}{\textit{ier}}) & r - m \\
call cg\_sol\_ptset\_write\_f(\textcolor{input}{fn}, \textcolor{input}{B}, \textcolor{input}{Z}, \textcolor{input}{solname}, \textcolor{input}{location}, \textcolor{input}{ptset\_type}, & - w m \\
~~~~~\textcolor{input}{npnts}, \textcolor{input}{pnts}, \textcolor{output}{\textit{S}}, \textcolor{output}{\textit{ier}}) & \\
call cg\_sol\_ptset\_info\_f(\textcolor{input}{fn}, \textcolor{input}{B}, \textcolor{input}{Z}, \textcolor{input}{S}, \textcolor{output}{\textit{ptset\_type}}, \textcolor{output}{\textit{npnts}}, \textcolor{output}{\textit{ier}}) & r - m \\
call cg\_sol\_ptset\_read\_f(\textcolor{input}{fn}, \textcolor{input}{B}, \textcolor{input}{Z}, \textcolor{input}{S}, \textcolor{output}{\textit{pnts}}, \textcolor{output}{\textit{ier}}) & r - m \\
call cg\_sol\_size\_f(\textcolor{input}{fn}, \textcolor{input}{B}, \textcolor{input}{Z}, \textcolor{input}{S}, \textcolor{output}{\textit{data\_dim}}, \textcolor{output}{\textit{dim\_vals}}, \textcolor{output}{\textit{ier}}) & r - m \\
\end{fctbox}

\noindent
\textbf{\textcolor{input}{Input}/\textcolor{output}{\textit{Output}}}

\noindent (Note that for Fortran calls, all integer arguments are integer*4 in 32-bit mode and integer*8 in 64-bit mode.
See ``64-bit Fortran Portability and Issues" section.)

\begin{Ventryi}{\texttt{ptset\_type}}\raggedright
\item [\texttt{fn}]
      CGNS file index number.
      (\textcolor{input}{Input})
\item [\texttt{B}]
      Base index number, where $1 \leq \text{\texttt{B}} \leq \text{\texttt{nbases}}$.
      (\textcolor{input}{Input})
\item [\texttt{Z}]
      Zone index number, where $1 \leq \text{\texttt{Z}} \leq \text{\texttt{nzones}}$.
      (\textcolor{input}{Input})
\item [\texttt{S}]
      Flow solution index number, where $1 \leq \text{\texttt{S}} \leq \text{\texttt{nsols}}$.
      (\textcolor{input}{Input} for \texttt{cg\_sol\_info};
      \textcolor{output}{\textit{output}} for \texttt{cg\_sol\_write})
\item [\texttt{nsols}]
      Number of flow solutions for zone \texttt{Z}.
      (\textcolor{output}{\textit{Output}})
\item [\texttt{solname}]
      Name of the flow solution.
      (\textcolor{input}{Input} for \texttt{cg\_sol\_write};
      \textcolor{output}{\textit{output}} for \texttt{cg\_sol\_info})
\item [\texttt{location}]
      Grid location where the solution is recorded.
      The current admissible locations are \texttt{Vertex},
      \texttt{CellCenter}, \texttt{IFaceCenter}, \texttt{JFaceCenter},
      and \texttt{KFaceCenter}.
      (\textcolor{input}{Input} for \texttt{cg\_sol\_write} and
      \texttt{cg\_sol\_ptset\_write};
      \textcolor{output}{\textit{output}} for \texttt{cg\_sol\_info})
\item [\texttt{ptset\_type}]
      Type of point set defining the interface in the current solution;
      either \texttt{PointRange} or \texttt{PointList}.
      (\textcolor{input}{Input} for \texttt{cg\_sol\_ptset\_write};
      \textcolor{output}{\textit{output}} for
      \texttt{cg\_sol\_ptset\_info})
\item [\texttt{npnts}]
      Number of points defining the interface in the current solution.
      For a \texttt{ptset\_type} of \texttt{PointRange}, npnts
      is always two. For a \texttt{ptset\_type} of \texttt{PointList},
      npnts is the number of points in the list. 
      (\textcolor{input}{Input} for \texttt{cg\_sol\_ptset\_write};
      \textcolor{output}{\textit{output}} for
      \texttt{cg\_sol\_ptset\_info})
\item [\texttt{pnts}]
      Array of points defining the interface in the current solution.
      (\textcolor{input}{Input} for \texttt{cg\_sol\_ptset\_write};
      \textcolor{output}{\textit{output}} for
      \texttt{cg\_sol\_ptset\_read})
\item [\texttt{data\_dim}]
      Number of dimensions defining the solution data.
      If a point set has been defined, this will be 1,
      otherwise this will be the current zone index dimension.
      (\textcolor{output}{\textit{Output}})
\item [\texttt{dim\_vals}]
      The array of \texttt{data\_dim} dimensions for the solution data.
      (\textcolor{output}{\textit{Output}})
\item [\texttt{ier}]
      Error status.
      (\textcolor{output}{\textit{Output}})
\end{Ventryi}

The above functions are used to create, and get information about,
\texttt{FlowSolution\_t} nodes.

Solution data may be specified over the entire zone, as in previous
versions of the library, using \texttt{cg\_sol\_write} or over
a patch region of the zone using \texttt{cg\_sol\_ptset\_write}.
The two are mutually exclusive. In the first case, the size of
the solution data is determined by the size of the zone, the
grid location, and rind data (if any) as in previous versions.
In the second case the solution data size is entirely determined
by the \texttt{PointList/PointRange} patch. In order to determine which of
the two forms in which the solution data was written, use
\texttt{cg\_sol\_ptset\_info}. If the solution is over the entire
zone, \texttt{ptset\_type} will be returned as \texttt{CG\_Null}
and \texttt{npnts} as 0. Otherwise, \texttt{ptset\_type} will
be either \texttt{PointList} or \texttt{PointRange}, and the
number of points will be returned in \texttt{npnts}.

To assist in determining the size of the solution data,
the function \texttt{cg\_sol\_size} has been added. For a solution
defined over the full zone, \texttt{data\_dim} returns the index
dimension for the zone, and \texttt{dim\_vals} specifies the
dimensions of the data, corrected for the grid location type
and rind data. If a point set patch has been specified,
\texttt{data\_dim} will be 1 and \texttt{dim\_vals} will contain
the total size of the patch.

Acceptable values of \texttt{GridLocation\_t} are \texttt{Vertex}
and \texttt{CellCenter}. If the base cell dimension is 2 or
greater (surface or volume), then \texttt{EdgeCenter} is
also allowed. For 3 dimensional bases, \texttt{FaceCenter}, and
for structured zones, \texttt{IFaceCenter}, \texttt{JFaceCenter}
and \texttt{KFaceCenter}, are also acceptable.

\newpage
\subsection{Flow Solution Data}
\label{s:flowsolutionarray}

\begin{fctbox}
\textcolor{output}{\textit{ier}} = cg\_field\_write(\textcolor{input}{int fn}, \textcolor{input}{int B}, \textcolor{input}{int Z}, \textcolor{input}{int S}, \textcolor{input}{DataType\_t datatype}, & - w m \\
~~~~~~\textcolor{input}{char *fieldname}, \textcolor{input}{void *solution\_array}, \textcolor{output}{\textit{int *F}}); & \\
\textcolor{output}{\textit{ier}} = cg\_field\_partial\_write(\textcolor{input}{int fn}, \textcolor{input}{int B}, \textcolor{input}{int Z}, \textcolor{input}{int S}, & - w m \\
~~~~~~\textcolor{input}{DataType\_t datatype}, \textcolor{input}{char *fieldname}, \textcolor{input}{cgsize\_t *range\_min}, & \\
~~~~~~\textcolor{input}{cgsize\_t *range\_max}, \textcolor{input}{void *solution\_array}, \textcolor{output}{\textit{int *F}}); & \\
\textcolor{output}{\textit{ier}} = cg\_nfields(\textcolor{input}{int fn}, \textcolor{input}{int B}, \textcolor{input}{int Z}, \textcolor{input}{int S}, \textcolor{output}{\textit{int *nfields}}); & r - m \\
\textcolor{output}{\textit{ier}} = cg\_field\_info(\textcolor{input}{int fn}, \textcolor{input}{int B}, \textcolor{input}{int Z}, \textcolor{input}{int S}, \textcolor{input}{int F}, & r - m \\
~~~~~~\textcolor{output}{\textit{DataType\_t *datatype}}, \textcolor{output}{\textit{char *fieldname}}); & \\
\textcolor{output}{\textit{ier}} = cg\_field\_read(\textcolor{input}{int fn}, \textcolor{input}{int B}, \textcolor{input}{int Z}, \textcolor{input}{int S}, \textcolor{input}{char *fieldname}, & r - m \\
~~~~~~\textcolor{input}{DataType\_t datatype}, \textcolor{input}{cgsize\_t *range\_min}, \textcolor{input}{cgsize\_t *range\_max}, & \\
~~~~~~\textcolor{output}{\textit{void *solution\_array}}); & \\
\hline
call cg\_field\_write\_f(\textcolor{input}{fn}, \textcolor{input}{B}, \textcolor{input}{Z}, \textcolor{input}{S}, \textcolor{input}{datatype}, \textcolor{input}{fieldname}, \textcolor{input}{solution\_array}, & - w m \\
~~~~~\textcolor{output}{\textit{F}}, \textcolor{output}{\textit{ier}}) & \\
call cg\_field\_partial\_write\_f(\textcolor{input}{fn}, \textcolor{input}{B}, \textcolor{input}{Z}, \textcolor{input}{S}, \textcolor{input}{datatype}, \textcolor{input}{fieldname}, & - w m \\
~~~~~\textcolor{input}{range\_min}, \textcolor{input}{range\_max}, \textcolor{input}{solution\_array}, \textcolor{output}{\textit{F}}, \textcolor{output}{\textit{ier}}) & \\
call cg\_nfields\_f(\textcolor{input}{fn}, \textcolor{input}{B}, \textcolor{input}{Z}, \textcolor{input}{S}, \textcolor{output}{\textit{nfields}}, \textcolor{output}{\textit{ier}}) & r - m \\
call cg\_field\_info\_f(\textcolor{input}{fn}, \textcolor{input}{B}, \textcolor{input}{Z}, \textcolor{input}{S}, \textcolor{input}{F}, \textcolor{output}{\textit{datatype}}, \textcolor{output}{\textit{fieldname}}, \textcolor{output}{\textit{ier}}) & r - m \\
call cg\_field\_read\_f(\textcolor{input}{fn}, \textcolor{input}{B}, \textcolor{input}{Z}, \textcolor{input}{S}, \textcolor{input}{fieldname}, \textcolor{input}{datatype}, \textcolor{input}{range\_min}, & r - m \\
~~~~~\textcolor{input}{range\_max}, \textcolor{output}{\textit{solution\_array}}, \textcolor{output}{\textit{ier}}) & \\
\end{fctbox}

\noindent
\textbf{\textcolor{input}{Input}/\textcolor{output}{\textit{Output}}}

\noindent (Note that for Fortran calls, all integer arguments are integer*4 in 32-bit mode and integer*8 in 64-bit mode.
See ``64-bit Fortran Portability and Issues" section.)

\begin{Ventryi}{\texttt{solution\_array}}\raggedright
\item [\texttt{fn}]
      CGNS file index number.
      (\textcolor{input}{Input})
\item [\texttt{B}]
      Base index number, where $1 \leq \text{\texttt{B}} \leq \text{\texttt{nbases}}$.
      (\textcolor{input}{Input})
\item [\texttt{Z}]
      Zone index number, where $1 \leq \text{\texttt{Z}} \leq \text{\texttt{nzones}}$.
      (\textcolor{input}{Input})
\item [\texttt{S}]
      Flow solution index number, where $1 \leq \text{\texttt{S}} \leq \text{\texttt{nsols}}$.
      (\textcolor{input}{Input})
\item [\texttt{F}]
      Solution array index number, where $1 \leq \text{\texttt{F}} \leq \text{\texttt{nfields}}$.
      (\textcolor{input}{Input} for \texttt{cg\_field\_info};
      \textcolor{output}{\textit{output}} for \texttt{cg\_field\_write})
\item [\texttt{nfields}]
      Number of data arrays in flow solution \texttt{S}.
      (\textcolor{output}{\textit{Output}})
\item [\texttt{datatype}]
      Data type in which the solution array is written.
      Admissible data types for a solution array are \texttt{Integer},
      \texttt{LongInteger}, \texttt{RealSingle}, and \texttt{RealDouble}.
      (\textcolor{input}{Input} for \texttt{cg\_field\_write},
      \texttt{cg\_field\_read};
      \textcolor{output}{\textit{output}} for \texttt{cg\_field\_info})
\item [\texttt{fieldname}]
      Name of the solution array.
      It is strongly advised to use the SIDS nomenclature conventions
      when naming the solution arrays to insure file compatibility.
      (\textcolor{input}{Input} for \texttt{cg\_field\_write},
      \texttt{cg\_field\_read};
      \textcolor{output}{\textit{output}} for \texttt{cg\_field\_info})
\item [\texttt{range\_min}]
      Lower range index (eg., \texttt{imin, jmin, kmin}).
      (\textcolor{input}{Input})
\item [\texttt{range\_max}]
      Upper range index (eg., \texttt{imax, jmax, kmax}).
      (\textcolor{input}{Input})
\item [\texttt{solution\_array}]
      Array of solution values for the range prescribed.
      (\textcolor{input}{Input} for \texttt{cg\_field\_write};
      \textcolor{output}{\textit{output}} for \texttt{cg\_field\_read})
\item [\texttt{ier}]
      Error status.
      (\textcolor{output}{\textit{Output}})
\end{Ventryi}

The above functions are used to read and write solution arrays stored
below a \texttt{FlowSolution\_t} node.

When writing, the function \texttt{cg\_field\_write} will
automatically write the full range of the solution (i.e., the entire
\texttt{solution\_array}).
The function \texttt{cg\_field\_partial\_write} may be used to write
only a subset of \texttt{solution\_array}.
When using the partial write, any existing data from \texttt{range\_min}
to \texttt{range\_max} will be overwritten by the new values.
All other values will not be affected.

The function \texttt{cg\_field\_read} returns the solution array
\texttt{fieldname}, for the range prescribed by \texttt{range\_min} and
\texttt{range\_max}.
The array is returned to the application in the data type requested in
\texttt{datatype}.
This data type does not need to be the same as the one in which the data is
stored in the file.
A solution array stored as double precision in the CGNS file can be
returned to the application as single precision, or vice versa.

In Fortran, when using \texttt{cg\_field\_read\_f} to read a 2D or 3D
solution, the extent of each dimension of \texttt{solution\_array} must
be consistent with the requested range.
When reading a 1D solution, the declared size can be larger than the
requested range.
For example, for a 2D zone with $100 \times 50$ vertices, if
\texttt{range\_min} and \texttt{range\_max} are set to (11,11) and (20,20)
to read a subset of the solution, then \texttt{solution\_array} must be
dimensioned (10,10).
If \texttt{solution\_array} is declared larger (e.g., (100,50)) the
indices for the returned array values will be wrong.

\newpage
\subsection{Discrete Data}
\label{s:discretedata}

\noindent
\textit{Node}: \texttt{DiscreteData\_t}

\begin{fctbox}
\textcolor{output}{\textit{ier}} = cg\_discrete\_write(\textcolor{input}{int fn}, \textcolor{input}{int B}, \textcolor{input}{int Z}, \textcolor{input}{char *DiscreteName}, & - w m \\
~~~~~~\textcolor{output}{\textit{int *D}}); & \\
\textcolor{output}{\textit{ier}} = cg\_ndiscrete(\textcolor{input}{int fn}, \textcolor{input}{int B}, \textcolor{input}{int Z}, \textcolor{output}{\textit{int *ndiscrete}}); & r - m \\
\textcolor{output}{\textit{ier}} = cg\_discrete\_read(\textcolor{input}{int fn}, \textcolor{input}{int B}, \textcolor{input}{int Z}, \textcolor{input}{int D}, & r - m \\
~~~~~~\textcolor{output}{\textit{char *DiscreteName}}); & \\
\textcolor{output}{\textit{ier}} = cg\_discrete\_ptset\_write(\textcolor{input}{int fn}, \textcolor{input}{int B}, \textcolor{input}{int Z}, \textcolor{input}{char *DiscreteName}, & - w m \\
~~~~~~\textcolor{input}{GridLocation\_t location}, \textcolor{input}{PointSetType\_t ptset\_type}, & \\
~~~~~~\textcolor{input}{cgsize\_t npnts}, \textcolor{input}{cgsize\_t *pnts}, \textcolor{output}{\textit{int *D}}); & \\
\textcolor{output}{\textit{ier}} = cg\_discrete\_ptset\_info(\textcolor{input}{int fn}, \textcolor{input}{int B}, \textcolor{input}{int Z}, \textcolor{input}{int D}, & r - m \\
~~~~~~\textcolor{output}{\textit{PointSetType\_t *ptset\_type}}, \textcolor{output}{\textit{cgsize\_t *npnts}}); & \\
\textcolor{output}{\textit{ier}} = cg\_discrete\_ptset\_read(\textcolor{input}{int fn}, \textcolor{input}{int B}, \textcolor{input}{int Z}, \textcolor{input}{int D}, & r - m \\
~~~~~~\textcolor{output}{\textit{cgsize\_t *pnts}}); & \\
\textcolor{output}{\textit{ier}} = cg\_discrete\_size(\textcolor{input}{int fn}, \textcolor{input}{int B}, \textcolor{input}{int Z}, \textcolor{input}{int D}, \textcolor{output}{\textit{int *data\_dim}}, & r - m \\
~~~~~~\textcolor{output}{\textit{cgsize\_t *dim\_vals}}); & \\
\hline
call cg\_discrete\_write\_f(\textcolor{input}{fn}, \textcolor{input}{B}, \textcolor{input}{Z}, \textcolor{input}{DiscreteName}, \textcolor{output}{\textit{D}}, \textcolor{output}{\textit{ier}}) & - w m \\
call cg\_ndiscrete\_f(\textcolor{input}{fn}, \textcolor{input}{B}, \textcolor{input}{Z}, \textcolor{output}{\textit{ndiscrete}}, \textcolor{output}{\textit{ier}}) & r - m \\
call cg\_discrete\_read\_f(\textcolor{input}{fn}, \textcolor{input}{B}, \textcolor{input}{Z}, \textcolor{input}{D}, \textcolor{output}{\textit{DiscreteName}}, \textcolor{output}{\textit{ier}}) & r - m \\
call cg\_discrete\_ptset\_write\_f(\textcolor{input}{fn}, \textcolor{input}{B}, \textcolor{input}{Z}, \textcolor{input}{DiscreteName}, \textcolor{input}{location}, & - w m \\
~~~~~\textcolor{input}{ptset\_type}, \textcolor{input}{npnts}, \textcolor{input}{pnts}, \textcolor{output}{\textit{D}}, \textcolor{output}{\textit{ier}}) & \\
call cg\_discrete\_ptset\_info\_f(\textcolor{input}{fn}, \textcolor{input}{B}, \textcolor{input}{Z}, \textcolor{input}{D}, \textcolor{output}{\textit{ptset\_type}}, \textcolor{output}{\textit{npnts}}, \textcolor{output}{\textit{ier}}) & r - m \\
call cg\_discrete\_ptset\_read\_f(\textcolor{input}{fn}, \textcolor{input}{B}, \textcolor{input}{Z}, \textcolor{input}{D}, \textcolor{output}{\textit{pnts}}, \textcolor{output}{\textit{ier}}) & r - m \\
call cg\_discrete\_size\_f(\textcolor{input}{fn}, \textcolor{input}{B}, \textcolor{input}{Z}, \textcolor{input}{D}, \textcolor{output}{\textit{data\_dim}}, \textcolor{output}{\textit{dim\_vals}}, \textcolor{output}{\textit{ier}}) & r - m \\
\end{fctbox}

\noindent
\textbf{\textcolor{input}{Input}/\textcolor{output}{\textit{Output}}}

\noindent (Note that for Fortran calls, all integer arguments are integer*4 in 32-bit mode and integer*8 in 64-bit mode.
See ``64-bit Fortran Portability and Issues" section.)

\begin{Ventryi}{\texttt{DiscreteName}}\raggedright
\item [\texttt{fn}]
      CGNS file index number.
      (\textcolor{input}{Input})
\item [\texttt{B}]
      Base index number, where $1 \leq \text{\texttt{B}} \leq \text{\texttt{nbases}}$.
      (\textcolor{input}{Input})
\item [\texttt{Z}]
      Zone index number, where $1 \leq \text{\texttt{Z}} \leq \text{\texttt{nzones}}$.
      (\textcolor{input}{Input})
\item [\texttt{D}]
      Discrete data index number, where $1 \leq \text{\texttt{D}} \leq \text{\texttt{ndiscrete}}$.
      (\textcolor{input}{Input} for \texttt{cg\_discrete\_read};
      \textcolor{output}{\textit{output}} for \texttt{cg\_discrete\_write})
\item [\texttt{ndiscrete}]
      Number of \texttt{DiscreteData\_t} data structures under zone \texttt{Z}.
      (\textcolor{output}{\textit{Output}})
\item [\texttt{DiscreteName}]
      Name of \texttt{DiscreteData\_t} data structure.
      (\textcolor{input}{Input} for \texttt{cg\_discrete\_write};
      \textcolor{output}{\textit{output}} for \texttt{cg\_discrete\_read})
\item [\texttt{location}]
      Grid location where the discrete data is recorded.
      The current admissible locations are \texttt{Vertex},
      \texttt{CellCenter}, \texttt{IFaceCenter}, \texttt{JFaceCenter},
      and \texttt{KFaceCenter}.
      (\textcolor{input}{Input} for \texttt{cg\_discrete\_ptset\_write};
      \textcolor{output}{\textit{output}} for \texttt{cg\_sol\_info})
\item [\texttt{ptset\_type}]
      Type of point set defining the interface for the discrete data;
      either \texttt{PointRange} or \texttt{PointList}.
      (\textcolor{input}{Input} for \texttt{cg\_sol\_ptset\_write};
      \textcolor{output}{\textit{output}} for
      \texttt{cg\_sol\_ptset\_info})
\item [\texttt{npnts}]
      Number of points defining the interface for the discrete data.
      For a \texttt{ptset\_type} of \texttt{PointRange}, npnts
      is always two. For a \texttt{ptset\_type} of \texttt{PointList},
      npnts is the number of points in the list. 
      (\textcolor{input}{Input} for \texttt{cg\_sol\_ptset\_write};
      \textcolor{output}{\textit{output}} for
      \texttt{cg\_sol\_ptset\_info})
\item [\texttt{pnts}]
      Array of points defining the interface for the discrete data.
      (\textcolor{input}{Input} for \texttt{cg\_sol\_ptset\_write};
      \textcolor{output}{\textit{output}} for
      \texttt{cg\_sol\_ptset\_read})
\item [\texttt{data\_dim}]
      Number of dimensions defining the discrete data.
      If a point set has been defined, this will be 1,
      otherwise this will be the current zone index dimension.
      (\textcolor{output}{\textit{Output}})
\item [\texttt{dim\_vals}]
      The array of \texttt{data\_dim} dimensions for the discrete data.
      (\textcolor{output}{\textit{Output}})
\item [\texttt{ier}]
      Error status.
      (\textcolor{output}{\textit{Output}})
\end{Ventryi}

\texttt{DiscreteData\_t} nodes are intended for the storage of fields
of data not usually identified as part of the flow solution, such as
fluxes or equation residuals.

The description for these functions is similar to the
\texttt{FlowSolution\_t} node (\autoref{s:flowsolution}).
To read and write the discrete data, use
\texttt{cg\_goto} (\autoref{s:navigating}) to access the
\texttt{DiscreteData\_t} node, then
\texttt{cg\_array\_read} and \texttt{cg\_array\_write} (\autoref{s:dataarray}).

\newpage
\subsection{Zone Subregions}
\label{s:subregion}

\noindent
\textit{Node}: \texttt{ZoneSubRegion\_t}

\begin{fctbox}
\textcolor{output}{\textit{ier}} = cg\_nsubregs(\textcolor{input}{int fn}, \textcolor{input}{int B}, \textcolor{input}{int Z}, \textcolor{output}{\textit{int *nsubregs}}); & r - m \\
\textcolor{output}{\textit{ier}} = cg\_subreg\_info(\textcolor{input}{int fn}, \textcolor{input}{int B}, \textcolor{input}{int Z}, \textcolor{input}{int S}, \textcolor{output}{\textit{char *regname}}, & r - m \\
~~~~~~\textcolor{output}{\textit{int *dimension}}, \textcolor{output}{\textit{GridLocation\_t *location}}, & \\
~~~~~~\textcolor{output}{\textit{PointSetType\_t *ptset\_type}}, \textcolor{output}{\textit{cgsize\_t *npnts}}, & \\
~~~~~~\textcolor{output}{\textit{int *bcname\_len}}, \textcolor{output}{\textit{int *bcname\_len}}); & \\
\textcolor{output}{\textit{ier}} = cg\_subreg\_ptset\_read(\textcolor{input}{int fn}, \textcolor{input}{int B}, \textcolor{input}{int Z}, \textcolor{input}{int S}, & r - m \\
~~~~~~\textcolor{output}{\textit{cgsize\_t *pnts}}); & \\
\textcolor{output}{\textit{ier}} = cg\_subreg\_bcname\_read(\textcolor{input}{int fn}, \textcolor{input}{int B}, \textcolor{input}{int Z}, \textcolor{input}{int S}, & r - m \\
~~~~~~\textcolor{output}{\textit{char *bcname}}); & \\
\textcolor{output}{\textit{ier}} = cg\_subreg\_gcname\_read(\textcolor{input}{int fn}, \textcolor{input}{int B}, \textcolor{input}{int Z}, \textcolor{input}{int S}, & r - m \\
~~~~~~\textcolor{output}{\textit{char *gcname}}); & \\
\textcolor{output}{\textit{ier}} = cg\_subreg\_ptset\_write(\textcolor{input}{int fn}, \textcolor{input}{int B}, \textcolor{input}{int Z}, \textcolor{input}{const char *regname}, & - w m \\
~~~~~~\textcolor{input}{int dimension}, \textcolor{input}{GridLocation\_t location}, & \\
~~~~~~\textcolor{input}{PointSetType\_t ptset\_type}, \textcolor{input}{cgsize\_t npnts}, & \\
~~~~~~\textcolor{input}{cgsize\_t *pnts}, \textcolor{output}{\textit{int *S}}); & \\
\textcolor{output}{\textit{ier}} = cg\_subreg\_bcname\_write(\textcolor{input}{int fn}, \textcolor{input}{int B}, \textcolor{input}{int Z}, \textcolor{input}{const char *regname}, & - w m \\
~~~~~~\textcolor{input}{int dimension}, \textcolor{input}{const char *bcname}, \textcolor{output}{\textit{int *S}}); & \\
\textcolor{output}{\textit{ier}} = cg\_subreg\_gcname\_write(\textcolor{input}{int fn}, \textcolor{input}{int B}, \textcolor{input}{int Z}, \textcolor{input}{const char *regname}, & - w m \\
~~~~~~\textcolor{input}{int dimension}, \textcolor{input}{const char *gcname}, \textcolor{output}{\textit{int *S}}); & \\
\hline
call cg\_nsubregs\_f(\textcolor{input}{fn}, \textcolor{input}{B}, \textcolor{input}{Z}, \textcolor{output}{\textit{nsubregs}}, \textcolor{output}{\textit{ier}}) & r - m \\
call cg\_subreg\_info\_f(\textcolor{input}{fn}, \textcolor{input}{B}, \textcolor{input}{Z}, \textcolor{input}{S}, \textcolor{output}{\textit{regname}}, \textcolor{output}{\textit{dimension}}, \textcolor{output}{\textit{location}}, & r - m \\
~~~~~\textcolor{output}{\textit{ptset\_type}}, \textcolor{output}{\textit{npnts}}, \textcolor{output}{\textit{bcname\_len}}, \textcolor{output}{\textit{bcname\_len}}, \textcolor{output}{\textit{ier}}) & \\
call cg\_subreg\_ptset\_read\_f(\textcolor{input}{fn}, \textcolor{input}{B}, \textcolor{input}{Z}, \textcolor{input}{S}, \textcolor{output}{\textit{pnts}}, \textcolor{output}{\textit{ier}}) & r - m \\
call cg\_subreg\_bcname\_read\_f(\textcolor{input}{fn}, \textcolor{input}{B}, \textcolor{input}{Z}, \textcolor{input}{S}, \textcolor{output}{\textit{bcname}}, \textcolor{output}{\textit{ier}}) & r - m \\
call cg\_subreg\_gcname\_read\_f(\textcolor{input}{fn}, \textcolor{input}{B}, \textcolor{input}{Z}, \textcolor{input}{S}, \textcolor{output}{\textit{gcname}}, \textcolor{output}{\textit{ier}}) & r - m \\
call cg\_subreg\_ptset\_write\_f(\textcolor{input}{fn}, \textcolor{input}{B}, \textcolor{input}{Z}, \textcolor{input}{regname},  \textcolor{input}{dimension}, \textcolor{input}{location}, & - w m \\
~~~~~\textcolor{input}{ptset\_type}, \textcolor{input}{npnts}, \textcolor{input}{pnts}, \textcolor{output}{\textit{S}}, \textcolor{output}{\textit{ier}}) & \\
call cg\_subreg\_bcname\_write\_f(\textcolor{input}{fn}, \textcolor{input}{B}, \textcolor{input}{Z}, \textcolor{input}{regname}, \textcolor{input}{dimension}, \textcolor{input}{bcname}, & - w m \\
~~~~~\textcolor{output}{\textit{S}}, \textcolor{output}{\textit{ier}}) & \\
call cg\_subreg\_gcname\_write\_f(\textcolor{input}{fn}, \textcolor{input}{B}, \textcolor{input}{Z}, \textcolor{input}{regname}, \textcolor{input}{dimension}, \textcolor{input}{gcname}, & - w m \\
~~~~~\textcolor{output}{\textit{S}}, \textcolor{output}{\textit{ier}}) & \\
\end{fctbox}

\noindent
\textbf{\textcolor{input}{Input}/\textcolor{output}{\textit{Output}}}

\noindent (Note that for Fortran calls, all integer arguments are integer*4 in 32-bit mode and integer*8 in 64-bit mode.
See ``64-bit Fortran Portability and Issues" section.)

\begin{Ventryi}{\texttt{ptset\_type}}\raggedright
\item [\texttt{fn}]
      CGNS file index number.
      (\textcolor{input}{Input})
\item [\texttt{B}]
      Base index number, where $1 \leq \text{\texttt{B}} \leq \text{\texttt{nbases}}$.
      (\textcolor{input}{Input})
\item [\texttt{Z}]
      Zone index number, where $1 \leq \text{\texttt{Z}} \leq \text{\texttt{nzones}}$.
      (\textcolor{input}{Input})
\item [\texttt{S}]
      ZoneSubRegion index number, where $1 \leq \text{\texttt{S}} \leq \text{\texttt{nsubregs}}$.
      (\textcolor{input}{Input} for \texttt{cg\_subreg\_info},
      \texttt{cg\_subreg\_ptset\_read}, \texttt{cg\_subreg\_bcname\_read} and
      \texttt{cg\_subreg\_gcname\_read};
      \textcolor{output}{\textit{output}} for \texttt{cg\_subreg\_ptset\_write},
      \texttt{cg\_subreg\_bcname\_write} and \texttt{cg\_subreg\_gcname\_write})
\item [\texttt{nsubregs}]
      Number of \texttt{ZoneSubRegion\_t} nodes under Zone \texttt{Z}.
      (\textcolor{output}{\textit{Output}})
\item [\texttt{regname}]
      Name of the \texttt{ZoneSubRegion\_t} node.
      (\textcolor{input}{Input} for \texttt{cg\_subreg\_ptset\_write},
      \texttt{cg\_subreg\_bcname\_write}, and
      \texttt{cg\_subreg\_gcname\_write};
      \textcolor{output}{\textit{output}} for \texttt{cg\_subreg\_info})
\item [\texttt{dimension}]
      Dimensionality of the subregion. 1 for lines, 2 for faces,
      3 for volumes.
      (\textcolor{input}{Input} for \texttt{cg\_subreg\_ptset\_write},
      \texttt{cg\_subreg\_bcname\_write}, and
      \texttt{cg\_subreg\_gcname\_write};
      \textcolor{output}{\textit{output}} for \texttt{cg\_subreg\_info})
\item [\texttt{location}]
      Grid location used in the definition of the point set.
      The current admissible locations are \texttt{Vertex} and
      \texttt{CellCenter}.
      (\textcolor{input}{Input} for \texttt{cg\_subreg\_ptset\_write};
      \textcolor{output}{\textit{output}} for \texttt{cg\_subreg\_info})
\item [\texttt{ptset\_type}]
      Type of point set defining the interface for the subregion data;
      either \texttt{PointRange} or \texttt{PointList}.
      (\textcolor{input}{Input} for \texttt{cg\_subreg\_ptset\_write};
      \textcolor{output}{\textit{output}} for
      \texttt{cg\_subreg\_info})
\item [\texttt{npnts}]
      Number of points defining the interface for the subregion data.
      For a \texttt{ptset\_type} of \texttt{PointRange}, npnts
      is always two. For a \texttt{ptset\_type} of \texttt{PointList},
      npnts is the number of points in the list. 
      (\textcolor{input}{Input} for \texttt{cg\_subreg\_ptset\_write};
      \textcolor{output}{\textit{output}} for
      \texttt{cg\_subreg\_info})
\item [\texttt{pnts}]
      Array of points defining the interface for the subregion data.
      (\textcolor{input}{Input} for \texttt{cg\_subreg\_ptset\_write};
      \textcolor{output}{\textit{output}} for
      \texttt{cg\_subreg\_ptset\_read})
\item [\texttt{bcname}]
      The name of a \texttt{BC\_t} node which defines the subregion.
      (\textcolor{input}{Input} for \texttt{cg\_subreg\_bcname\_write};
      \textcolor{output}{\textit{output}} for
      \texttt{cg\_subreg\_bcname\_read})
\item [\texttt{gcname}]
      The name of a \texttt{GridConnectivity\_t} or
      \texttt{GridConnectivity1to1\_t} node which defines the subregion.
      (\textcolor{input}{Input} for \texttt{cg\_subreg\_gcname\_write};
      \textcolor{output}{\textit{output}} for
      \texttt{cg\_subreg\_gcname\_read})
\item [\texttt{bcname\_len}]
      String length of \texttt{bcname}.
      (\textcolor{output}{\textit{Output}})
\item [\texttt{gcname\_len}]
      String length of \texttt{gcname}.
      (\textcolor{output}{\textit{Output}})
\item [\texttt{ier}]
      Error status.
      (\textcolor{output}{\textit{Output}})
\end{Ventryi}

These functions allow for the specification of Zone subregions. The
subregion may be specified as either the name of an existing
\texttt{BC\_t} node (\texttt{cg\_subreg\_bcname\_write}), an
existing \texttt{GridConnectivity\_t} or \texttt{GridConnectivity1to1\_t}
node (\texttt{cg\_subreg\_gcname\_write}), or as a
\texttt{PointSet/PointRange} (\texttt{cg\_subreg\_ptset\_write}).
These specifications are mutually exclusive.
To determine the type of the subregion, use
\texttt{cg\_subreg\_info}. If the subregion is a point set,
then \texttt{ptset\_type} will indicate the point set type (either
\texttt{PointList} or \texttt{PointRange}) and \texttt{npts} will be
set to the number of points to define the region. Otherwise,
\texttt{ptset\_type} will be set to \texttt{CG\_Null} and
\texttt{npnts} will be 0. In this case, one of \texttt{bcname\_len}
or \texttt{gcname\_len} will be non-zero, indicating whether the
subregion references a \texttt{BC\_t} node
(\texttt{bcname\_len} non-zero) or
\texttt{GridConnectivity\_t} node (\texttt{gcname\_len} non-zero).



\clearemptydoublepage
\section{Grid Connectivity}
\label{s:connectivity}
\thispagestyle{plain}

\subsection{One-to-One Connectivity}
\label{s:1to1}

\noindent
\textit{Node}: \texttt{GridConnectivity1to1\_t}

\begin{fctbox}
\textcolor{output}{\textit{ier}} = cg\_n1to1\_global(\textcolor{input}{int fn}, \textcolor{input}{int B}, \textcolor{output}{\textit{int *n1to1\_global}}); & r - m \\
\textcolor{output}{\textit{ier}} = cg\_1to1\_read\_global(\textcolor{input}{int fn}, \textcolor{input}{int B}, \textcolor{output}{\textit{char **connectname}}, & r - m \\
~~~~~~\textcolor{output}{\textit{char **zonename}}, \textcolor{output}{\textit{char **donorname}}, \textcolor{output}{\textit{cgsize\_t **range}}, & \\
~~~~~~\textcolor{output}{\textit{cgsize\_t **donor\_range}}, \textcolor{output}{\textit{int **transform}}); & \\
\hline
call cg\_n1to1\_global\_f(\textcolor{input}{fn}, \textcolor{input}{B}, \textcolor{output}{\textit{n1to1\_global}}, \textcolor{output}{\textit{ier}}) & r - m \\
call cg\_1to1\_read\_global\_f(\textcolor{input}{fn}, \textcolor{input}{B}, \textcolor{output}{\textit{connectname}}, \textcolor{output}{\textit{zonename}}, \textcolor{output}{\textit{donorname}}, & r - m \\
~~~~~\textcolor{output}{\textit{range}}, \textcolor{output}{\textit{donor\_range}}, \textcolor{output}{\textit{transform}}, \textcolor{output}{\textit{ier}}) & \\
\end{fctbox}

\noindent
\textbf{\textcolor{input}{Input}/\textcolor{output}{\textit{Output}}}

\noindent (Note that for Fortran calls, all integer arguments are integer*4 in 32-bit mode and integer*8 in 64-bit mode.
See ``64-bit Fortran Portability and Issues" section.)

\begin{Ventryi}{\texttt{n1to1\_global}}\raggedright
\item [\texttt{fn}]
      CGNS file index number.
      (\textcolor{input}{Input})
\item [\texttt{B}]
      Base index number, where $1 \leq \text{\texttt{B}} \leq \text{\texttt{nbases}}$.
      (\textcolor{input}{Input})
\item [\texttt{n1to1\_global}]
      Total number of one-to-one interfaces in base \texttt{B},
      stored under \texttt{GridConnectivity1to1\_t} nodes.
      (I.e., this does not include one-to-one interfaces that may be
      stored under \texttt{GridConnectivity\_t} nodes, used for
      generalized zone interfaces.)
      Note that the function \texttt{cg\_n1to1} (described below) may be
      used to get the number of one-to-one interfaces in a specific
      zone.
      (\textcolor{output}{\textit{Output}})
\item [\texttt{connectname}]
      Name of the interface.
      (\textcolor{output}{\textit{Output}})
\item [\texttt{zonename}]
      Name of the first zone, for all one-to-one interfaces in base
      \texttt{B}.
      (\textcolor{output}{\textit{Output}})
\item [\texttt{donorname}]
      Name of the second zone, for all one-to-one interfaces in base
      \texttt{B}.
      (\textcolor{output}{\textit{Output}})
\item [\texttt{range}]
      Range of points for the first zone, for all one-to-one interfaces
      in base \texttt{B}.
      (\textcolor{output}{\textit{Output}})
\item [\texttt{donor\_range}]
      Range of points for the current zone, for all one-to-one interfaces
      in base \texttt{B}.
      (\textcolor{output}{\textit{Output}})
\item [\texttt{transform}]
      Short hand notation for the transformation matrix defining the
      relative orientation of the two zones.
      This transformation is given for all one-to-one interfaces in
      base \texttt{B}.
      See the description of \texttt{GridConnectivity1to1\_t} in the
      \href{../sids/sids.pdf}{SIDS manual} for details.
      (\textcolor{output}{\textit{Output}})
\item [\texttt{ier}]
      Error status.
      (\textcolor{output}{\textit{Output}})
\end{Ventryi}

The above functions may be used to get information about all the
one-to-one zone interfaces in a CGNS database.

\begin{fctbox}
\textcolor{output}{\textit{ier}} = cg\_1to1\_write(\textcolor{input}{int fn}, \textcolor{input}{int B}, \textcolor{input}{int Z}, \textcolor{input}{char *connectname}, & - w m \\
~~~~~~\textcolor{input}{char *donorname}, \textcolor{input}{cgsize\_t *range}, \textcolor{input}{cgsize\_t *donor\_range}, & \\
~~~~~~\textcolor{input}{int *transform}, \textcolor{output}{\textit{int *I}}); & \\
\textcolor{output}{\textit{ier}} = cg\_n1to1(\textcolor{input}{int fn}, \textcolor{input}{int B}, \textcolor{input}{int Z}, \textcolor{output}{\textit{int *n1to1}}); & r - m \\
\textcolor{output}{\textit{ier}} = cg\_1to1\_read(\textcolor{input}{int fn}, \textcolor{input}{int B}, \textcolor{input}{int Z}, \textcolor{input}{int I}, \textcolor{output}{\textit{char *connectname}}, & r - m \\
~~~~~~\textcolor{output}{\textit{char *donorname}}, \textcolor{output}{\textit{cgsize\_t *range}}, \textcolor{output}{\textit{cgsize\_t *donor\_range}}, & \\
~~~~~~\textcolor{output}{\textit{int *transform}}); & \\
\hline
call cg\_1to1\_write\_f(\textcolor{input}{fn}, \textcolor{input}{B}, \textcolor{input}{Z}, \textcolor{input}{connectname}, \textcolor{input}{donorname}, \textcolor{input}{range}, & - w m \\
~~~~~\textcolor{input}{donor\_range}, \textcolor{input}{transform}, \textcolor{output}{\textit{I}}, \textcolor{output}{\textit{ier}}) & \\
call cg\_n1to1\_f(\textcolor{input}{fn}, \textcolor{input}{B}, \textcolor{input}{Z}, \textcolor{output}{\textit{n1to1}}, \textcolor{output}{\textit{ier}}) & r - m \\
call cg\_1to1\_read\_f(\textcolor{input}{fn}, \textcolor{input}{B}, \textcolor{input}{Z}, \textcolor{input}{I}, \textcolor{output}{\textit{connectname}}, \textcolor{output}{\textit{donorname}}, \textcolor{output}{\textit{range}}, & r - m \\
~~~~~\textcolor{output}{\textit{donor\_range}}, \textcolor{output}{\textit{transform}}, \textcolor{output}{\textit{ier}}) & \\
\end{fctbox}

\noindent
\textbf{\textcolor{input}{Input}/\textcolor{output}{\textit{Output}}}

\noindent (Note that for Fortran calls, all integer arguments are integer*4 in 32-bit mode and integer*8 in 64-bit mode.
See ``64-bit Fortran Portability and Issues" section.)

\begin{Ventryi}{\texttt{connectname}}\raggedright
\item [\texttt{fn}]
      CGNS file index number.
      (\textcolor{input}{Input})
\item [\texttt{B}]
      Base index number, where $1 \leq \text{\texttt{B}} \leq \text{\texttt{nbases}}$.
      (\textcolor{input}{Input})
\item [\texttt{Z}]
      Zone index number, where $1 \leq \text{\texttt{Z}} \leq \text{\texttt{nzones}}$.
      (\textcolor{input}{Input})
\item [\texttt{I}]
      Interface index number, where $1 \leq \text{\texttt{I}} \leq \text{\texttt{n1to1}}$.
      (\textcolor{input}{Input} for \texttt{cg\_1to1\_read};
      \textcolor{output}{\textit{output}} for \texttt{cg\_1to1\_write})
\item [\texttt{n1to1}]
      Number of one-to-one interfaces in zone \texttt{Z}, stored under
      \texttt{GridConnectivity1to1\_t} nodes.
      (I.e., this does not include one-to-one interfaces that may be
      stored under \texttt{GridConnectivity\_t} nodes, used for
      generalized zone interfaces.)
      (\textcolor{output}{\textit{Output}})
\item [\texttt{connectname}]
      Name of the interface.
      (\textcolor{input}{Input} for \texttt{cg\_1to1\_write};
      \textcolor{output}{\textit{output}} for \texttt{cg\_1to1\_read})
\item [\texttt{donorname}]
      Name of the zone interfacing with the current zone.
      (\textcolor{input}{Input} for \texttt{cg\_1to1\_write};
      \textcolor{output}{\textit{output}} for \texttt{cg\_1to1\_read})
\item [\texttt{range}]
      Range of points for the current zone.
      (\textcolor{input}{Input} for \texttt{cg\_1to1\_write};
      \textcolor{output}{\textit{output}} for \texttt{cg\_1to1\_read})
\item [\texttt{donor\_range}]
      Range of points for the donor zone.
      (\textcolor{input}{Input} for \texttt{cg\_1to1\_write};
      \textcolor{output}{\textit{output}} for \texttt{cg\_1to1\_read})
\item [\texttt{transform}]
      Short hand notation for the transformation matrix defining the
      relative orientation of the two zones.
      See the description of \texttt{GridConnectivity1to1\_t} in the
      \href{../sids/sids.pdf}{SIDS manual} for details.
      (\textcolor{input}{Input} for \texttt{cg\_1to1\_write};
      \textcolor{output}{\textit{output}} for \texttt{cg\_1to1\_read})
\item [\texttt{ier}]
      Error status.
      (\textcolor{output}{\textit{Output}})
\end{Ventryi}

The above functions are used to read and write one-to-one connectivity
data for a specific zone.

\newpage
\subsection{Generalized Connectivity}
\label{s:generalized}

\noindent
\textit{Node}: \texttt{GridConnectivity\_t}

\begin{fctbox}
\textcolor{output}{\textit{ier}} = cg\_conn\_write(\textcolor{input}{int fn}, \textcolor{input}{int B}, \textcolor{input}{int Z}, \textcolor{input}{char *connectname}, & - w m \\
~~~~~~\textcolor{input}{GridLocation\_t location}, \textcolor{input}{GridConnectivityType\_t connect\_type}, & \\
~~~~~~\textcolor{input}{PointSetType\_t ptset\_type}, \textcolor{input}{cgsize\_t npnts}, \textcolor{input}{cgsize\_t *pnts}, & \\
~~~~~~\textcolor{input}{char *donorname}, \textcolor{input}{ZoneType\_t donor\_zonetype}, & \\
~~~~~~\textcolor{input}{PointSetType\_t donor\_ptset\_type}, \textcolor{input}{DataType\_t donor\_datatype}, & \\
~~~~~~\textcolor{input}{cgsize\_t ndata\_donor}, \textcolor{input}{cgsize\_t *donor\_data}, \textcolor{output}{\textit{int *I}}); & \\
\textcolor{output}{\textit{ier}} = cg\_conn\_write\_short(\textcolor{input}{int fn}, \textcolor{input}{int B}, \textcolor{input}{int Z}, \textcolor{input}{char *connectname}, & - w m \\
~~~~~~\textcolor{input}{GridLocation\_t location}, \textcolor{input}{GridConnectivityType\_t connect\_type}, & \\
~~~~~~\textcolor{input}{PointSetType\_t ptset\_type}, \textcolor{input}{cgsize\_t npnts}, \textcolor{input}{cgsize\_t *pnts}, & \\
~~~~~~\textcolor{input}{char *donorname}, \textcolor{output}{\textit{int *I}}); & \\
\textcolor{output}{\textit{ier}} = cg\_nconns(\textcolor{input}{int fn}, \textcolor{input}{int B}, \textcolor{input}{int Z}, \textcolor{output}{\textit{int *nconns}}); & r - m \\
\textcolor{output}{\textit{ier}} = cg\_conn\_info(\textcolor{input}{int fn}, \textcolor{input}{int B}, \textcolor{input}{int Z}, \textcolor{input}{int I}, \textcolor{output}{\textit{char *connectname}}, & r - m \\
~~~~~~\textcolor{output}{\textit{GridLocation\_t *location}}, \textcolor{output}{\textit{GridConnectivityType\_t *connect\_type}}, & \\
~~~~~~\textcolor{output}{\textit{PointSetType\_t *ptset\_type}}, \textcolor{output}{\textit{cgsize\_t *npnts}}, \textcolor{output}{\textit{char *donorname}}, & \\
~~~~~~\textcolor{output}{\textit{ZoneType\_t *donor\_zonetype}}, \textcolor{output}{\textit{PointSetType\_t *donor\_ptset\_type}}, & \\
~~~~~~\textcolor{output}{\textit{DataType\_t *donor\_datatype}}, \textcolor{output}{\textit{cgsize\_t *ndata\_donor}}); & \\
\textcolor{output}{\textit{ier}} = cg\_conn\_read(\textcolor{input}{int fn}, \textcolor{input}{int B}, \textcolor{input}{int Z}, \textcolor{input}{int I}, \textcolor{output}{\textit{cgsize\_t *pnts}}, & r - m \\
~~~~~~\textcolor{input}{DataType\_t donor\_datatype}, \textcolor{output}{\textit{cgsize\_t *donor\_data}}); & \\
\textcolor{output}{\textit{ier}} = cg\_conn\_read\_short(\textcolor{input}{int fn}, \textcolor{input}{int B}, \textcolor{input}{int Z}, \textcolor{input}{int I}, \textcolor{output}{\textit{cgsize\_t *pnts}}); & r - m \\
\hline
call cg\_conn\_write\_f(\textcolor{input}{fn}, \textcolor{input}{B}, \textcolor{input}{Z}, \textcolor{input}{connectname}, \textcolor{input}{location}, \textcolor{input}{connect\_type}, & - w m \\
~~~~~\textcolor{input}{ptset\_type}, \textcolor{input}{npnts}, \textcolor{input}{pnts}, \textcolor{input}{donorname}, \textcolor{input}{donor\_zonetype}, & \\
~~~~~\textcolor{input}{donor\_ptset\_type}, \textcolor{input}{donor\_datatype}, \textcolor{input}{ndata\_donor}, \textcolor{input}{donor\_data}, \textcolor{output}{\textit{I}}, & \\
~~~~~\textcolor{output}{\textit{ier}}) & \\
call cg\_conn\_write\_short\_f(\textcolor{input}{fn}, \textcolor{input}{B}, \textcolor{input}{Z}, \textcolor{input}{connectname}, \textcolor{input}{location}, & - w m \\
~~~~~\textcolor{input}{connect\_type}, \textcolor{input}{ptset\_type}, \textcolor{input}{npnts}, \textcolor{input}{pnts}, \textcolor{input}{donorname}, \textcolor{output}{\textit{I}}, \textcolor{output}{\textit{ier}}) & \\
call cg\_nconns\_f(\textcolor{input}{fn}, \textcolor{input}{B}, \textcolor{input}{Z}, \textcolor{output}{\textit{nconns}}, \textcolor{output}{\textit{ier}}) & r - m \\
call cg\_conn\_info\_f(\textcolor{input}{fn}, \textcolor{input}{B}, \textcolor{input}{Z}, \textcolor{input}{I}, \textcolor{output}{\textit{connectname}}, \textcolor{output}{\textit{location}}, \textcolor{output}{\textit{connect\_type}}, & r - m \\
~~~~~\textcolor{output}{\textit{ptset\_type}}, \textcolor{output}{\textit{npnts}}, \textcolor{output}{\textit{donorname}}, \textcolor{output}{\textit{donor\_zonetype}}, \textcolor{output}{\textit{donor\_ptset\_type}}, & \\
~~~~~\textcolor{output}{\textit{donor\_datatype}}, \textcolor{output}{\textit{ndata\_donor}}, \textcolor{output}{\textit{ier}}) & \\
call cg\_conn\_read\_f(\textcolor{input}{fn}, \textcolor{input}{B}, \textcolor{input}{Z}, \textcolor{input}{I}, \textcolor{output}{\textit{pnts}}, \textcolor{input}{donor\_datatype}, \textcolor{output}{\textit{donor\_data}}, & r - m \\
~~~~~\textcolor{output}{\textit{ier}}) & \\
call cg\_conn\_read\_short\_f(\textcolor{input}{fn}, \textcolor{input}{B}, \textcolor{input}{Z}, \textcolor{input}{I}, \textcolor{output}{\textit{pnts}}, \textcolor{output}{\textit{ier}}) & \\
\end{fctbox}

\noindent
\textbf{\textcolor{input}{Input}/\textcolor{output}{\textit{Output}}}

\noindent (Note that for Fortran calls, all integer arguments are integer*4 in 32-bit mode and integer*8 in 64-bit mode.
See ``64-bit Fortran Portability and Issues" section.)

\begin{Ventryi}{\texttt{donor\_ptset\_type}}\raggedright
\item [\texttt{fn}]
      CGNS file index number.
      (\textcolor{input}{Input})
\item [\texttt{B}]
      Base index number, where $1 \leq \text{\texttt{B}} \leq \text{\texttt{nbases}}$.
      (\textcolor{input}{Input})
\item [\texttt{Z}]
      Zone index number, where $1 \leq \text{\texttt{Z}} \leq \text{\texttt{nzones}}$.
      (\textcolor{input}{Input})
\item [\texttt{I}]
      Discrete data index number, where $1 \leq \text{\texttt{I}} \leq \text{\texttt{nconns}}$.
      (\textcolor{input}{Input} for \texttt{cg\_conn\_info},
      \texttt{cg\_conn\_read};
      \textcolor{output}{\textit{output}} for \texttt{cg\_conn\_write},
      \texttt{cg\_conn\_write\_short})
\item [\texttt{nconns}]
      Number of interfaces for zone \texttt{Z}.
      (\textcolor{output}{\textit{Output}})
\item [\texttt{connectname}]
      Name of the interface.
      (\textcolor{input}{Input} for \texttt{cg\_conn\_write},
      \texttt{cg\_conn\_write\_short};
      \textcolor{output}{\textit{output}} for \texttt{cg\_conn\_info})
\item [\texttt{location}]
      Grid location used in the definition of the point set.
      The currently admissible locations are \texttt{Vertex} and
      \texttt{CellCenter}.
      (\textcolor{input}{Input} for \texttt{cg\_conn\_write},
      \texttt{cg\_conn\_write\_short};
      \textcolor{output}{\textit{output}} for \texttt{cg\_conn\_info})
\item [\texttt{connect\_type}]
      Type of interface being defined.
      The admissible types are \texttt{Overset}, \texttt{Abutting},
      and \texttt{Abutting1to1}.
      (\textcolor{input}{Input} for \texttt{cg\_conn\_write},
      \texttt{cg\_conn\_write\_short};
      \textcolor{output}{\textit{output}} for \texttt{cg\_conn\_info})
\item [\texttt{ptset\_type}]
      Type of point set defining the interface in the current zone;
      either \texttt{PointRange} or \texttt{PointList}.
      (\textcolor{input}{Input} for \texttt{cg\_conn\_write},
      \texttt{cg\_conn\_write\_short};
      \textcolor{output}{\textit{output}} for \texttt{cg\_conn\_info})
\item [\texttt{donor\_ptset\_type}]
      Type of point set defining the interface in the donor zone;
      either \texttt{PointListDonor} or \texttt{CellListDonor}.
      (\textcolor{input}{Input} for \texttt{cg\_conn\_write};
      \textcolor{output}{\textit{output}} for \texttt{cg\_conn\_info})
\item [\texttt{npnts}]
      Number of points defining the interface in the current zone.
      For a \texttt{ptset\_type} of \texttt{PointRange}, \texttt{npnts}
      is always two.
      For a \texttt{ptset\_type} of \texttt{PointList}, \texttt{npnts} is
      the number of points in the \texttt{PointList}.
      (\textcolor{input}{Input} for \texttt{cg\_conn\_write},
      \texttt{cg\_conn\_write\_short};
      \textcolor{output}{\textit{output}} for \texttt{cg\_conn\_info})
\item [\texttt{ndata\_donor}]
      Number of points or cells in the current zone.
      These are paired with points, cells, or fractions thereof in the
      donor zone.
      (\textcolor{input}{Input} for \texttt{cg\_conn\_write};
      \textcolor{output}{\textit{output}} for \texttt{cg\_conn\_info})
\item [\texttt{donorname}]
      Name of the zone interfacing with the current zone.
      (\textcolor{input}{Input} for \texttt{cg\_conn\_write};
      \textcolor{output}{\textit{output}} for \texttt{cg\_conn\_info})
\item [\texttt{donor\_datatype}]
      Data type in which the donor points are stored in the file.
      As of Version 3.0, this value is ignored when writing, and on
      reading it will return either \texttt{Integer} or \texttt{LongInteger}
      depending on whether the file was written using 32 or 64-bit.
      The \texttt{donor\_datatype} argument was left in these functions
      only for backward compatibility. The donot data is always read
      as \texttt{cgsize\_t}.
      (\textcolor{input}{Input} for \texttt{cg\_conn\_write},
      \texttt{cg\_conn\_read};
      \textcolor{output}{\textit{output}} for \texttt{cg\_conn\_info})
\item [\texttt{pnts}]
      Array of points defining the interface in the current zone.
      (\textcolor{input}{Input} for \texttt{cg\_conn\_write},
      \texttt{cg\_conn\_write\_short};
      \textcolor{output}{\textit{output}} for \texttt{cg\_conn\_read})
\item [\texttt{donor\_data}]
      Array of donor points or cells corresponding to
      \texttt{ndata\_donor}.
      Note that it is possible that the same donor point or cell may
      be used multiple times.
      (\textcolor{input}{Input} for \texttt{cg\_conn\_write};
      \textcolor{output}{\textit{output}} for \texttt{cg\_conn\_read})
\item [\texttt{donor\_zonetype}]
      Type of the donor zone.
      The admissible types are \texttt{Structured} and \texttt{Unstructured}.
      (\textcolor{input}{Input} for \texttt{cg\_conn\_write};
      \textcolor{output}{\textit{output}} for \texttt{cg\_conn\_info})
\item [\texttt{ier}]
      Error status.
      (\textcolor{output}{\textit{Output}})
\end{Ventryi}

Note that the interpolation factors stored
in the \texttt{InterpolantsDonor} data array are accessed using the
\texttt{cg\_goto} and \texttt{cg\_array\_\textit{xxx}} functions, described
in \autoref{s:navigating} and \autoref{s:dataarray}, respectively.

\newpage
\subsection{Special Grid Connectivity Properties}
\label{s:connproperty}

\noindent
\textit{Node}: \texttt{GridConnectivityProperty\_t}

\begin{fctbox}
\textcolor{output}{\textit{ier}} = cg\_conn\_periodic\_write(\textcolor{input}{int fn}, \textcolor{input}{int B}, \textcolor{input}{int Z}, \textcolor{input}{int I}, & - w m \\
~~~~~~\textcolor{input}{float *RotationCenter}, \textcolor{input}{float *RotationAngle}, \textcolor{input}{float *Translation}); & \\
\textcolor{output}{\textit{ier}} = cg\_conn\_average\_write(\textcolor{input}{int fn}, \textcolor{input}{int B}, \textcolor{input}{int Z}, \textcolor{input}{int I}, & - w m \\
~~~~~~\textcolor{input}{AverageInterfaceType\_t AverageInterfaceType}); & \\
\textcolor{output}{\textit{ier}} = cg\_1to1\_periodic\_write(\textcolor{input}{int fn}, \textcolor{input}{int B}, \textcolor{input}{int Z}, \textcolor{input}{int I}, & - w m \\
~~~~~~\textcolor{input}{float *RotationCenter}, \textcolor{input}{float *RotationAngle}, \textcolor{input}{float *Translation}); & \\
\textcolor{output}{\textit{ier}} = cg\_1to1\_average\_write(\textcolor{input}{int fn}, \textcolor{input}{int B}, \textcolor{input}{int Z}, \textcolor{input}{int I}, & - w m \\
~~~~~~\textcolor{input}{AverageInterfaceType\_t AverageInterfaceType}); & \\
\textcolor{output}{\textit{ier}} = cg\_conn\_periodic\_read(\textcolor{input}{int fn}, \textcolor{input}{int B}, \textcolor{input}{int Z}, \textcolor{input}{int I}, & r - m \\
~~~~~~\textcolor{output}{\textit{float *RotationCenter}}, \textcolor{output}{\textit{float *RotationAngle}}, & \\
~~~~~~\textcolor{output}{\textit{float *Translation}}); & \\
\textcolor{output}{\textit{ier}} = cg\_conn\_average\_read(\textcolor{input}{int fn}, \textcolor{input}{int B}, \textcolor{input}{int Z}, \textcolor{input}{int I}, & r - m \\
~~~~~~\textcolor{output}{\textit{AverageInterfaceType\_t *AverageInterfaceType}}); & \\
\textcolor{output}{\textit{ier}} = cg\_1to1\_periodic\_read(\textcolor{input}{int fn}, \textcolor{input}{int B}, \textcolor{input}{int Z}, \textcolor{input}{int I}, & r - m \\
~~~~~~\textcolor{output}{\textit{float *RotationCenter}}, \textcolor{output}{\textit{float *RotationAngle}}, & \\
~~~~~~\textcolor{output}{\textit{float *Translation}}); & \\
\textcolor{output}{\textit{ier}} = cg\_1to1\_average\_read(\textcolor{input}{int fn}, \textcolor{input}{int B}, \textcolor{input}{int Z}, \textcolor{input}{int I}, & r - m \\
~~~~~~\textcolor{output}{\textit{AverageInterfaceType\_t *AverageInterfaceType}}); & \\
\hline
call cg\_conn\_periodic\_write\_f(\textcolor{input}{fn}, \textcolor{input}{B}, \textcolor{input}{Z}, \textcolor{input}{I}, \textcolor{input}{RotationCenter}, & - w m \\
~~~~~\textcolor{input}{RotationAngle}, \textcolor{input}{Translation}, \textcolor{output}{\textit{ier}}) & \\
call cg\_conn\_average\_write\_f(\textcolor{input}{fn}, \textcolor{input}{B}, \textcolor{input}{Z}, \textcolor{input}{I}, \textcolor{input}{AverageInterfaceType}, \textcolor{output}{\textit{ier}}) & - w m \\
call cg\_1to1\_periodic\_write\_f(\textcolor{input}{fn}, \textcolor{input}{B}, \textcolor{input}{Z}, \textcolor{input}{I}, \textcolor{input}{RotationCenter}, & - w m \\
~~~~~\textcolor{input}{RotationAngle}, \textcolor{input}{Translation}, \textcolor{output}{\textit{ier}}) & \\
call cg\_1to1\_average\_write\_f(\textcolor{input}{fn}, \textcolor{input}{B}, \textcolor{input}{Z}, \textcolor{input}{I}, \textcolor{input}{AverageInterfaceType}, \textcolor{output}{\textit{ier}}) & - w m \\
call cg\_conn\_periodic\_read\_f(\textcolor{input}{fn}, \textcolor{input}{B}, \textcolor{input}{Z}, \textcolor{input}{I}, \textcolor{output}{\textit{RotationCenter}}, & r - m \\
~~~~~\textcolor{output}{\textit{RotationAngle}}, \textcolor{output}{\textit{Translation}}, \textcolor{output}{\textit{ier}}) & \\
call cg\_conn\_average\_read\_f(\textcolor{input}{fn}, \textcolor{input}{B}, \textcolor{input}{Z}, \textcolor{input}{I}, \textcolor{output}{\textit{AverageInterfaceType}}, \textcolor{output}{\textit{ier}}) & r - m \\
call cg\_1to1\_periodic\_read\_f(\textcolor{input}{fn}, \textcolor{input}{B}, \textcolor{input}{Z}, \textcolor{input}{I}, \textcolor{output}{\textit{RotationCenter}}, & r - m \\
~~~~~\textcolor{output}{\textit{RotationAngle}}, \textcolor{output}{\textit{Translation}}, \textcolor{output}{\textit{ier}}) & \\
call cg\_1to1\_average\_read\_f(\textcolor{input}{fn}, \textcolor{input}{B}, \textcolor{input}{Z}, \textcolor{input}{I}, \textcolor{output}{\textit{AverageInterfaceType}}, \textcolor{output}{\textit{ier}}) & r - m \\
\end{fctbox}

\noindent
\textbf{\textcolor{input}{Input}/\textcolor{output}{\textit{Output}}}

\noindent (Note that for Fortran calls, all integer arguments are integer*4 in 32-bit mode and integer*8 in 64-bit mode.
See ``64-bit Fortran Portability and Issues" section.)

\begin{Ventryi}{\texttt{AverageInterfaceType}}\raggedright
\item [\texttt{fn}]
      CGNS file index number.
      (\textcolor{input}{Input})
\item [\texttt{B}]
      Base index number, where $1 \leq \text{\texttt{B}} \leq \text{\texttt{nbases}}$.
      (\textcolor{input}{Input})
\item [\texttt{Z}]
      Zone index number, where $1 \leq \text{\texttt{Z}} \leq \text{\texttt{nzones}}$.
      (\textcolor{input}{Input})
\item [\texttt{I}]
      Grid connectivity index number, where
      $1 \leq \text{\texttt{I}} \leq \text{\texttt{nconns}}$
      for the ``\texttt{cg\_conn}'' functions, and
      $1 \leq \text{\texttt{I}} \leq \text{\texttt{n1to1}}$
      for the ``\texttt{cg\_1to1}'' functions.
      (\textcolor{input}{Input})
\item [\texttt{RotationCenter}]
      An array of size \texttt{phys\_dim} defining the coordinates of
      the origin for defining the rotation angle between the periodic
      interfaces.
      (\texttt{phys\_dim} is the number of coordinates required to define
      a vector in the field.)
      (In Fortran, this is an array of Real*4 values.)
      (\textcolor{input}{Input} for \texttt{cg\_conn\_periodic\_write},
      \texttt{cg\_1to1\_periodic\_write};
      \textcolor{output}{\textit{output}} for \texttt{cg\_conn\_periodic\_read},
      \texttt{cg\_1to1\_periodic\_read})
\item [\texttt{RotationAngle}]
      An array of size \texttt{phys\_dim} defining the rotation angle from
      the current interface to the connecting interface.
      If rotating about more than one axis, the rotation is performed
      first about the x-axis, then the y-axis, then the z-axis.
      (In Fortran, this is an array of Real*4 values.)
      (\textcolor{input}{Input} for \texttt{cg\_conn\_periodic\_write},
      \texttt{cg\_1to1\_periodic\_write};
      \textcolor{output}{\textit{output}} for \texttt{cg\_conn\_periodic\_read},
      \texttt{cg\_1to1\_periodic\_read})
\item [\texttt{Translation}]
      An array of size \texttt{phys\_dim} defining the translation from
      the current interface to the connecting interface.
      (In Fortran, this is an array of Real*4 values.)
      (\textcolor{input}{Input} for \texttt{cg\_conn\_periodic\_write},
      \texttt{cg\_1to1\_periodic\_write};
      \textcolor{output}{\textit{output}} for \texttt{cg\_conn\_periodic\_read},
      \texttt{cg\_1to1\_periodic\_read})
\item [\texttt{AverageInterfaceType}]
      The type of averaging to be done.
      Valid types are \texttt{CG\_Null}, \texttt{CG\_UserDefined}, \texttt{AverageAll},
      \texttt{AverageCircumferential}, \texttt{AverageRadial}, \texttt{AverageI},
      \texttt{AverageJ}, and \texttt{AverageK}.
      (\textcolor{input}{Input} for \texttt{cg\_conn\_average\_write},
      \texttt{cg\_1to1\_average\_write};
      \textcolor{output}{\textit{output}} for \texttt{cg\_conn\_average\_read},
      \texttt{cg\_1to1\_average\_read})
\item [\texttt{ier}]
      Error status.
      (\textcolor{output}{\textit{Output}})
\end{Ventryi}

These functions may be used to store special grid connectivity properties.
The ``\texttt{cg\_conn}'' functions apply to generalized grid
connectivity nodes (i.e., \texttt{GridConnectivity\_t}), and the
``\texttt{cg\_1to1}'' functions apply to 1-to-1 grid connectivity nodes
(i.e., \texttt{GridConnectivity1to1\_t}).

The ``\texttt{write}'' functions will create the
\texttt{GridConnectivityProperty\_t} node if it doesn't already exist,
then add the appropriate connectivity property.
Multiple connectivity properties may be recorded under the same
\texttt{GridConnectivityProperty\_t} node.

The ``\texttt{read}'' functions will return with $\text{\texttt{ier}} = 2 =
\text{\texttt{CG\_NODE\_NOT\_FOUND}}$ if the requested connectivity property,
or the \texttt{GridConnectivityProperty\_t} node itself, doesn't exist.

\newpage
\subsection{Overset Holes}
\label{s:overset}

\noindent
\textit{Node}: \texttt{OversetHoles\_t}

\begin{fctbox}
\textcolor{output}{\textit{ier}} = cg\_hole\_write(\textcolor{input}{int fn}, \textcolor{input}{int B}, \textcolor{input}{int Z}, \textcolor{input}{char *holename}, & - w m \\
~~~~~~\textcolor{input}{GridLocation\_t location}, \textcolor{input}{PointSetType\_t ptset\_type}, \textcolor{input}{int nptsets}, & \\
~~~~~~\textcolor{input}{cgsize\_t npnts}, \textcolor{input}{cgsize\_t *pnts}, \textcolor{output}{\textit{int *I}}); & \\
\textcolor{output}{\textit{ier}} = cg\_nholes(\textcolor{input}{int fn}, \textcolor{input}{int B}, \textcolor{input}{int Z}, \textcolor{output}{\textit{int *nholes}}); & r - m \\
\textcolor{output}{\textit{ier}} = cg\_hole\_info(\textcolor{input}{int fn}, \textcolor{input}{int B}, \textcolor{input}{int Z}, \textcolor{input}{int I}, \textcolor{output}{\textit{char *holename}}, & r - m \\
~~~~~~\textcolor{output}{\textit{GridLocation\_t *location}}, \textcolor{output}{\textit{PointSetType\_t *ptset\_type}}, & \\
~~~~~~\textcolor{output}{\textit{int *nptsets}}, \textcolor{output}{\textit{cgsize\_t *npnts}}); & \\
\textcolor{output}{\textit{ier}} = cg\_hole\_read(\textcolor{input}{int fn}, \textcolor{input}{int B}, \textcolor{input}{int Z}, \textcolor{input}{int I}, \textcolor{output}{\textit{cgsize\_t *pnts}}); & r - m \\
\hline
call cg\_hole\_write\_f(\textcolor{input}{fn}, \textcolor{input}{B}, \textcolor{input}{Z}, \textcolor{input}{holename}, \textcolor{input}{location}, \textcolor{input}{ptset\_type}, & - w m \\
~~~~~\textcolor{input}{nptsets}, \textcolor{input}{npnts}, \textcolor{input}{pnts}, \textcolor{output}{\textit{I}}, \textcolor{output}{\textit{ier}}) & \\
call cg\_nholes\_f(\textcolor{input}{fn}, \textcolor{input}{B}, \textcolor{input}{Z}, \textcolor{output}{\textit{nholes}}, \textcolor{output}{\textit{ier}}) & r - m \\
call cg\_hole\_info\_f(\textcolor{input}{fn}, \textcolor{input}{B}, \textcolor{input}{Z}, \textcolor{input}{I}, \textcolor{output}{\textit{holename}}, \textcolor{output}{\textit{location}}, \textcolor{output}{\textit{ptset\_type}}, & r - m \\
~~~~~\textcolor{output}{\textit{nptsets}}, \textcolor{output}{\textit{npnts}}, \textcolor{output}{\textit{ier}}) & \\
call cg\_hole\_read\_f(\textcolor{input}{fn}, \textcolor{input}{B}, \textcolor{input}{Z}, \textcolor{input}{I}, \textcolor{output}{\textit{pnts}}, \textcolor{output}{\textit{ier}}) & r - m \\
\end{fctbox}

\noindent
\textbf{\textcolor{input}{Input}/\textcolor{output}{\textit{Output}}}

\noindent (Note that for Fortran calls, all integer arguments are integer*4 in 32-bit mode and integer*8 in 64-bit mode.
See ``64-bit Fortran Portability and Issues" section.)

\begin{Ventryi}{\texttt{ptset\_type}}\raggedright
\item [\texttt{fn}]
      CGNS file index number.
      (\textcolor{input}{Input})
\item [\texttt{B}]
      Base index number, where $1 \leq \text{\texttt{B}} \leq \text{\texttt{nbases}}$.
      (\textcolor{input}{Input})
\item [\texttt{Z}]
      Zone index number, where $1 \leq \text{\texttt{Z}} \leq \text{\texttt{nzones}}$.
      (\textcolor{input}{Input})
\item [\texttt{I}]
      Overset hole index number, where $1 \leq \text{\texttt{I}} \leq \text{\texttt{nholes}}$.
      (\textcolor{input}{Input} for \texttt{cg\_hole\_info},
      \texttt{cg\_hole\_read};
      \textcolor{output}{\textit{output}} for \texttt{cg\_hole\_write})
\item [\texttt{nholes}]
      Number of overset holes in zone \texttt{Z}.
      (\textcolor{output}{\textit{Output}})
\item [\texttt{holename}]
      Name of the overset hole.
      (\textcolor{input}{Input} for \texttt{cg\_hole\_write};
      \textcolor{output}{\textit{output}} for \texttt{cg\_hole\_info})
\item [\texttt{location}]
      Grid location used in the definition of the point set.
      The currently admissible locations are \texttt{Vertex} and
      \texttt{CellCenter}.
      (\textcolor{input}{Input} for \texttt{cg\_hole\_write};
      \textcolor{output}{\textit{output}} for \texttt{cg\_hole\_info})
\item [\texttt{ptset\_type}]
      The extent of the overset hole may be defined using a range of
      points or cells, or using a discrete list of all points or cells
      in the overset hole.
      If a range of points or cells is used, \texttt{ptset\_type} is set
      to \texttt{PointRange}.
      When a discrete list of points or cells is used,
      \texttt{ptset\_type} equals \texttt{PointList}.
      (\textcolor{input}{Input} for \texttt{cg\_hole\_write};
      \textcolor{output}{\textit{output}} for \texttt{cg\_hole\_info})
\item [\texttt{nptsets}]
      Number of point sets used to define the hole.
      If \texttt{ptset\_type} is \texttt{PointRange}, several point sets
      may be used.
      If \texttt{ptset\_type} is \texttt{PointList}, only one point set
      is allowed.
      (\textcolor{input}{Input} for \texttt{cg\_hole\_write};
      \textcolor{output}{\textit{output}} for \texttt{cg\_hole\_info})
\item [\texttt{npnts}]
      Number of points (or cells) in the point set.
      For a \texttt{ptset\_type} of \texttt{PointRange}, \texttt{npnts}
      is always two.
      For a \texttt{ptset\_type} of \texttt{PointList}, \texttt{npnts} is
      the number of points or cells in the \texttt{PointList}.
      (\textcolor{input}{Input} for \texttt{cg\_hole\_write};
      \textcolor{output}{\textit{output}} for \texttt{cg\_hole\_info})
\item [\texttt{pnts}]
      Array of points or cells in the point set.
      (\textcolor{input}{Input} for \texttt{cg\_hole\_write};
      \textcolor{output}{\textit{output}} for \texttt{cg\_hole\_read})
\item [\texttt{ier}]
      Error status.
      (\textcolor{output}{\textit{Output}})
\end{Ventryi}


\clearemptydoublepage
\section{Boundary Conditions}
\label{s:bc}
\thispagestyle{plain}

\subsection{Boundary Condition Type and Location}
\label{s:bctype}

\noindent
\textit{Node}: \texttt{BC\_t}

\begin{fctbox}
\textcolor{output}{\textit{ier}} = cg\_boco\_write(\textcolor{input}{int fn}, \textcolor{input}{int B}, \textcolor{input}{int Z}, \textcolor{input}{char *boconame}, & - w m \\
~~~~~~\textcolor{input}{BCType\_t bocotype}, \textcolor{input}{PointSetType\_t ptset\_type}, \textcolor{input}{cgsize\_t npnts}, & \\
~~~~~~\textcolor{input}{cgsize\_t *pnts}, \textcolor{output}{\textit{int *BC}}); & \\
\textcolor{output}{\textit{ier}} = cg\_boco\_normal\_write(\textcolor{input}{int fn}, \textcolor{input}{int B}, \textcolor{input}{int Z}, \textcolor{input}{int BC}, & - w m \\
~~~~~~\textcolor{input}{int *NormalIndex}, \textcolor{input}{int NormalListFlag}, & \\
~~~~~~\textcolor{input}{DataType\_t NormalDataType}, \textcolor{input}{void *NormalList}); & \\
\textcolor{output}{\textit{ier}} = cg\_boco\_gridlocation\_write(\textcolor{input}{int fn}, \textcolor{input}{int B}, \textcolor{input}{int Z}, \textcolor{input}{int BC}, & - w m \\
~~~~~~\textcolor{input}{\textit{GridLocation\_t location}}); & \\
\textcolor{output}{\textit{ier}} = cg\_nbocos(\textcolor{input}{int fn}, \textcolor{input}{int B}, \textcolor{input}{int Z}, \textcolor{output}{\textit{int *nbocos}}); & r - m \\
\textcolor{output}{\textit{ier}} = cg\_boco\_info(\textcolor{input}{int fn}, \textcolor{input}{int B}, \textcolor{input}{int Z}, \textcolor{input}{int BC}, \textcolor{output}{\textit{char *boconame}}, & r - m \\
~~~~~~\textcolor{output}{\textit{BCType\_t *bocotype}}, \textcolor{output}{\textit{PointSetType\_t *ptset\_type}}, \textcolor{output}{\textit{cgsize\_t *npnts}}, & \\
~~~~~~\textcolor{output}{\textit{int *NormalIndex}}, \textcolor{output}{\textit{cgsize\_t *NormalListFlag}}, & \\
~~~~~~\textcolor{output}{\textit{DataType\_t *NormalDataType}}, \textcolor{output}{\textit{int *ndataset}}); & \\
\textcolor{output}{\textit{ier}} = cg\_boco\_read(\textcolor{input}{int fn}, \textcolor{input}{int B}, \textcolor{input}{int Z}, \textcolor{input}{int BC}, \textcolor{output}{\textit{cgsize\_t *pnts}}, & r - m \\
~~~~~~\textcolor{output}{\textit{void *NormalList}}); & \\
\textcolor{output}{\textit{ier}} = cg\_boco\_gridlocation\_read(\textcolor{input}{int fn}, \textcolor{input}{int B}, \textcolor{input}{int Z}, \textcolor{input}{int BC}, & r - m \\
~~~~~~\textcolor{output}{\textit{GridLocation\_t *location}}); & \\
\hline
call cg\_boco\_write\_f(\textcolor{input}{fn}, \textcolor{input}{B}, \textcolor{input}{Z}, \textcolor{input}{boconame}, \textcolor{input}{bocotype}, \textcolor{input}{ptset\_type}, \textcolor{input}{npnts}, & - w m \\
~~~~~\textcolor{input}{pnts}, \textcolor{output}{\textit{BC}}, \textcolor{output}{\textit{ier}}) & \\
call cg\_boco\_normal\_write\_f(\textcolor{input}{fn}, \textcolor{input}{B}, \textcolor{input}{Z}, \textcolor{input}{BC}, \textcolor{input}{NormalIndex}, \textcolor{input}{NormalListFlag}, & - w m \\
~~~~~\textcolor{input}{NormalDataType}, \textcolor{input}{NormalList}, \textcolor{output}{\textit{ier}}) & \\
call cg\_boco\_gridlocation\_write\_f(\textcolor{input}{fn}, \textcolor{input}{B}, \textcolor{input}{Z}, \textcolor{input}{BC}, \textcolor{input}{location}, \textcolor{output}{\textit{ier}}); & - w m \\
call cg\_nbocos\_f(\textcolor{input}{fn}, \textcolor{input}{B}, \textcolor{input}{Z}, \textcolor{output}{\textit{nbocos}}, \textcolor{output}{\textit{ier}}) & r - m \\
call cg\_boco\_info\_f(\textcolor{input}{fn}, \textcolor{input}{B}, \textcolor{input}{Z}, \textcolor{input}{BC}, \textcolor{output}{\textit{boconame}}, \textcolor{output}{\textit{bocotype}}, \textcolor{output}{\textit{ptset\_type}}, & r - m \\
~~~~~\textcolor{output}{\textit{npnts}}, \textcolor{output}{\textit{NormalIndex}}, \textcolor{output}{\textit{NormalListFlag}}, \textcolor{output}{\textit{NormalDataType}}, \textcolor{output}{\textit{ndataset}}, & \\
~~~~~\textcolor{output}{\textit{ier}}) & \\
call cg\_boco\_read\_f(\textcolor{input}{fn}, \textcolor{input}{B}, \textcolor{input}{Z}, \textcolor{input}{BC}, \textcolor{output}{\textit{pnts}}, \textcolor{output}{\textit{NormalList}}, \textcolor{output}{\textit{ier}}) & r - m \\
call cg\_boco\_gridlocation\_read\_f(\textcolor{input}{fn}, \textcolor{input}{B}, \textcolor{input}{Z}, \textcolor{input}{BC}, \textcolor{output}{\textit{location}}, \textcolor{output}{\textit{ier}}); & r - m \\
\end{fctbox}

\noindent
\textbf{\textcolor{input}{Input}/\textcolor{output}{\textit{Output}}}

\begin{Ventryi}{\texttt{NormalListFlag}}\raggedright
\item [\texttt{fn}]
      CGNS file index number.
      (\textcolor{input}{Input})
\item [\texttt{B}]
      Base index number, where $1 \leq \text{\texttt{B}} \leq \text{\texttt{nbases}}$.
      (\textcolor{input}{Input})
\item [\texttt{Z}]
      Zone index number, where $1 \leq \text{\texttt{Z}} \leq \text{\texttt{nzones}}$.
      (\textcolor{input}{Input})
\item [\texttt{BC}]
      Boundary condition index number, where $1 \leq \text{\texttt{BC}} \leq \text{\texttt{nbocos}}$.
      (\textcolor{input}{Input} for \texttt{cg\_boco\_normal\_write},
      \texttt{cg\_boco\_info}, \texttt{cg\_boco\_read};
      \textcolor{output}{\textit{output}} for \texttt{cg\_boco\_write})
\item [\texttt{nbocos}]
      Number of boundary conditions in zone \texttt{Z}.
      (\textcolor{output}{\textit{Output}})
\item [\texttt{boconame}]
      Name of the boundary condition.
      (\textcolor{input}{Input} for \texttt{cg\_boco\_write};
      \textcolor{output}{\textit{output}} for \texttt{cg\_boco\_info})
\item [\texttt{bocotype}]
      Type of boundary condition defined.
      See the eligible types for \texttt{BCType\_t} in \autoref{s:typedefs}.
      Note that if \texttt{bocotype} is \texttt{FamilySpecified}
      the boundary condition type is being specified for the family
      to which the boundary belongs.
      The boundary condition type for the family may be read and written
      using \texttt{cg\_fambc\_read} and \texttt{cg\_fambc\_write},
      as described in \autoref{s:familybc}.
      (\textcolor{input}{Input} for \texttt{cg\_boco\_write};
      \textcolor{output}{\textit{output}} for \texttt{cg\_boco\_info})
\item [\texttt{ptset\_type}]
      The extent of the boundary condition may be defined using a range
      of points or elements using \texttt{PointRange}, or using a
      discrete list of all points or elements at which the boundary
      condition is applied using \texttt{PointList}.

      When the boundary condition is to be applied anywhere other than points,
      then \texttt{GridLocation\_t} under the \texttt{BC\_t} node must
      be used to indicate this.
      The value of \texttt{GridLocation\_t} may be read or written by
      \texttt{cg\_boco\_gridlocation\_read} and
      \texttt{cg\_boco\_gridlocation\_write}.
      As in previous versions of the library, this may also be done by
      first using \texttt{cg\_goto} (\autoref{s:navigating})
      to access the \texttt{BC\_t} node, then using
      \texttt{cg\_gridlocation\_read} or
      \texttt{cg\_gridlocation\_write} (\autoref{s:gridlocation}).
      (\textcolor{input}{Input} for \texttt{cg\_boco\_write};
      \textcolor{output}{\textit{output}} for \texttt{cg\_boco\_info})
\item [\texttt{npnts}]
      Number of points or elements defining the boundary
      condition region.
      For a \texttt{ptset\_type} of \texttt{PointRange},
      \texttt{npnts} is always two.
      For a \texttt{ptset\_type} of \texttt{PointList},
      \texttt{npnts} is the number of points or elements in the list.
      (\textcolor{input}{Input} for \texttt{cg\_boco\_write};
      \textcolor{output}{\textit{output}} for \texttt{cg\_boco\_info})
\item [\texttt{pnts}]
      Array of point or element indices defining the boundary condition
      region.
      There should be \texttt{npnts} values, each of dimension
      \texttt{IndexDimension} (i.e., 1 for unstructured grids,
      and 2 or 3 for structured grids with 2-D or 3-D elements,
      respectively).
      (\textcolor{input}{Input} for \texttt{cg\_boco\_write};
      \textcolor{output}{\textit{output}} for \texttt{cg\_boco\_read})
\item [\texttt{NormalIndex}]
      Index vector indicating the computational coordinate direction
      of the boundary condition patch normal.
      (\textcolor{input}{Input} for \texttt{cg\_boco\_normal\_write};
      \textcolor{output}{\textit{output}} for \texttt{cg\_boco\_info})
\item [\texttt{NormalListFlag}]
      For \texttt{cg\_boco\_normal\_write}, \texttt{NormalListFlag} is a
      flag indicating if the normals are defined in \texttt{NormalList};
      1 if they are defined, 0 if they're not.

      For \texttt{cg\_boco\_info}, if the normals are defined in
      \texttt{NormalList}, \texttt{NormalListFlag} is the number of points
      in the patch times \texttt{phys\_dim}, the number of coordinates
      required to define a vector in the field.
      If the normals are not defined in \texttt{NormalList},
      \texttt{NormalListFlag} is 0.
      (\textcolor{input}{Input} for \texttt{cg\_boco\_normal\_write};
      \textcolor{output}{\textit{output}} for \texttt{cg\_boco\_info})
\item [\texttt{NormalDataType}]
      Data type used in the definition of the normals.
      Admissible data types for the normals are \texttt{RealSingle} and
      \texttt{RealDouble}.
      (\textcolor{input}{Input} for \texttt{cg\_boco\_normal\_write};
      \textcolor{output}{\textit{output}} for \texttt{cg\_boco\_info})
\item [\texttt{NormalList}]
      List of vectors normal to the boundary condition patch pointing
      into the interior of the zone.
      (\textcolor{input}{Input} for \texttt{cg\_boco\_normal\_write};
      \textcolor{output}{\textit{output}} for \texttt{cg\_boco\_read})
\item [\texttt{ndataset}]
      Number of boundary condition datasets for the current boundary
      condition.
      (\textcolor{output}{\textit{Output}})
\item [\texttt{location}]
      Grid location used in the definition of the point set.
      The currently admissible locations are \texttt{Vertex} (the default
      if not given). For 2-D grids, \texttt{EdgeCenter} is also allowed,
      and for 3-D grids, the additional values of \texttt{FaceCenter},
      \texttt{IFaceCenter}, \texttt{JFaceCenter}, and
      \texttt{KFaceCenter} may be used.
      (\textcolor{input}{Input} for \texttt{cg\_boco\_gridlocation\_write};
      \textcolor{output}{\textit{output}} for
      \texttt{cg\_boco\_gridlocation\_read})
\item [\texttt{ier}]
      Error status.
      (\textcolor{output}{\textit{Output}})
\end{Ventryi}

\noindent
\textit{Notes}: (see CPEX 0031)
\begin{itemize}
\item The use of \texttt{ElementList} and \texttt{ElementRange} for
      \texttt{ptset\_type} is deprecated and should not be used
      in new code. These are still currently accepted, but will be internally
      replaced with the appropriate values of \texttt{PointList/PointRange}
      and \texttt{GridLocation\_t}.
\item \texttt{CellCenter} for \texttt{GridLocation\_t} is also deprecated.
      If used, the value will be replaced by \texttt{EdgeCenter} for
      2-D grids or \texttt{FaceCenter} for 3-D grids.
\end{itemize}

\subsection{Boundary Condition Datasets}
\label{s:bcdataset}

\noindent
\textit{Node}: \texttt{BCDataSet\_t}

\begin{fctbox}
\textcolor{output}{\textit{ier}} = cg\_dataset\_write(\textcolor{input}{int fn}, \textcolor{input}{int B}, \textcolor{input}{int Z}, \textcolor{input}{int BC}, \textcolor{input}{char *DatasetName}, & - w m \\
~~~~~~\textcolor{input}{BCType\_t BCType}, \textcolor{output}{\textit{int *Dset}}); & \\
\textcolor{output}{\textit{ier}} = cg\_dataset\_read(\textcolor{input}{int fn}, \textcolor{input}{int B}, \textcolor{input}{int Z}, \textcolor{input}{int BC}, \textcolor{input}{int Dset}, & r - m \\
~~~~~~\textcolor{output}{\textit{char *DatasetName}}, \textcolor{output}{\textit{BCType\_t *BCType}}, \textcolor{output}{\textit{int *DirichletFlag}}, & \\
~~~~~~\textcolor{output}{\textit{int *NeumannFlag}}); & \\
\hline
call cg\_dataset\_write\_f(\textcolor{input}{fn}, \textcolor{input}{B}, \textcolor{input}{Z}, \textcolor{input}{BC}, \textcolor{input}{DatasetName}, \textcolor{input}{BCType}, \textcolor{output}{\textit{Dset}}, \textcolor{output}{\textit{ier}}) & - w m \\
call cg\_dataset\_read\_f(\textcolor{input}{fn}, \textcolor{input}{B}, \textcolor{input}{Z}, \textcolor{input}{BC}, \textcolor{input}{Dset}, \textcolor{output}{\textit{DatasetName}}, \textcolor{output}{\textit{BCType}}, & r - m \\
~~~~~~\textcolor{output}{\textit{DirichletFlag}}, \textcolor{output}{\textit{NeumannFlag}}, \textcolor{output}{\textit{ier}}) & \\
\end{fctbox}

\noindent
\textbf{\textcolor{input}{Input}/\textcolor{output}{\textit{Output}}}

\begin{Ventryi}{\texttt{DirichletFlag}}\raggedright
\item [\texttt{fn}]
      CGNS file index number.
      (\textcolor{input}{Input})
\item [\texttt{B}]
      Base index number, where $1 \leq \text{\texttt{B}} \leq \text{\texttt{nbases}}$.
      (\textcolor{input}{Input})
\item [\texttt{Z}]
      Zone index number, where $1 \leq \text{\texttt{Z}} \leq \text{\texttt{nzones}}$.
      (\textcolor{input}{Input})
\item [\texttt{BC}]
      Boundary condition index number, where $1 \leq \text{\texttt{BC}} \leq \text{\texttt{nbocos}}$.
      (\textcolor{input}{Input})
\item [\texttt{Dset}]
      Dataset index number, where $1 \leq \text{\texttt{Dset}} \leq \text{\texttt{ndataset}}$.
      (\textcolor{input}{Input} for \texttt{cg\_dataset\_read};
      \textcolor{output}{\textit{output}} for \texttt{cg\_dataset\_write})
\item [\texttt{DatasetName}]
      Name of dataset.
      (\textcolor{input}{Input} for \texttt{cg\_dataset\_write};
      \textcolor{output}{\textit{output}} for \texttt{cg\_dataset\_read})
\item [\texttt{BCType}]
      Simple boundary condition type for the dataset.
      The supported types are listed in the table of ``Simple
      Boundary Condition Types'' in the SIDS manual, but note that
      \texttt{FamilySpecified} does not apply here.
      (\textcolor{input}{Input} for \texttt{cg\_dataset\_write};
      \textcolor{output}{\textit{output}} for \texttt{cg\_dataset\_read})
\item [\texttt{DirichletFlag}]
      Flag indicating if the dataset contains Dirichlet data.
      (\textcolor{output}{\textit{Output}})
\item [\texttt{NeumannFlag}]
      Flag indicating if the dataset contains Neumann data.
      (\textcolor{output}{\textit{Output}})
\item [\texttt{ier}]
      Error status.
      (\textcolor{output}{\textit{Output}})
\end{Ventryi}

The above functions are applicable to \texttt{BCDataSet\_t} nodes
that are children of \texttt{BC\_t} nodes.

For \texttt{BCDataSet\_t} nodes that are children of a \texttt{BC\_t}
node, after accessing a particular \texttt{BCDataSet\_t} node using
\texttt{cg\_goto}, the Point Set functions described in \autoref{s:ptset}
may be used to read or write the locations at which the boundary
conditions are to be applied.
This is only applicable when the boundary conditions are to be applied
at locations different from those used with \texttt{cg\_boco\_write} to
define the boundary condition region (e.g., when the region is being
defined by specification of vertices, but the boundary conditions are to
be applied at face centers).

When writing point set data to a \texttt{BCDataSet\_t} node, in addition
to the specification of the indices using \texttt{cg\_ptset\_write},
the function \texttt{cg\_gridlocation\_write} must also be used to
specify the location of the data with respect to the grid (e.g.,
\texttt{Vertex} or \texttt{FaceCenter}).

\begin{fctbox}
\textcolor{output}{\textit{ier}} = cg\_bcdataset\_write(\textcolor{input}{char *DatasetName}, \textcolor{input}{BCType\_t BCType}, & - w m \\
~~~~~~\textcolor{input}{BCDataType\_t BCDataType}); & \\
\textcolor{output}{\textit{ier}} = cg\_bcdataset\_info(\textcolor{output}{\textit{int *ndataset}}); & \\
\textcolor{output}{\textit{ier}} = cg\_bcdataset\_read(\textcolor{input}{int Dset}, \textcolor{output}{\textit{char *DatasetName}}, \textcolor{output}{\textit{BCType\_t *BCType}}, & r - m \\
~~~~~~\textcolor{output}{\textit{int *DirichletFlag}}, \textcolor{output}{\textit{int *NeumannFlag}}); & \\
\hline
call cg\_bcdataset\_write\_f(\textcolor{input}{DatasetName}, \textcolor{input}{BCType}, \textcolor{input}{BCDataType\_t BCDataType}, \textcolor{output}{\textit{ier}}) & - w m \\
call cg\_bcdataset\_info\_f(\textcolor{output}{\textit{int *ndataset}}, \textcolor{output}{\textit{ier}}) & \\
call cg\_bcdataset\_read\_f(\textcolor{input}{Dset}, \textcolor{output}{\textit{DatasetName}}, \textcolor{output}{\textit{BCType}}, \textcolor{output}{\textit{DirichletFlag}}, & r - m \\
~~~~~~\textcolor{output}{\textit{NeumannFlag}}, \textcolor{output}{\textit{ier}}) & \\
\end{fctbox}

\noindent
\textbf{\textcolor{input}{Input}/\textcolor{output}{\textit{Output}}}

\begin{Ventryi}{\texttt{DirichletFlag}}\raggedright
\item [\texttt{Dset}]
      Dataset index number, where $1 \leq \text{\texttt{Dset}} \leq \text{\texttt{ndataset}}$.
      (\textcolor{input}{Input})
\item [\texttt{DatasetName}]
      Name of dataset.
      (\textcolor{input}{Input} for \texttt{cg\_bcdataset\_write};
      \textcolor{output}{\textit{output}} for \texttt{cg\_bcdataset\_read})
\item [\texttt{BCType}]
      Simple boundary condition type for the dataset.
      The supported types are listed in the table of ``Simple
      Boundary Condition Types'' in the SIDS manual, but note that
      \texttt{FamilySpecified} does not apply here.
      (\textcolor{input}{Input} for \texttt{cg\_bcdataset\_write};
      \textcolor{output}{\textit{output}} for \texttt{cg\_bcdataset\_read})
\item [\texttt{BCDataType}]
      Type of boundary condition in the dataset (i.e., for a
      \texttt{BCData\_t} child node).
      Admissible types are \texttt{Dirichlet} and \texttt{Neumann}.
      (\textcolor{input}{Input})
\item [\texttt{ndataset}]
      Number of \texttt{BCDataSet} nodes under the current
      \texttt{FamilyBC\_t} node.
      (\textcolor{output}{\textit{Output}})
\item [\texttt{DirichletFlag}]
      Flag indicating if the dataset contains Dirichlet data.
      (\textcolor{output}{\textit{Output}})
\item [\texttt{NeumannFlag}]
      Flag indicating if the dataset contains Neumann data.
      (\textcolor{output}{\textit{Output}})
\item [\texttt{ier}]
      Error status.
      (\textcolor{output}{\textit{Output}})
\end{Ventryi}

The above functions are applicable to \texttt{BCDataSet\_t} nodes
that are used to define boundary conditions for a CFD family, and thus
are children of a \texttt{FamilyBC\_t} node.
The \texttt{FamilyBC\_t} node must first be accessed using \texttt{cg\_goto}.

The first time \texttt{cg\_bcdataset\_write} is called with a particular
\texttt{DatasetName}, \texttt{BCType}, and \texttt{BCDataType}, a new
\texttt{BCDataSet\_t} node is created, with a child \texttt{BCData\_t} node.
Subsequent calls with the same \texttt{DatasetName} and \texttt{BCType}
may be made to add additional \texttt{BCData\_t} nodes, of type
\texttt{BCDataType}, to the existing \texttt{BCDataSet\_t} node.

\subsection{Boundary Condition Data}
\label{s:bcdata}

\noindent
\textit{Node}: \texttt{BCData\_t}

\begin{fctbox}
\textcolor{output}{\textit{ier}} = cg\_bcdata\_write(\textcolor{input}{int fn}, \textcolor{input}{int B}, \textcolor{input}{int Z}, \textcolor{input}{int BC}, \textcolor{input}{int Dset}, & - w m \\
~~~~~~\textcolor{input}{BCDataType\_t BCDataType}); & \\
\hline
call cg\_bcdata\_write\_f(\textcolor{input}{fn}, \textcolor{input}{B}, \textcolor{input}{Z}, \textcolor{input}{BC}, \textcolor{input}{Dset}, \textcolor{input}{BCDataType}, \textcolor{output}{\textit{ier}}) & - w m \\
\end{fctbox}

\noindent
\textbf{\textcolor{input}{Input}/\textcolor{output}{\textit{Output}}}

\begin{Ventryi}{\texttt{ptset\_type}}\raggedright
\item [\texttt{fn}]
      CGNS file index number.
      (\textcolor{input}{Input})
\item [\texttt{B}]
      Base index number, where $1 \leq \text{\texttt{B}} \leq \text{\texttt{nbases}}$.
      (\textcolor{input}{Input})
\item [\texttt{Z}]
      Zone index number, where $1 \leq \text{\texttt{Z}} \leq \text{\texttt{nzones}}$.
      (\textcolor{input}{Input})
\item [\texttt{BC}]
      Boundary condition index number, where $1 \leq \text{\texttt{BC}} \leq \text{\texttt{nbocos}}$.
      (\textcolor{input}{Input})
\item [\texttt{Dset}]
      Dataset index number, where $1 \leq \text{\texttt{Dset}} \leq \text{\texttt{ndataset}}$.
      (\textcolor{input}{Input})
\item [\texttt{BCDataType}]
      Type of boundary condition in the dataset.
      Admissible boundary condition types are \texttt{Dirichlet} and
      \texttt{Neumann}.
      (\textcolor{input}{Input})
\item [\texttt{ier}]
      Error status.
      (\textcolor{output}{\textit{Output}})
\end{Ventryi}

To write the boundary condition data itself, after creating the
\texttt{BCData\_t} node using the function \texttt{cg\_bcdata\_write}, use
\texttt{cg\_goto} to access the node, then \texttt{cg\_array\_write} to
write the data.
Note that when using \texttt{cg\_goto} to access a \texttt{BCData\_t}
node, the node index should be specified as either \texttt{Dirichlet} or
\texttt{Neumann}, depending on the type of boundary condition.
See the description of \texttt{cg\_goto} in \autoref{s:navigating} for
details.

\newpage
\subsection{Special Boundary Condition Properties}
\label{s:bcproperty}

\noindent
\textit{Node}: \texttt{BCProperty\_t}

\begin{fctbox}
\textcolor{output}{\textit{ier}} = cg\_bc\_wallfunction\_write(\textcolor{input}{int fn}, \textcolor{input}{int B}, \textcolor{input}{int Z}, \textcolor{input}{int BC}, & - w m \\
~~~~~~\textcolor{input}{WallFunctionType\_t WallFunctionType}); & \\
\textcolor{output}{\textit{ier}} = cg\_bc\_area\_write(\textcolor{input}{int fn}, \textcolor{input}{int B}, \textcolor{input}{int Z}, \textcolor{input}{int BC}, & - w m \\
~~~~~~\textcolor{input}{AreaType\_t AreaType}, \textcolor{input}{float SurfaceArea}, \textcolor{input}{\textit{char *RegionName}}); & \\
\textcolor{output}{\textit{ier}} = cg\_bc\_wallfunction\_read(\textcolor{input}{int fn}, \textcolor{input}{int B}, \textcolor{input}{int Z}, \textcolor{input}{int BC}, & r - m \\
~~~~~~\textcolor{output}{\textit{WallFunctionType\_t *WallFunctionType}}); & \\
\textcolor{output}{\textit{ier}} = cg\_bc\_area\_read(\textcolor{input}{int fn}, \textcolor{input}{int B}, \textcolor{input}{int Z}, \textcolor{input}{int BC}, & r - m \\
~~~~~~\textcolor{output}{\textit{AreaType\_t *AreaType}}, \textcolor{output}{\textit{float *SurfaceArea}}, \textcolor{output}{\textit{char *RegionName}}); & \\
\hline
call cg\_bc\_wallfunction\_write\_f(\textcolor{input}{fn}, \textcolor{input}{B}, \textcolor{input}{Z}, \textcolor{input}{BC}, \textcolor{input}{WallFunctionType}, \textcolor{output}{\textit{ier}}) & - w m \\
call cg\_bc\_area\_write\_f(\textcolor{input}{fn}, \textcolor{input}{B}, \textcolor{input}{Z}, \textcolor{input}{BC}, \textcolor{input}{AreaType}, \textcolor{input}{SurfaceArea}, & - w m \\
~~~~~\textcolor{input}{RegionName}, \textcolor{output}{\textit{ier}}) & \\
call cg\_bc\_wallfunction\_read\_f(\textcolor{input}{fn}, \textcolor{input}{B}, \textcolor{input}{Z}, \textcolor{input}{BC}, \textcolor{output}{\textit{WallFunctionType}}, \textcolor{output}{\textit{ier}}) & r - m \\
call cg\_bc\_area\_read\_f(\textcolor{input}{fn}, \textcolor{input}{B}, \textcolor{input}{Z}, \textcolor{input}{BC}, \textcolor{output}{\textit{AreaType}}, \textcolor{output}{\textit{SurfaceArea}}, & r - m \\
~~~~~\textcolor{output}{\textit{RegionName}}, \textcolor{output}{\textit{ier}}) & \\
\end{fctbox}

\noindent
\textbf{\textcolor{input}{Input}/\textcolor{output}{\textit{Output}}}

\begin{Ventryi}{\texttt{WallFunctionType}}\raggedright
\item [\texttt{fn}]
      CGNS file index number.
      (\textcolor{input}{Input})
\item [\texttt{B}]
      Base index number, where $1 \leq \text{\texttt{B}} \leq \text{\texttt{nbases}}$.
      (\textcolor{input}{Input})
\item [\texttt{Z}]
      Zone index number, where $1 \leq \text{\texttt{Z}} \leq \text{\texttt{nzones}}$.
      (\textcolor{input}{Input})
\item [\texttt{BC}]
      Boundary condition index number, where $1 \leq \text{\texttt{BC}} \leq \text{\texttt{nbocos}}$.
      (\textcolor{input}{Input})
\item [\texttt{WallFunctionType}]
      The wall function type.
      Valid types are \texttt{CG\_Null}, \texttt{CG\_UserDefined}, and \texttt{Generic}.
      (\textcolor{input}{Input} for \texttt{cg\_bc\_wallfunction\_write};
      \textcolor{output}{\textit{output}} for
      \texttt{cg\_bc\_wallfunction\_read})
\item [\texttt{AreaType}]
      The type of area.
      Valid types are \texttt{CG\_Null}, \texttt{CG\_UserDefined}, \texttt{BleedArea},
      and \texttt{CaptureArea}.
      (\textcolor{input}{Input} for \texttt{cg\_bc\_area\_write};
      \textcolor{output}{\textit{output}} for \texttt{cg\_bc\_area\_read})
\item [\texttt{SurfaceArea}]
      The size of the area. (In Fortran, this is a Real*4 value.)
      (\textcolor{input}{Input} for \texttt{cg\_bc\_area\_write};
      \textcolor{output}{\textit{output}} for \texttt{cg\_bc\_area\_read})
\item [\texttt{RegionName}]
      The name of the region, 32 characters max.
      (\textcolor{input}{Input} for \texttt{cg\_bc\_area\_write};
      \textcolor{output}{\textit{output}} for \texttt{cg\_bc\_area\_read})
\item [\texttt{ier}]
      Error status.
      (\textcolor{output}{\textit{Output}})
\end{Ventryi}

The ``\texttt{write}'' functions will create the \texttt{BCProperty\_t}
node if it doesn't already exist, then add the appropriate boundary
condition property.
Multiple boundary condition properties may be recorded under the same
\texttt{BCProperty\_t} node.

The ``\texttt{read}'' functions will return with $\text{\texttt{ier}} = 2 =
\text{\texttt{CG\_NODE\_NOT\_FOUND}}$ if the requested boundary condition
property, or the \texttt{BCProperty\_t} node itself, doesn't exist.


\clearemptydoublepage
\section{Equation Specification}
\label{s:equation}
\thispagestyle{plain}

\subsection{Flow Equation Set}
\label{s:flowequationset}

\noindent
\textit{Node}: \texttt{FlowEquationSet\_t}

\begin{fctbox}
\textcolor{output}{\textit{ier}} = cg\_equationset\_write(\textcolor{input}{int EquationDimension}); & - w m \\
\textcolor{output}{\textit{ier}} = cg\_equationset\_read(\textcolor{output}{\textit{int *EquationDimension}}, & r - m \\
~~~~~~\textcolor{output}{\textit{int *GoverningEquationsFlag}}, \textcolor{output}{\textit{int *GasModelFlag}}, & \\
~~~~~~\textcolor{output}{\textit{int *ViscosityModelFlag}}, \textcolor{output}{\textit{int *ThermalConductModelFlag}}, & \\
~~~~~~\textcolor{output}{\textit{int *TurbulenceClosureFlag}}, \textcolor{output}{\textit{int *TurbulenceModelFlag}}); & \\
\textcolor{output}{\textit{ier}} = cg\_equationset\_chemistry\_read(\textcolor{output}{\textit{int *ThermalRelaxationFlag}}, & r - m \\
~~~~~~\textcolor{output}{\textit{int *ChemicalKineticsFlag}}); & \\
\textcolor{output}{\textit{ier}} = cg\_equationset\_elecmagn\_read(\textcolor{output}{\textit{int *ElecFldModelFlag}}, & r - m \\
~~~~~~\textcolor{output}{\textit{int *MagnFldModelFlag}}, \textcolor{output}{\textit{ConductivityModelFlag}}); & \\
\hline
call cg\_equationset\_write\_f(\textcolor{input}{EquationDimension}, \textcolor{output}{\textit{ier}}) & - w m \\
call cg\_equationset\_read\_f(\textcolor{output}{\textit{EquationDimension}}, & r - m \\
~~~~~\textcolor{output}{\textit{GoverningEquationsFlag}}, \textcolor{output}{\textit{GasModelFlag}}, \textcolor{output}{\textit{ViscosityModelFlag}}, & \\
~~~~~\textcolor{output}{\textit{ThermalConductModelFlag}}, \textcolor{output}{\textit{TurbulenceClosureFlag}}, & \\
~~~~~\textcolor{output}{\textit{TurbulenceModelFlag}}, \textcolor{output}{\textit{ier}}) & \\
call cg\_equationset\_chemistry\_read\_f(\textcolor{output}{\textit{ThermalRelaxationFlag}}, & r - m \\
~~~~~\textcolor{output}{\textit{ChemicalKineticsFlag}}, \textcolor{output}{\textit{ier}}) & \\
call cg\_equationset\_elecmagn\_read\_f(\textcolor{output}{\textit{ElecFldModelFlag}}, & r - m \\
~~~~~\textcolor{output}{\textit{MagnFldModelFlag}}, \textcolor{output}{\textit{ConductivityModelFlag}}, \textcolor{output}{\textit{ier}}) & \\
\end{fctbox}

\noindent
\textbf{\textcolor{input}{Input}/\textcolor{output}{\textit{Output}}}

\noindent (Note that for Fortran calls, all integer arguments are integer*4 in 32-bit mode and integer*8 in 64-bit mode.
See ``64-bit Fortran Portability and Issues" section.)

\begin{Ventryi}{\texttt{ThermalConductModelFlag}}\raggedright
\item [\texttt{EquationDimension}]
      Dimensionality of the governing equations; it is the number of
      spatial variables describing the flow.
      (\textcolor{input}{Input} for \texttt{cg\_equationset\_write};
      \textcolor{output}{\textit{output}} for \texttt{cg\_equationset\_info})
\item [\texttt{GoverningEquationsFlag}]
      Flag indicating whether or not this \texttt{FlowEquationSet\_t}
      node includes the definition of the governing equations;
      0 if it doesn't, 1 if it does.
      (\textcolor{output}{\textit{Output}})
\item [\texttt{GasModelFlag}]
      Flag indicating whether or not this \texttt{FlowEquationSet\_t}
      node includes the definition of a gas model;
      0 if it doesn't, 1 if it does.
      (\textcolor{output}{\textit{Output}})
\item [\texttt{ViscosityModelFlag}]
      Flag indicating whether or not this \texttt{FlowEquationSet\_t}
      node includes the definition of a viscosity model;
      0 if it doesn't, 1 if it does.
      (\textcolor{output}{\textit{Output}})
\item [\texttt{ThermalConductModelFlag}]
      Flag indicating whether or not this \texttt{FlowEquationSet\_t}
      node includes the definition of a thermal conductivity model;
      0 if it doesn't, 1 if it does.
      (\textcolor{output}{\textit{Output}})
\item [\texttt{TurbulenceClosureFlag}]
      Flag indicating whether or not this \texttt{FlowEquationSet\_t}
      node includes the definition of the turbulence closure;
      0 if it doesn't, 1 if it does.
      (\textcolor{output}{\textit{Output}})
\item [\texttt{TurbulenceModelFlag}]
      Flag indicating whether or not this \texttt{FlowEquationSet\_t}
      node includes the definition of a turbulence model;
      0 if it doesn't, 1 if it does.
      (\textcolor{output}{\textit{Output}})
\item [\texttt{ThermalRelaxationFlag}]
      Flag indicating whether or not this \texttt{FlowEquationSet\_t}
      node includes the definition of thae thermal relaxation model;
      0 if it doesn't, 1 if it does.
      (\textcolor{output}{\textit{Output}})
\item [\texttt{ChemicalKineticsFlag}]
      Flag indicating whether or not this \texttt{FlowEquationSet\_t}
      node includes the definition of a chemical kinetics model;
      0 if it doesn't, 1 if it does.
      (\textcolor{output}{\textit{Output}})
\item [\texttt{ElecFldModelFlag}]
      Flag indicating whether or not this \texttt{FlowEquationSet\_t}
      node includes the definition of an electric field model for
      electromagnetic flows;
      0 if it doesn't, 1 if it does.
      (\textcolor{output}{\textit{Output}})
\item [\texttt{MagnFldModelFlag}]
      Flag indicating whether or not this \texttt{FlowEquationSet\_t}
      node includes the definition of a magnetic field model for
      electromagnetic flows;
      0 if it doesn't, 1 if it does.
      (\textcolor{output}{\textit{Output}})
\item [\texttt{ConductivityModelFlag}]
      Flag indicating whether or not this \texttt{FlowEquationSet\_t}
      node includes the definition of a conductivity model for
      electromagnetic flows;
      0 if it doesn't, 1 if it does.
      (\textcolor{output}{\textit{Output}})
\item [\texttt{ier}]
      Error status.
      (\textcolor{output}{\textit{Output}})
\end{Ventryi}

\subsection{Governing Equations}
\label{s:governingequations}

\noindent
\textit{Node}: \texttt{GoverningEquations\_t}

\begin{fctbox}
\textcolor{output}{\textit{ier}} = cg\_governing\_write(\textcolor{input}{GoverningEquationsType\_t Equationstype}); & - w m \\
\textcolor{output}{\textit{ier}} = cg\_governing\_read(\textcolor{output}{\textit{GoverningEquationsType\_t *EquationsType}}); & r - m \\
\textcolor{output}{\textit{ier}} = cg\_diffusion\_write(\textcolor{input}{int *diffusion\_model}); & - w m \\
\textcolor{output}{\textit{ier}} = cg\_diffusion\_read(\textcolor{output}{\textit{int *diffusion\_model}}); & r - m \\
\hline
call cg\_governing\_write\_f(\textcolor{input}{EquationsType}, \textcolor{output}{\textit{ier}}) & - w m \\
call cg\_governing\_read\_f(\textcolor{output}{\textit{EquationsType}}, \textcolor{output}{\textit{ier}}) & r - m \\
call cg\_diffusion\_write\_f(\textcolor{input}{diffusion\_model}, \textcolor{output}{\textit{ier}}) & - w m \\
call cg\_diffusion\_read\_f(\textcolor{output}{\textit{diffusion\_model}}, \textcolor{output}{\textit{ier}}) & r - m \\
\end{fctbox}

\noindent
\textbf{\textcolor{input}{Input}/\textcolor{output}{\textit{Output}}}

\noindent (Note that for Fortran calls, all integer arguments are integer*4 in 32-bit mode and integer*8 in 64-bit mode.
See ``64-bit Fortran Portability and Issues" section.)

\begin{Ventryi}{\texttt{diffusion\_model}}\raggedright
\item [\texttt{EquationsType}]
      Type of governing equations.
      The admissible types are \texttt{CG\_Null}, \texttt{CG\_UserDefined},
      \texttt{FullPotential}, \texttt{Euler}, \texttt{NSLaminar},
      \texttt{NSTurbulent}, \texttt{NSLaminarIncompressible}, and
      \texttt{NSTurbulentIncompressible}.
      (\textcolor{input}{Input} for \texttt{cg\_governing\_write};
      \textcolor{output}{\textit{output}} for \texttt{cg\_governing\_read})
\item [\texttt{diffusion\_model}]
      Flags defining which diffusion terms are included in the governing
      equations.
      This is only applicable to the Navier-Stokes equations with
      structured grids.
      See the discussion of \texttt{GoverningEquations\_t} in the
      \href{../sids/sids.pdf}{SIDS manual} for details.
      (\textcolor{input}{Input} for \texttt{cg\_diffusion\_write};
      \textcolor{output}{\textit{output}} for \texttt{cg\_diffusion\_read})
\item [\texttt{ier}]
      Error status.
      (\textcolor{output}{\textit{Output}})
\end{Ventryi}

\subsection{Auxiliary Models}
\label{s:auxiliarymodels}

\noindent
\textit{Nodes}: \texttt{GasModel\_t}, 
                \texttt{ViscosityModel\_t},
                \texttt{ThermalConductivityModel\_t},
                \texttt{TurbulenceClosure\_t},
                \texttt{TurbulenceModel\_t},
                \texttt{ThermalRelaxationModel\_t},
                \texttt{ChemicalKineticsModel\_t},
                \texttt{EMElectricFieldModel\_t},
                \texttt{EMMagneticFieldModel\_t},
                \texttt{EMConductivityModel\_t}

\begin{fctbox}
\textcolor{output}{\textit{ier}} = cg\_model\_write(\textcolor{input}{char *ModelLabel}, \textcolor{input}{ModelType\_t ModelType}); & - w m \\
\textcolor{output}{\textit{ier}} = cg\_model\_read(\textcolor{input}{char *ModelLabel}, \textcolor{output}{\textit{ModelType\_t *ModelType}}); & r - m \\
\hline
call cg\_model\_write\_f(\textcolor{input}{ModelLabel}, \textcolor{input}{ModelType}, \textcolor{output}{\textit{ier}}) & - w m \\
call cg\_model\_read\_f(\textcolor{input}{ModelLabel}, \textcolor{output}{\textit{ModelType}}, \textcolor{output}{\textit{ier}}) & r - m \\
\end{fctbox}

\noindent
\textbf{\textcolor{input}{Input}/\textcolor{output}{\textit{Output}}}

\noindent (Note that for Fortran calls, all integer arguments are integer*4 in 32-bit mode and integer*8 in 64-bit mode.
See ``64-bit Fortran Portability and Issues" section.)

\begin{Ventryi}{\texttt{ModelLabel}}\raggedright
\item [\texttt{ModelLabel}]
      The CGNS label for the model being defined.
      The models supported by CGNS are:
      \begin{itemize*}
      \item \texttt{GasModel\_t}
      \item \texttt{ViscosityModel\_t}
      \item \texttt{ThermalConductivityModel\_t}
      \item \texttt{TurbulenceClosure\_t}
      \item \texttt{TurbulenceModel\_t}
      \item \texttt{ThermalRelaxationModel\_t}
      \item \texttt{ChemicalKineticsModel\_t}
      \item \texttt{EMElectricFieldModel\_t}
      \item \texttt{EMMagneticFieldModel\_t}
      \item \texttt{EMConductivityModel\_t}
      \end{itemize*}
      (\textcolor{input}{Input})
\item [\texttt{ModelType}]
      One of the model types (listed below) allowed for the
      \texttt{ModelLabel} selected.
      (\textcolor{input}{Input} for \texttt{cg\_model\_write};
      \textcolor{output}{\textit{output}} for \texttt{cg\_model\_read})
\item [\texttt{ier}]
      Error status.
      (\textcolor{output}{\textit{Output}})
\end{Ventryi}

\noindent
The types allowed for the various models are:

\begin{Ventryi}{\texttt{ThermalConductivityModel\_t}}\raggedright
\item [\texttt{GasModel\_t}]
      \texttt{CG\_Null, CG\_UserDefined, Ideal, VanderWaals, CaloricallyPerfect,
      ThermallyPerfect, ConstantDensity, RedlichKwong}
\item [\texttt{ViscosityModel\_t}]
      \texttt{CG\_Null, CG\_UserDefined, Constant, PowerLaw, SutherlandLaw}
\item [\texttt{ThermalConductivityModel\_t}]
      \texttt{CG\_Null, CG\_UserDefined, PowerLaw, SutherlandLaw, ConstantPrandtl}
\item [\texttt{TurbulenceModel\_t}]
      \texttt{CG\_Null, CG\_UserDefined, Algebraic\_BaldwinLomax,
      Algebraic\_CebeciSmith, HalfEquation\_JohnsonKing,
      OneEquation\_BaldwinBarth, OneEquation\_SpalartAllmaras,
      TwoEquation\_JonesLaunder, TwoEquation\_MenterSST,
      TwoEquation\_Wilcox}
\item [\texttt{TurbulenceClosure\_t}]
      \texttt{CG\_Null, CG\_UserDefined, EddyViscosity, ReynoldsStress,
      ReynoldsStressAlgebraic}
\item [\texttt{ThermalRelaxationModel\_t}]
      \texttt{CG\_Null, CG\_UserDefined, Frozen, ThermalEquilib,
      ThermalNonequilib}
\item [\texttt{ChemicalKineticsModel\_t}]
      \texttt{CG\_Null, CG\_UserDefined, Frozen,
      ChemicalEquilibCurveFit, ChemicalEquilibMinimization,
      ChemicalNonequilib}
\item [\texttt{EMElectricFieldModel\_t}]
      \texttt{CG\_Null, CG\_UserDefined, Constant, Frozen, Interpolated,
      Voltage}
\item [\texttt{EMMagneticFieldModel\_t}]
      \texttt{CG\_Null, CG\_UserDefined, Constant, Frozen, Interpolated}
\item [\texttt{EMConductivityModel\_t}]
      \texttt{CG\_Null, CG\_UserDefined, Constant, Frozen, Equilibrium\_LinRessler,
      Chemistry\_LinRessler}
\end{Ventryi}



\clearemptydoublepage
\section{Families}
\label{s:families}
\thispagestyle{plain}

\subsection{Family Definition}
\label{s:family}

\noindent
\textit{Node}: \fort{Family\_t}

\begin{fctbox}
\textcolor{output}{\textit{ier}} = cg\_family\_write(\textcolor{input}{int fn}, \textcolor{input}{int B}, \textcolor{input}{char *FamilyName}, \textcolor{output}{\textit{int *Fam}}); & - w m \\
\textcolor{output}{\textit{ier}} = cg\_nfamilies(\textcolor{input}{int fn}, \textcolor{input}{int B}, \textcolor{output}{\textit{int *nfamilies}}); & r - m \\
\textcolor{output}{\textit{ier}} = cg\_family\_read(\textcolor{input}{int fn}, \textcolor{input}{int B}, \textcolor{input}{int Fam}, \textcolor{output}{\textit{char *FamilyName}}, & r - m \\
~~~~~~\textcolor{output}{\textit{int *nFamBC}}, \textcolor{output}{\textit{int *nGeo}}); & \\
\hline
call cg\_family\_write\_f(\textcolor{input}{fn}, \textcolor{input}{B}, \textcolor{input}{FamilyName}, \textcolor{output}{\textit{Fam}}, \textcolor{output}{\textit{ier}}) & - w m \\
call cg\_nfamilies\_f(\textcolor{input}{fn}, \textcolor{input}{B}, \textcolor{output}{\textit{nfamilies}}, \textcolor{output}{\textit{ier}}) & r - m \\
call cg\_family\_read\_f(\textcolor{input}{fn}, \textcolor{input}{B}, \textcolor{input}{Fam}, \textcolor{output}{\textit{FamilyName}}, \textcolor{output}{\textit{nFamBC}}, \textcolor{output}{\textit{nGeo}}, \textcolor{output}{\textit{ier}}) & r - m \\
\end{fctbox}

\noindent
\textbf{\textcolor{input}{Input}/\textcolor{output}{\textit{Output}}}

\begin{Ventryi}{\fort{FamilyName}}\raggedright
\item [\fort{fn}]
      CGNS file index number.
      (\textcolor{input}{Input})
\item [\fort{B}]
      Base index number, where $1 \leq \text{\fort{B}} \leq \text{\fort{nbases}}$.
      (\textcolor{input}{Input})
\item [\fort{nfamilies}]
      Number of families in base \fort{B}.
      (\textcolor{output}{\textit{Output}})
\item [\fort{Fam}]
      Family index number, where $1 \leq \text{\fort{Fam}} \leq \text{\fort{nfamilies}}$.
      (\textcolor{input}{Input} for \fort{cg\_family\_read};
      \textcolor{output}{\textit{output}} for \fort{cg\_family\_write})
\item [\fort{FamilyName}]
      Name of the family.
      (\textcolor{input}{Input} for \fort{cg\_family\_write};
      \textcolor{output}{\textit{output}} for \fort{cg\_family\_read})
\item [\fort{nFamBC}]
      Number of boundary conditions for this family.
      This should be either 0 or 1.
      (\textcolor{output}{\textit{Output}})
\item [\fort{nGeo}]
      Number of geometry references for this family.
      (\textcolor{output}{\textit{Output}})
\item [\fort{ier}]
      Error status.
      (\textcolor{output}{\textit{Output}})
\end{Ventryi}

\newpage
\subsection{Geometry Reference}
\label{s:geometry}

\noindent
\textit{Node}: \fort{GeometryReference\_t}

\begin{fctbox}
\textcolor{output}{\textit{ier}} = cg\_geo\_write(\textcolor{input}{int fn}, \textcolor{input}{int B}, \textcolor{input}{int Fam}, \textcolor{input}{char *GeoName}, & - w m \\
~~~~~~\textcolor{input}{char *FileName}, \textcolor{input}{char *CADSystem}, \textcolor{output}{\textit{int *G}}); & \\
\textcolor{output}{\textit{ier}} = cg\_geo\_read(\textcolor{input}{int fn}, \textcolor{input}{int B}, \textcolor{input}{int Fam}, \textcolor{input}{int G}, \textcolor{output}{\textit{char *GeoName}}, & r - m \\
~~~~~~\textcolor{output}{\textit{char **FileName}}, \textcolor{output}{\textit{char *CADSystem}}, \textcolor{output}{\textit{int *nparts}}); & \\
\textcolor{output}{\textit{ier}} = cg\_part\_write(\textcolor{input}{int fn}, \textcolor{input}{int B}, \textcolor{input}{int Fam}, \textcolor{input}{int G}, \textcolor{input}{char *PartName}, & - w m \\
~~~~~~\textcolor{output}{\textit{int *P}}); & \\
\textcolor{output}{\textit{ier}} = cg\_part\_read(\textcolor{input}{int fn}, \textcolor{input}{int B}, \textcolor{input}{int Fam}, \textcolor{input}{int G}, \textcolor{input}{int P}, & r - m \\
~~~~~~\textcolor{output}{\textit{char *PartName}}); & \\
\hline
call cg\_geo\_write\_f(\textcolor{input}{fn}, \textcolor{input}{B}, \textcolor{input}{Fam}, \textcolor{input}{GeoName}, \textcolor{input}{FileName}, \textcolor{input}{CADSystem}, \textcolor{output}{\textit{G}}, \textcolor{output}{\textit{ier}}) & - w m \\
call cg\_geo\_read\_f(\textcolor{input}{fn}, \textcolor{input}{B}, \textcolor{input}{Fam}, \textcolor{input}{G}, \textcolor{output}{\textit{GeoName}}, \textcolor{output}{\textit{FileName}}, \textcolor{output}{\textit{CADSystem}}, & r - m \\
~~~~~\textcolor{output}{\textit{nparts}}, \textcolor{output}{\textit{ier}}) & \\
call cg\_part\_write\_f(\textcolor{input}{fn}, \textcolor{input}{B}, \textcolor{input}{Fam}, \textcolor{input}{G}, \textcolor{input}{PartName}, \textcolor{output}{\textit{P}}, \textcolor{output}{\textit{ier}}) & - w m \\
call cg\_part\_read\_f(\textcolor{input}{fn}, \textcolor{input}{B}, \textcolor{input}{Fam}, \textcolor{input}{G}, \textcolor{input}{P}, \textcolor{output}{\textit{PartName}}, \textcolor{output}{\textit{ier}}) & r - m \\
\end{fctbox}

\noindent
\textbf{\textcolor{input}{Input}/\textcolor{output}{\textit{Output}}}

\begin{Ventryi}{\fort{CADSystem}}\raggedright
\item [\fort{fn}]
      CGNS file index number.
      (\textcolor{input}{Input})
\item [\fort{B}]
      Base index number, where $1 \leq \text{\fort{B}} \leq \text{\fort{nbases}}$.
      (\textcolor{input}{Input})
\item [\fort{Fam}]
      Family index number, where $1 \leq \text{\fort{Fam}} \leq \text{\fort{nfamilies}}$.
      (\textcolor{input}{Input})
\item [\fort{G}]
      Geometry reference index number, where $1 \leq \text{\fort{G}} \leq \text{\fort{nGeo}}$.
      (\textcolor{input}{Input} for \fort{cg\_geo\_read},
      \fort{cg\_part\_write}, \fort{cg\_part\_read};
      \textcolor{output}{\textit{output}} for \fort{cg\_geo\_write})
\item [\fort{P}]
      Geometry entity index number, where $1 \leq \text{\fort{P}} \leq \text{\fort{nparts}}$.
      (\textcolor{input}{Input} for \fort{cg\_part\_read};
      \textcolor{output}{\textit{output}} for \fort{cg\_part\_write})
\item [\fort{GeoName}]
      Name of \fort{GeometryReference\_t} node.
      (\textcolor{input}{Input} for \fort{cg\_geo\_write};
      \textcolor{output}{\textit{output}} for \fort{cg\_geo\_read})
\item [\fort{FileName}]
      Name of geometry file.
      (\textcolor{input}{Input} for \fort{cg\_geo\_write};
      \textcolor{output}{\textit{output}} for \fort{cg\_geo\_read})
\item [\fort{CADSystem}]
      Geometry format.
      (\textcolor{input}{Input} for \fort{cg\_geo\_write};
      \textcolor{output}{\textit{output}} for \fort{cg\_geo\_read})
\item [\fort{nparts}]
      Number of geometry entities.
      (\textcolor{output}{\textit{Output}})
\item [\fort{PartName}]
      Name of a geometry entity in the file \fort{FileName}.
      (\textcolor{input}{Input} for \fort{cg\_part\_write};
      \textcolor{output}{\textit{output}} for \fort{cg\_part\_read})
\item [\fort{ier}]
      Error status.
      (\textcolor{output}{\textit{Output}})
\end{Ventryi}

Note that with \texttt{cg\_geo\_read} the memory for the filename
character string, \texttt{FileName}, will be allocated by the Mid-Level
Library.
The application code is responsible for releasing this memory when it is
no longer needed by calling \texttt{cg\_free(FileName)}, described in
\autoref{s:free}.

\newpage
\subsection{Family Boundary Condition}
\label{s:familybc}

\noindent
\textit{Node}: \fort{FamilyBC\_t}

\begin{fctbox}
\textcolor{output}{\textit{ier}} = cg\_fambc\_write(\textcolor{input}{int fn}, \textcolor{input}{int B}, \textcolor{input}{int Fam}, \textcolor{input}{char *FamBCName}, & - w m \\
~~~~~~\textcolor{input}{BCType\_t BCType}, \textcolor{output}{\textit{int *BC}}); & \\
\textcolor{output}{\textit{ier}} = cg\_fambc\_read(\textcolor{input}{int fn}, \textcolor{input}{int B}, \textcolor{input}{int Fam}, \textcolor{input}{int BC}, \textcolor{output}{\textit{char *FamBCName}}, & r - m \\
~~~~~~\textcolor{output}{\textit{BCType\_t *BCType}}); & \\
\hline
call cg\_fambc\_write\_f(\textcolor{input}{fn}, \textcolor{input}{B}, \textcolor{input}{Fam}, \textcolor{input}{FamBCName}, \textcolor{input}{BCType}, \textcolor{output}{\textit{BC}}, \textcolor{output}{\textit{ier}}) & - w m \\
call cg\_fambc\_read\_f(\textcolor{input}{fn}, \textcolor{input}{B}, \textcolor{input}{Fam}, \textcolor{input}{BC}, \textcolor{output}{\textit{FamBCName}}, \textcolor{output}{\textit{BCType}}, \textcolor{output}{\textit{ier}}) & r - m \\
\end{fctbox}

\noindent
\textbf{\textcolor{input}{Input}/\textcolor{output}{\textit{Output}}}

\begin{Ventryi}{\fort{FamBCName}}\raggedright
\item [\fort{fn}]
      CGNS file index number.
      (\textcolor{input}{Input})
\item [\fort{B}]
      Base index number, where $1 \leq \text{\fort{B}} \leq \text{\fort{nbases}}$.
      (\textcolor{input}{Input})
\item [\fort{Fam}]
      Family index number, where $1 \leq \text{\fort{Fam}} \leq \text{\fort{nfamilies}}$.
      (\textcolor{input}{Input})
\item [\fort{BC}]
      Family boundary condition index number.
      This must be equal to 1.
      (\textcolor{input}{Input} for \fort{cg\_fambc\_read};
      \textcolor{output}{\textit{output}} for \fort{cg\_fambc\_write})
\item [\fort{FamBCName}]
      Name of the \fort{FamilyBC\_t} node.
      (\textcolor{input}{Input} for \fort{cg\_fambc\_write};
      \textcolor{output}{\textit{output}} for \fort{cg\_fambc\_read})
\item [\fort{BCType}]
      Boundary condition type for the family.
      See the eligible types for \fort{BCType\_t} in \autoref{s:typedefs}.
      (\textcolor{input}{Input} for \fort{cg\_fambc\_write};
      \textcolor{output}{\textit{output}} for \fort{cg\_fambc\_read})
\item [\fort{ier}]
      Error status.
      (\textcolor{output}{\textit{Output}})
\end{Ventryi}

\subsection{Family Name}
\label{s:familyname}

\noindent
\textit{Node}: \fort{FamilyName\_t}

\begin{fctbox}
\textcolor{output}{\textit{ier}} = cg\_famname\_write(\textcolor{input}{char *FamilyName}); & - w m \\
\textcolor{output}{\textit{ier}} = cg\_famname\_read(\textcolor{output}{\textit{char *FamilyName}}); & r - m \\
\hline
call cg\_famname\_write\_f(\textcolor{input}{FamilyName}, \textcolor{output}{\textit{ier}}) & - w m \\
call cg\_famname\_read\_f(\textcolor{output}{\textit{FamilyName}}, \textcolor{output}{\textit{ier}}) & r - m \\
\end{fctbox}

\noindent
\textbf{\textcolor{input}{Input}/\textcolor{output}{\textit{Output}}}

\begin{Ventryi}{\fort{ptset\_type}}\raggedright
\item [\fort{FamilyName}]
      Family name.
      (\textcolor{input}{Input} for \fort{cg\_famname\_write};
      \textcolor{output}{\textit{output}} for \fort{cg\_famname\_read})
\item [\fort{ier}]
      Error status.
      (\textcolor{output}{\textit{Output}})
\end{Ventryi}


\clearemptydoublepage
\section{Time-Dependent Data}
\label{s:timedep}
\thispagestyle{plain}

\subsection{Base Iterative Data}
\label{s:biter}

\noindent
\textit{Node}: \texttt{BaseIterativeData\_t}

\begin{fctbox}
\textcolor{output}{\textit{ier}} = cg\_biter\_write(\textcolor{input}{int fn}, \textcolor{input}{int B}, \textcolor{input}{char *BaseIterName}, \textcolor{input}{int Nsteps}); & - w m \\
\textcolor{output}{\textit{ier}} = cg\_biter\_read(\textcolor{input}{int fn}, \textcolor{input}{int B}, \textcolor{output}{\textit{char *BaseIterName}}, \textcolor{output}{\textit{int *Nsteps}}); & r - m \\
\hline
call cg\_biter\_write\_f(\textcolor{input}{fn}, \textcolor{input}{B}, \textcolor{input}{BaseIterName}, \textcolor{input}{Nsteps}, \textcolor{output}{\textit{ier}}) & - w m \\
call cg\_biter\_read\_f(\textcolor{input}{fn}, \textcolor{input}{B}, \textcolor{output}{\textit{BaseIterName}}, \textcolor{output}{\textit{Nsteps}}, \textcolor{output}{\textit{ier}}) & r - m \\
\end{fctbox}

\noindent
\textbf{\textcolor{input}{Input}/\textcolor{output}{\textit{Output}}}

\begin{Ventryi}{\texttt{BaseIterName}}\raggedright
\item [\texttt{fn}]
      CGNS file index number.
      (\textcolor{input}{Input})
\item [\texttt{B}]
      Base index number, where $1 \leq \text{\texttt{B}} \leq \text{\texttt{nbases}}$.
      (\textcolor{input}{Input})
\item [\texttt{BaseIterName}]
      Name of the \texttt{BaseIterativeData\_t} node.
      (\textcolor{input}{Input} for \texttt{cg\_biter\_write};
      \textcolor{output}{\textit{output}} for \texttt{cg\_biter\_read})
\item [\texttt{Nsteps}]
      Number of time steps or iterations.
      (\textcolor{input}{Input} for \texttt{cg\_biter\_write};
      \textcolor{output}{\textit{output}} for \texttt{cg\_biter\_read})
\item [\texttt{ier}]
      Error status.
      (\textcolor{output}{\textit{Output}})
\end{Ventryi}

\subsection{Zone Iterative Data}
\label{s:ziter}

\noindent
\textit{Node}: \texttt{ZoneIterativeData\_t}

\begin{fctbox}
\textcolor{output}{\textit{ier}} = cg\_ziter\_write(\textcolor{input}{int fn}, \textcolor{input}{int B}, \textcolor{input}{int Z}, \textcolor{input}{char *ZoneIterName}); & - w m \\
\textcolor{output}{\textit{ier}} = cg\_ziter\_read(\textcolor{input}{int fn}, \textcolor{input}{int B}, \textcolor{input}{int Z}, \textcolor{output}{\textit{char *ZoneIterName}}); & r - m \\
\hline
call cg\_ziter\_write\_f(\textcolor{input}{fn}, \textcolor{input}{B}, \textcolor{input}{Z}, \textcolor{input}{ZoneIterName}, \textcolor{output}{\textit{ier}}) & - w m \\
call cg\_ziter\_read\_f(\textcolor{input}{fn}, \textcolor{input}{B}, \textcolor{input}{Z}, \textcolor{output}{\textit{ZoneIterName}}, \textcolor{output}{\textit{ier}}) & r - m \\
\end{fctbox}

\noindent
\textbf{\textcolor{input}{Input}/\textcolor{output}{\textit{Output}}}

\begin{Ventryi}{\texttt{ZoneIterName}}\raggedright
\item [\texttt{fn}]
      CGNS file index number.
      (\textcolor{input}{Input})
\item [\texttt{B}]
      Base index number, where $1 \leq \text{\texttt{B}} \leq \text{\texttt{nbases}}$.
      (\textcolor{input}{Input})
\item [\texttt{Z}]
      Family index number, where $1 \leq \text{\texttt{Z}} \leq \text{\texttt{nzones}}$.
      (\textcolor{input}{Input})
\item [\texttt{ZoneIterName}]
      Name of the \texttt{ZoneIterativeData\_t} node.
      (\textcolor{input}{Input} for \texttt{cg\_ziter\_write};
      \textcolor{output}{\textit{output}} for \texttt{cg\_ziter\_read})
\item [\texttt{ier}]
      Error status.
      (\textcolor{output}{\textit{Output}})
\end{Ventryi}

\subsection{Rigid Grid Motion}
\label{s:rigid}

\noindent
\textit{Node}: \texttt{RigidGridMotion\_t}

\begin{fctbox}
\textcolor{output}{\textit{ier}} = cg\_rigid\_motion\_write(\textcolor{input}{int fn}, \textcolor{input}{int B}, \textcolor{input}{int Z}, & - w m \\
~~~~~~\textcolor{input}{char *RigidGridMotionName}, & \\
~~~~~~\textcolor{input}{RigidGridMotionType\_t RigidGridMotionType}, \textcolor{output}{\textit{int *R}}); & \\
\textcolor{output}{\textit{ier}} = cg\_n\_rigid\_motions(\textcolor{input}{int fn}, \textcolor{input}{int B}, \textcolor{input}{int Z}, \textcolor{output}{\textit{int *n\_rigid\_motions}}); & r - m \\
\textcolor{output}{\textit{ier}} = cg\_rigid\_motion\_read(\textcolor{input}{int fn}, \textcolor{input}{int B}, \textcolor{input}{int Z}, \textcolor{input}{int R}, & r - m \\
~~~~~~\textcolor{output}{\textit{char *RigidGridMotionName}}, & \\
~~~~~~\textcolor{output}{\textit{RigidGridMotionType\_t RigidGridMotionType}}); & \\
\hline
call cg\_rigid\_motion\_write\_f(\textcolor{input}{fn}, \textcolor{input}{B}, \textcolor{input}{Z}, \textcolor{input}{RigidGridMotionName}, & - w m \\
~~~~~\textcolor{input}{RigidGridMotionType}, \textcolor{output}{\textit{R}}, \textcolor{output}{\textit{ier}}) & \\
call cg\_n\_rigid\_motions\_f(\textcolor{input}{fn}, \textcolor{input}{B}, \textcolor{input}{Z}, \textcolor{output}{\textit{n\_rigid\_motions}}, \textcolor{output}{\textit{ier}}) & r - m \\
call cg\_rigid\_motion\_read\_f(\textcolor{input}{fn}, \textcolor{input}{B}, \textcolor{input}{Z}, \textcolor{input}{R}, \textcolor{output}{\textit{RigidGridMotionName}}, & r - m \\
~~~~~\textcolor{output}{\textit{RigidGridMotionType}}, \textcolor{output}{\textit{ier}}) & \\
\end{fctbox}

\noindent
\textbf{\textcolor{input}{Input}/\textcolor{output}{\textit{Output}}}

\begin{Ventryi}{\texttt{RigidGridMotionName}}\raggedright
\item [\texttt{fn}]
      CGNS file index number.
      (\textcolor{input}{Input})
\item [\texttt{B}]
      Base index number, where $1 \leq \text{\texttt{B}} \leq \text{\texttt{nbases}}$.
      (\textcolor{input}{Input})
\item [\texttt{Z}]
      Family index number, where $1 \leq \text{\texttt{Z}} \leq \text{\texttt{nzones}}$.
      (\textcolor{input}{Input})
\item [\texttt{RigidGridMotionName}]
      Name of the \texttt{RigidGridMotion\_t} node.
      (\textcolor{input}{Input} for \texttt{cg\_rigid\_motion\_write};
      \textcolor{output}{\textit{output}} for \texttt{cg\_rigid\_motion\_read})
\item [\texttt{RigidGridMotionType}]
      Type of rigid grid motion.
      The admissible types are \texttt{CG\_Null}, \texttt{CG\_UserDefined},
      \texttt{ConstantRate}, and \texttt{VariableRate}.
      (\textcolor{input}{Input} for \texttt{cg\_rigid\_motion\_write};
      \textcolor{output}{\textit{output}} for \texttt{cg\_rigid\_motion\_read})
\item [\texttt{n\_rigid\_motions}]
      Number of \texttt{RigidGridMotion\_t} nodes under zone \texttt{Z}.
      (\textcolor{output}{\textit{Output}})
\item [\texttt{R}]
      Rigid rotation index number, where $1 \leq \text{\texttt{R}} \leq \text{\texttt{n\_rigid\_motions}}$.
      (\textcolor{input}{Input} for \texttt{cg\_rigid\_motion\_read};
      \textcolor{output}{\textit{output}} for \texttt{cg\_rigid\_motion\_write})
\item [\texttt{ier}]
      Error status.
      (\textcolor{output}{\textit{Output}})
\end{Ventryi}

\newpage
\subsection{Arbitrary Grid Motion}
\label{s:arbitrary}

\noindent
\textit{Node}: \texttt{ArbitraryGridMotion\_t}

\begin{fctbox}
\textcolor{output}{\textit{ier}} = cg\_arbitrary\_motion\_write(\textcolor{input}{int fn}, \textcolor{input}{int B}, \textcolor{input}{int Z}, & - w m \\
~~~~~~\textcolor{input}{char *ArbitraryGridMotionName}, & \\
~~~~~~\textcolor{input}{ArbitraryGridMotionType\_t ArbitraryGridMotionType}, \textcolor{output}{\textit{int *A}}); & \\
\textcolor{output}{\textit{ier}} = cg\_n\_arbitrary\_motions(\textcolor{input}{int fn}, \textcolor{input}{int B}, \textcolor{input}{int Z}, & r - m \\
~~~~~~\textcolor{output}{\textit{int *n\_arbitrary\_motions}}); & \\
\textcolor{output}{\textit{ier}} = cg\_arbitrary\_motion\_read(\textcolor{input}{int fn}, \textcolor{input}{int B}, \textcolor{input}{int Z}, \textcolor{input}{int A}, & r - m \\
~~~~~~\textcolor{output}{\textit{char *ArbitraryGridMotionName}}, & \\
~~~~~~\textcolor{output}{\textit{ArbitraryGridMotionType\_t ArbitraryGridMotionType}}); & \\
\hline
call cg\_arbitrary\_motion\_write\_f(\textcolor{input}{fn}, \textcolor{input}{B}, \textcolor{input}{Z}, \textcolor{input}{ArbitraryGridMotionName}, & - w m \\
~~~~~\textcolor{input}{ArbitraryGridMotionType}, \textcolor{output}{\textit{A}}, \textcolor{output}{\textit{ier}}) & \\
call cg\_n\_arbitrary\_motions\_f(\textcolor{input}{fn}, \textcolor{input}{B}, \textcolor{input}{Z}, \textcolor{output}{\textit{n\_arbitrary\_motions}}, \textcolor{output}{\textit{ier}}) & r - m \\
call cg\_arbitrary\_motion\_read\_f(\textcolor{input}{fn}, \textcolor{input}{B}, \textcolor{input}{Z}, \textcolor{input}{A}, \textcolor{output}{\textit{ArbitraryGridMotionName}}, & r - m \\
~~~~~\textcolor{output}{\textit{ArbitraryGridMotionType}}, \textcolor{output}{\textit{ier}}) & \\
\end{fctbox}

\noindent
\textbf{\textcolor{input}{Input}/\textcolor{output}{\textit{Output}}}

\begin{Ventryi}{\texttt{ArbitraryGridMotionName}}\raggedright
\item [\texttt{fn}]
      CGNS file index number.
      (\textcolor{input}{Input})
\item [\texttt{B}]
      Base index number, where $1 \leq \text{\texttt{B}} \leq \text{\texttt{nbases}}$.
      (\textcolor{input}{Input})
\item [\texttt{Z}]
      Family index number, where $1 \leq \text{\texttt{Z}} \leq \text{\texttt{nzones}}$.
      (\textcolor{input}{Input})
\item [\texttt{ArbitraryGridMotionName}]
      Name of the \texttt{ArbitraryGridMotion\_t} node.
      (\textcolor{input}{Input} for \texttt{cg\_arbitrary\_motion\_write};
      \textcolor{output}{\textit{output}} for \texttt{cg\_arbitrary\_motion\_read})
\item [\texttt{ArbitraryGridMotionType}]
      Type of arbitrary grid motion.
      The admissible types are \texttt{CG\_Null}, \texttt{CG\_UserDefined},
      \texttt{NonDeformingGrid}, and \texttt{DeformingGrid}.
      (\textcolor{input}{Input} for \texttt{cg\_arbitrary\_motion\_write};
      \textcolor{output}{\textit{output}} for \texttt{cg\_arbitrary\_motion\_read})
\item [\texttt{n\_arbitrary\_motions}]
      Number of \texttt{ArbitraryGridMotion\_t} nodes under zone \texttt{Z}.
      (\textcolor{output}{\textit{Output}})
\item [\texttt{A}]
      Arbitrary grid motion index number, where $1 \leq \text{\texttt{A}} \leq \text{\texttt{n\_arbitrary\_motions}}$.
      (\textcolor{input}{Input} for \texttt{cg\_arbitrary\_motion\_read};
      \textcolor{output}{\textit{output}} for \texttt{cg\_arbitrary\_motion\_write})
\item [\texttt{ier}]
      Error status.
      (\textcolor{output}{\textit{Output}})
\end{Ventryi}


\clearemptydoublepage
\section{Links}
\label{s:links}
\thispagestyle{plain}

\begin{fctbox}
\textcolor{output}{\textit{ier}} = cg\_link\_write(\textcolor{input}{char *nodename}, \textcolor{input}{char *filename}, & - w m \\
~~~~~~\textcolor{input}{char *name\_in\_file}); & \\
\textcolor{output}{\textit{ier}} = cg\_is\_link(\textcolor{output}{\textit{int *path\_length}}); & r - m \\
\textcolor{output}{\textit{ier}} = cg\_link\_read(\textcolor{output}{\textit{char **filename}}, \textcolor{output}{\textit{char **link\_path}}); & r - m \\
\hline
call cg\_link\_write\_f(\textcolor{input}{nodename}, \textcolor{input}{filename}, \textcolor{input}{name\_in\_file}, \textcolor{output}{\textit{ier}}) & - w m \\
call cg\_is\_link\_f(\textcolor{output}{\textit{path\_length}}, \textcolor{output}{\textit{ier}}) & r - m \\
call cg\_link\_read\_f(\textcolor{output}{\textit{filename}}, \textcolor{output}{\textit{link\_path}}, \textcolor{output}{\textit{ier}}) & r - m \\
\end{fctbox}

\noindent
\textbf{\textcolor{input}{Input}/\textcolor{output}{\textit{Output}}}

\begin{Ventryi}{\fort{name\_in\_file}}\raggedright
\item [\fort{nodename}]
      Name of the link node to create, e.g., \fort{GridCoordinates}.
      (\textcolor{input}{Input})
\item [\fort{filename}]
      Name of the linked file, or empty string if the link is within the
      same file.
      (\textcolor{input}{Input} for \fort{cg\_link\_write};
      \textcolor{output}{\textit{output}} for \fort{cg\_link\_read})
\item [\fort{name\_in\_file}]
      Path name of the node which the link points to.
      This can be a simple or a compound name, e.g.,
      \fort{Base/Zone 1/GridCoordinates}.
      (\textcolor{input}{Input})
\item [\fort{path\_length}]
      Length of the path name of the linked node.
      The value 0 is returned if the node is not a link.
      (\textcolor{output}{\textit{Output}})
\item [\fort{link\_path}]
      Path name of the node which the link points to.
      (\textcolor{output}{\textit{Output}})
\item [\fort{ier}]
      Error status.
      (\textcolor{output}{\textit{Output}})
\end{Ventryi}

Use \fort{cg\_goto(\_f)}, described in \autoref{s:navigating}, to
position to a location in the file prior to calling these routines.

When using \fort{cg\_link\_write}, the node being linked to does not have 
to exist when the link is created.
However, when the link is used, an error will occur if the linked-to
node does not exist.

Only nodes that support child nodes will support links.

It is assumed that the CGNS version for the file containing the link,
as determined by the \fort{CGNSLibraryVersion\_t} node, is also
applicable to \fort{filename}, the file containing the linked node.

Memory is allocated by the library for the return values of the C
function \fort{cg\_link\_read}.
This memory should be freed by the user when no longer needed by calling
\texttt{cg\_free(filename)} and \texttt{cg\_free(link\_path)}, described
in \autoref{s:free}.


\end{document}
